% ═══════════════════════════════════════════════════════════════
% Lista S01 — O que é Inteligência Artificial?
% Prof. Marcelo Keese Albertini · FACOM · UFU · GBC063
%
% Compilação:
%   Questões:  pdflatex -jobname=lista_s01_introducao_ia "\def\GABARITO{0}% ═══════════════════════════════════════════════════════════════
% Lista S01 — O que é Inteligência Artificial?
% Prof. Marcelo Keese Albertini · FACOM · UFU · GBC063
%
% Compilação:
%   Questões:  pdflatex -jobname=lista_s01_introducao_ia "\def\GABARITO{0}% ═══════════════════════════════════════════════════════════════
% Lista S01 — O que é Inteligência Artificial?
% Prof. Marcelo Keese Albertini · FACOM · UFU · GBC063
%
% Compilação:
%   Questões:  pdflatex -jobname=lista_s01_introducao_ia "\def\GABARITO{0}% ═══════════════════════════════════════════════════════════════
% Lista S01 — O que é Inteligência Artificial?
% Prof. Marcelo Keese Albertini · FACOM · UFU · GBC063
%
% Compilação:
%   Questões:  pdflatex -jobname=lista_s01_introducao_ia "\def\GABARITO{0}\input{lista_s01_introducao_ia.tex}"
%   Gabarito:  pdflatex -jobname=gabarito_lista_s01_introducao_ia "\def\GABARITO{1}\input{lista_s01_introducao_ia.tex}"
% ═══════════════════════════════════════════════════════════════

\documentclass[12pt,a4paper]{article}

\usepackage[utf8]{inputenc}
\usepackage[T1]{fontenc}
\usepackage[brazil]{babel}
\usepackage{geometry}
\geometry{margin=2.2cm}
\usepackage{enumitem}
\usepackage{amsmath,amssymb}
\usepackage{xcolor}
\usepackage{hyperref}
% Using standard fonts available on any TeX distribution
\usepackage{mathpazo}
\usepackage[scaled=0.9]{helvet}
\usepackage{courier}
\renewcommand{\familydefault}{\rmdefault}

% ── Conditional: gabarito on/off ──
\ifdefined\GABARITO
  \ifnum\GABARITO=1
    \newcommand{\sol}[1]{\textcolor{blue}{#1}}
    \newcommand{\errocomum}[1]{\textcolor{orange}{⚠️ Erro comum: #1}}
    \newcommand{\competencia}[1]{\textcolor{gray}{Competência: #1}}
  \else
    \newcommand{\sol}[1]{}
    \newcommand{\errocomum}[1]{}
    \newcommand{\competencia}[1]{}
  \fi
\else
  \newcommand{\sol}[1]{}
  \newcommand{\errocomum}[1]{}
  \newcommand{\competencia}[1]{}
\fi

% ── Question formatting ──
\newcommand{\nivel}[1]{\textcolor{gray}{[#1]}}
\newenvironment{questao}[2]{
  \vspace{0.5cm}
  \noindent\textbf{Q#1} \nivel{#2} \quad
}{
  \vspace{0.2cm}
}

\begin{document}

% ════════════════ HEADER ════════════════
\begin{center}
  {\sffamily\Large GBC063 · Semana 1 · Inteligência Artificial}

  \vspace{0.2cm}
  {\small Prof. Marcelo Keese Albertini · FACOM · UFU}

  \vspace{0.2cm}
  {\footnotesize Cobre as competências do resumo S01 ·
  \href{../semana01/material_introducao_ia.html}{material de estudo} ·
  \href{../semana01/resumo_s01_introducao_ia.md}{resumo}}

  \ifdefined\GABARITO
    \ifnum\GABARITO=1
      \vspace{0.3cm}
      \fbox{\textcolor{red}{\textbf{GABARITO — Documento Exclusivo do Professor}}}
    \fi
  \fi
\end{center}

\vspace{0.3cm}
\hrule
\vspace{0.5cm}

% ════════════════ NÍVEL ● (REATIVAÇÃO) ════════════════
\section*{Nível ● \quad Reativação}

\begin{questao}{1}{●}
  Defina \textbf{agente racional} segundo Russell \& Norvig e explique,
  em no máximo 3 linhas, por que esta definição — e não o Teste de Turing —
  é adotada como definição operacional do curso.

  \sol{
    \textbf{Resposta:} Um agente racional é aquele que, para cada percepção possível,
    seleciona a ação que maximiza a métrica de desempenho esperada. O curso adota esta
    definição porque ela (1) funciona com qualquer métrica de sucesso, (2) não depende
    de replicar comportamento humano, e (3) unifica desde um termostato até o AlphaGo.

    \textbf{Resolução:} Russell \& Norvig (2013), Cap. 2. A definição de agente racional
    é independente de humanidade — um agente pode ser racional e não-humano (ex.: AlphaGo
    joga Go de forma sobre-humana).

    \errocomum{Definir agente racional como "agente que age como um humano". Isso é a visão
    "agir humanamente" (Teste de Turing), não "agir racionalmente".}

    \competencia{Explicar a diferença entre IA simbólica e aprendizado de máquina —
    e definir agente racional.}
  }
\end{questao}

\begin{questao}{2}{●}
  As quatro visões clássicas da IA organizam-se em uma matriz 2×2:
  pensamento \textit{vs.} ação × desempenho humano \textit{vs.} racionalidade.
  Classifique cada um dos sistemas abaixo em \textbf{um} dos quatro quadrantes
  e justifique em 1 linha:

  \begin{enumerate}[label=(\alph*),nosep]
    \item ELIZA (Weizenbaum, 1966)
    \item Deep Blue (IBM, 1997)
    \item O agente aspirador de pó da Aula B
    \item Um modelo de linguagem treinado para prever o próximo token
  \end{enumerate}

  \sol{
    \textbf{Respostas:}
    (a) Agir humanamente — ELIZA simulava um terapeuta humano sem raciocinar.
    (b) Agir racionalmente — Deep Blue selecionava lances que maximizavam chance de vitória,
    sem imitar o pensamento de um grande mestre humano.
    (c) Agir racionalmente — o aspirador maximiza a limpeza dada sua percepção.
    (d) Pensar racionalmente (durante o treino: otimiza uma função de perda) e
    agir humanamente (na inferência: gera texto indistinguível do humano).

    \errocomum{Classificar ELIZA como "pensar humanamente". ELIZA não modelava pensamento —
    apenas casava padrões de string.}

    \competencia{Classificar qualquer sistema de IA nas quatro visões e justificar.}
  }
\end{questao}

\begin{questao}{3}{●}
  Verdadeiro ou falso. \textbf{Justifique cada resposta em até 2 linhas.}

  \begin{enumerate}[label=(\alph*),nosep]
    \item "O primeiro inverno da IA (1969–1980) foi causado exclusivamente por falta
    de poder computacional."
    \item "Redes neurais artificiais são uma ideia dos anos 2010."
  \end{enumerate}

  \sol{
    \textbf{Respostas:}
    (a) \textbf{Falso.} O inverno foi causado principalmente pelo livro \textit{Perceptrons}
    (Minsky \& Papert, 1969), que provou limitações do perceptron simples e levou o campo
    a concluir — erroneamente — que redes neurais eram um beco sem saída. A falta de
    computação contribuiu, mas o fator decisivo foi a perda de credibilidade e financiamento.

    (b) \textbf{Falso.} O primeiro modelo de neurônio artificial é de McCulloch \& Pitts (1943).
    O perceptron é de Rosenblatt (1959). O backpropagation que as treina é de 1986
    (Rumelhart, Hinton \& Williams). As ideias fundamentais têm 40–80 anos; o que mudou
    em 2012 foi a escala (dados + GPUs).

    \errocomum{Para (a): atribuir o inverno só à falta de hardware. O fator principal foi
    sociológico — a comunidade perdeu confiança nas redes neurais.}
    \errocomum{Para (b): achar que deep learning é uma invenção recente. As ideias são antigas;
    o que é novo é a escala.}

    \competencia{Enunciar a Bitter Lesson e dar exemplos históricos.}
  }
\end{questao}

% ════════════════ NÍVEL ●● (EXECUÇÃO) ════════════════
\section*{Nível ●● \quad Execução}

\begin{questao}{4}{●●}
  Considere o seguinte ambiente de aspirador de pó com \textbf{três} cômodos
  em linha (A–B–C). O agente começa na posição B. O cômodo A está sujo; B e C
  estão limpos. O agente segue exatamente a lógica abaixo:

  \begin{verbatim}
  def agente_aspirador(percepcao):
      posicao, sujo = percepcao
      if sujo:        return "ASPIRAR"
      if posicao == "A": return "DIREITA"
      if posicao == "B": return "DIREITA"
      return "ESQUERDA"
  \end{verbatim}

  Preencha a tabela de execução para os \textbf{6 primeiros passos}, indicando:
  posição atual, sensor (sujo/limpo), ação tomada, e estado dos cômodos após a ação.
  O agente limpa todos os cômodos? Em quantos passos?

  \sol{
    \textbf{Resposta:}

    \begin{tabular}{c|c|c|c|c}
      Passo & Posição & Sensor & Ação & Cômodos [A,B,C] \\
      \hline
      0 & B & limpo & DIREITA & [sujo, limpo, limpo] \\
      1 & C & limpo & ESQUERDA & [sujo, limpo, limpo] \\
      2 & B & limpo & DIREITA & [sujo, limpo, limpo] \\
      3 & C & limpo & ESQUERDA & [sujo, limpo, limpo] \\
      \dots & \dots & \dots & \dots & \dots \\
    \end{tabular}

    O agente \textbf{nunca alcança o cômodo A} porque sua lógica só vai de A→B→C→B→C→B→C\ldots
    Ele fica preso no ciclo B↔C. Para alcançar A, seria preciso uma regra adicional
    (ex.: "se estou em C e ambos A e B estão sujos, vá para B").

    \textbf{Resolução:} O bug está na lógica de movimento: quando a posição é B ou C,
    o agente nunca se move para A. A condição \texttt{posicao == "A"} nunca é verdadeira
    se o agente não consegue chegar a A.

    \errocomum{Assumir que o agente "dá a volta". O ambiente é uma linha, não um ciclo —
    de C só se vai para B.}

    \competencia{Implementar o agente aspirador e prever seu comportamento
    com 3 cômodos.}
  }
\end{questao}

\begin{questao}{5}{●●}
  O motor de encadeamento progressivo abaixo recebe uma base de fatos inicial
  e um conjunto de regras. Execute o algoritmo \textbf{passo a passo} com os
  dados fornecidos e indique o conjunto final de fatos.

  \begin{verbatim}
  fatos = {"chovendo", "temperatura < 15"}
  regras = [
    ({"chovendo", "temperatura < 15"}, "risco de gelo"),
    ({"risco de gelo"},              "fechar pista"),
    ({"chovendo"},                   "pista molhada"),
  ]
  \end{verbatim}

  \sol{
    \textbf{Resposta:} fatos finais = \{\texttt{"chovendo"}, \texttt{"temperatura < 15"},
    \texttt{"risco de gelo"}, \texttt{"fechar pista"}, \texttt{"pista molhada"}\}

    \textbf{Resolução:}
    \begin{enumerate}[nosep]
      \item Iteração 1: chovendo + temperatura < 15 → adiciona \texttt{"risco de gelo"}.
      \item Iteração 2: risco de gelo → adiciona \texttt{"fechar pista"}.
      \item Iteração 3: chovendo → adiciona \texttt{"pista molhada"}.
      \item Iteração 4: nenhuma regra nova dispara → fim.
    \end{enumerate}

    \errocomum{Esquecer que as regras disparam em qualquer ordem — o algoritmo testa TODAS
    as regras a cada iteração, não para na primeira que dispara.}

    \competencia{Explicar o funcionamento de um motor de inferência e por que regras
    não escalam.}
  }
\end{questao}

\begin{questao}{6}{●●}
  Em 2023, um LLM foi capaz de passar no exame da ordem dos advogados (Bar Exam)
  dos EUA sem ter recebido uma única regra explícita sobre direito. Usando os
  conceitos da Semana 1, explique \textbf{por que isso é um fato tão significativo}
  na história da IA. Se não obteve uma resposta completa, use o roteiro:
  (1) como um sistema especialista resolveria o Bar Exam? (2) por que essa abordagem
  falharia? (3) o que o LLM fez de diferente?

  \sol{
    \textbf{Resposta:} Um sistema especialista exigiria que engenheiros de conhecimento
    codificassem milhares de regras jurídicas à mão — um gargalo intransponível. O LLM
    não recebeu regra nenhuma: aprendeu padrões estatísticos de centenas de gigabytes
    de texto jurídico e geral. Isso é a Bitter Lesson em ação: conhecimento aprendido
    de dados venceu conhecimento artesanal — em um domínio (direito) que se supunha
    exigir raciocínio simbólico explícito.

    \errocomum{Dizer que o LLM "decorou" as respostas. LLMs não têm memória explícita do
    corpus de treino — eles aprendem padrões estatísticos, não fazem lookup de frases.}

    \competencia{Enunciar a Bitter Lesson e dar exemplos que a corroboram.}
  }
\end{questao}

% ════════════════ NÍVEL ●●● (ANÁLISE) ════════════════
\section*{Nível ●●● \quad Análise}

\begin{questao}{7}{●●●}
  Considere a afirmação de Sutton (2019):

  \begin{quote}
  \textit{``Métodos gerais que alavancam computação são, em última instância,
  os mais efetivos — e por uma grande margem. A maioria do esforço em IA é
  gasto tentando construir conhecimento em agentes; isso ajuda no curto prazo,
  mas no longo prazo atinge um platô.''}
  \end{quote}

  \begin{enumerate}[label=(\alph*),nosep]
    \item Por que conhecimento artesanal \textit{necessariamente} atinge um platô,
    enquanto métodos que aprendem de dados não? (Dica: pense no que acontece quando
    você dobra a quantidade de dados ou de computação em cada caso.)
    \item Dê um exemplo \textbf{concreto} de um sistema que atingiu esse platô e
    explique \textit{o que especificamente} o limitou.
  \end{enumerate}

  \sol{
    \textbf{Respostas:}
    (a) Conhecimento artesanal atinge um platô porque sua melhoria depende de \textbf{humanos
    adicionando mais regras} — um processo que é linear, caro e sujeito a contradições.
    Métodos que aprendem melhoram quando você adiciona mais dados ou computação, e esse
    processo é (em princípio) ilimitado: mais dados → melhores gradientes → melhor modelo.
    Não há um humano no loop ditando o teto.

    (b) Sistemas especialistas médicos (ex.: MYCIN, anos 80). Eram melhores que médicos
    juniores em diagnósticos específicos, mas cada nova doença exigia novas regras escritas
    por especialistas. Quando o número de regras passou de centenas, as interações entre
    elas tornaram-se imprevisíveis. O sistema não melhorava com mais dados de pacientes —
    só com mais regras escritas à mão. Enquanto isso, métodos de ML atuais melhoram
    automaticamente com mais exames, prontuários e artigos.

    \errocomum{Em (a), dizer que conhecimento artesanal atinge platô "porque humanos
    são limitados". A resposta precisa identificar o mecanismo: a taxa de melhoria é
    ditada pela velocidade com que humanos escrevem regras, não pela velocidade com
    que computação cresce.}

    \competencia{Enunciar a Bitter Lesson e explicar o mecanismo por trás dela.}
  }
\end{questao}

\begin{questao}{8}{●●●}
  O efeito ELIZA — atribuir inteligência onde há apenas casamento de padrões —
  foi documentado em 1966. Ele persiste em 2025 com LLMs?

  \begin{enumerate}[label=(\alph*),nosep]
    \item Aponte \textbf{uma semelhança} entre ELIZA e um LLM moderno que torna
    o efeito ELIZA relevante hoje.
    \item Aponte \textbf{uma diferença fundamental} que torna a comparação enganosa.
    \item Em no máximo 4 linhas: um LLM é ``só casamento de padrões''? Justifique.
  \end{enumerate}

  \sol{
    \textbf{Respostas:}
    (a) Ambos produzem texto fluente sem garantia de compreensão do conteúdo semântico.
    Um LLM pode gerar uma explicação convincente de um conceito que ele não ``entende''
    no sentido humano — assim como ELIZA simulava empatia sem sentir nada. O efeito
    ELIZA ocorre quando o usuário projeta compreensão onde há apenas correlação estatística.

    (b) ELIZA usava algumas dezenas de regras de reescrita escritas à mão. Um LLM aprende
    de centenas de bilhões de parâmetros ajustados sobre trilhões de tokens. A escala e
    a origem do conhecimento (artesanal vs. aprendido) são de natureza completamente diferente.

    (c) Não. Reduzir um LLM a ``só casamento de padrões'' é como reduzir um cérebro a
    ``só sinapses disparando'' — tecnicamente verdadeiro, mas enganoso. A questão relevante
    não é se o mecanismo básico é simples (é: multiplicação de matrizes + não-linearidades),
    mas se a escala produz comportamentos emergentes que merecem ser chamados de raciocínio.
    Esta é uma pergunta em aberto na pesquisa em 2025.

    \errocomum{Responder (c) com um simples ``sim'' ou ``não''. A pergunta exige distinguir
    mecanismo de comportamento — o que o sistema \textit{faz} vs. como ele \textit{faz}.}

    \competencia{Classificar sistemas de IA e analisar criticamente alegações de inteligência.}
  }
\end{questao}

% ════════════════ NÍVEL ★ (SÍNTESE) ════════════════
\section*{Nível ★ \quad Síntese}

\begin{questao}{9}{★}
  Projete um experimento mental: você herda um sistema especialista médico com
  10.000 regras, desenvolvido ao longo de 15 anos por uma equipe de 20 engenheiros
  de conhecimento. O sistema tem 92\% de acurácia em diagnóstico.

  \begin{enumerate}[label=(\alph*),nosep]
    \item Identifique \textbf{três} cenários em que este sistema falharia de forma
    sistemática que um modelo de ML treinado sobre os mesmos dados não falharia.
    \item Para cada cenário, explique \textbf{por que} a natureza artesanal das
    regras é a causa raiz da falha.
    \item Proponha uma arquitetura híbrida (simbólico + aprendizado) que combine
    o melhor dos dois mundos para este domínio. Descreva em no máximo 6 linhas.
  \end{enumerate}

  \sol{
    \textbf{Respostas:}
    (a) e (b):
    \begin{enumerate}[nosep]
      \item \textbf{Nova doença:} o SE falha porque nenhuma regra cobre o fenômeno novo.
      O ML pode detectar que os sintomas formam um cluster distinto nos dados e sinalizar
      uma anomalia — mesmo sem rótulos. Causa raiz: SE só reconhece o que foi explicitamente
      programado.
      \item \textbf{Interação entre condições:} um paciente com diabetes \textit{e} COVID
      pode ativar regras contraditórias (uma recomenda corticoides, outra contraindica).
      O SE não tem como ponderar evidências conflitantes. O ML aprende pesos relativos
      dos features. Causa raiz: regras são discretas e não-compostas; não modelam interações.
      \item \textbf{Deriva de distribuição:} se a população de pacientes mudar (ex.:
      envelhecimento), a acurácia do SE cai porque as regras foram escritas para a
      distribuição antiga. O ML pode ser retreinado com dados novos. Causa raiz: SE
      não se adapta sem reengenharia manual.
    \end{enumerate}

    (c) \textbf{Arquitetura híbrida:} usar um modelo de ML (ex.: gradient-boosted trees
    ou rede neural) para produzir o diagnóstico base e a estimativa de confiança. As
    10.000 regras viram um sistema de \textbf{verificação}: antes de emitir o diagnóstico,
    o sistema checa se ele viola alguma regra de segurança crítica (ex.: ``nunca prescrever
    droga X para grávidas''). Se violar, escala para um médico humano. O ML lida com a
    generalização; as regras lidam com as restrições rígidas de segurança.

    \errocomum{Propor um híbrido em que as regras \textit{guiam} o aprendizado. Isso
    reintroduz o gargalo: as regras limitam o que o modelo pode aprender. O melhor
    híbrido usa regras como \textit{restrições de segurança}, não como indutores de viés.}

    \competencia{Sintetizar os conceitos de IA simbólica e aprendizado em uma arquitetura
    fundamentada.}
  }
\end{questao}

\begin{questao}{10}{★ 🌉}
  O treinamento do GPT usa a mesma estrutura de loop que o agente aspirador da Aula B:
  \textbf{percepção → decisão → ação}. A diferença está na complexidade de \texttt{decidir()}.

  \begin{enumerate}[label=(\alph*),nosep]
    \item No aspirador, \texttt{decidir()} é uma função com 4 regras \texttt{if}.
    No GPT, \texttt{decidir()} é uma rede neural com centenas de bilhões de
    parâmetros. \textbf{Por que}, conceitualmente, a mesma interface de 4 linhas
    funciona para ambos? O que isso revela sobre o paradigma de agentes?
    \item A Bitter Lesson argumenta que conhecimento artesanal atinge um platô.
    O prompt engineering — escrever instruções cuidadosas para um LLM — é uma
    forma de conhecimento artesanal? Justifique em até 5 linhas.
  \end{enumerate}

  \sol{
    \textbf{Respostas:}
    (a) A interface \texttt{obs → decidir(obs) → ação} funciona para ambos porque ela
    abstrai \textbf{o que} o agente faz de \textbf{como} ele decide. O paradigma de
    agentes é uma \textit{casca}: define a assinatura da decisão, não sua implementação.
    O que muda entre o aspirador e o GPT é apenas a função \texttt{decidir()} — e essa
    função pode ser desde 4 ifs até uma rede de 175 bilhões de parâmetros. Isso revela
    que o paradigma de agentes é \textbf{universal}: ele unifica toda a IA sob uma
    única abstração.

    (b) Sim, prompt engineering é conhecimento artesanal — mas de um tipo novo. É
    artesanal porque um humano está manualmente ajustando a entrada para melhorar
    a saída, sem que o modelo aprenda com esse ajuste. Mas difere do artesanal
    clássico porque (1) o modelo subjacente foi aprendido de dados, não programado,
    e (2) o custo de ajustar um prompt é ordens de grandeza menor que o de escrever
    regras. Dito isso, a Bitter Lesson sugere que, no longo prazo, métodos que
    \textit{aprendem} a fazer prompting (auto-prompting, RLHF) substituirão o
    prompting manual — e já há evidências disso.

    \errocomum{Em (b), responder simplesmente ``não'' porque o modelo foi treinado com ML.
    O prompt em si é conhecimento artesanal injetado \textit{na inferência}, não no treino.}

    \competencia{Reconhecer o loop percepção→decisão→ação como a interface universal
    da IA e aplicar a Bitter Lesson a desenvolvimentos recentes.}
  }
\end{questao}

\vspace{1cm}
\hrule
\vspace{0.3cm}
\begin{center}
  {\footnotesize Lista S01 · GBC063 · Prof. Marcelo Keese Albertini · FACOM/UFU}
  \ifdefined\GABARITO
    \ifnum\GABARITO=1
      \\[0.2cm] {\footnotesize\color{red}\textbf{GABARITO — USO EXCLUSIVO DO PROFESSOR}}
    \fi
  \fi
\end{center}

\end{document}
"
%   Gabarito:  pdflatex -jobname=gabarito_lista_s01_introducao_ia "\def\GABARITO{1}% ═══════════════════════════════════════════════════════════════
% Lista S01 — O que é Inteligência Artificial?
% Prof. Marcelo Keese Albertini · FACOM · UFU · GBC063
%
% Compilação:
%   Questões:  pdflatex -jobname=lista_s01_introducao_ia "\def\GABARITO{0}\input{lista_s01_introducao_ia.tex}"
%   Gabarito:  pdflatex -jobname=gabarito_lista_s01_introducao_ia "\def\GABARITO{1}\input{lista_s01_introducao_ia.tex}"
% ═══════════════════════════════════════════════════════════════

\documentclass[12pt,a4paper]{article}

\usepackage[utf8]{inputenc}
\usepackage[T1]{fontenc}
\usepackage[brazil]{babel}
\usepackage{geometry}
\geometry{margin=2.2cm}
\usepackage{enumitem}
\usepackage{amsmath,amssymb}
\usepackage{xcolor}
\usepackage{hyperref}
% Using standard fonts available on any TeX distribution
\usepackage{mathpazo}
\usepackage[scaled=0.9]{helvet}
\usepackage{courier}
\renewcommand{\familydefault}{\rmdefault}

% ── Conditional: gabarito on/off ──
\ifdefined\GABARITO
  \ifnum\GABARITO=1
    \newcommand{\sol}[1]{\textcolor{blue}{#1}}
    \newcommand{\errocomum}[1]{\textcolor{orange}{⚠️ Erro comum: #1}}
    \newcommand{\competencia}[1]{\textcolor{gray}{Competência: #1}}
  \else
    \newcommand{\sol}[1]{}
    \newcommand{\errocomum}[1]{}
    \newcommand{\competencia}[1]{}
  \fi
\else
  \newcommand{\sol}[1]{}
  \newcommand{\errocomum}[1]{}
  \newcommand{\competencia}[1]{}
\fi

% ── Question formatting ──
\newcommand{\nivel}[1]{\textcolor{gray}{[#1]}}
\newenvironment{questao}[2]{
  \vspace{0.5cm}
  \noindent\textbf{Q#1} \nivel{#2} \quad
}{
  \vspace{0.2cm}
}

\begin{document}

% ════════════════ HEADER ════════════════
\begin{center}
  {\sffamily\Large GBC063 · Semana 1 · Inteligência Artificial}

  \vspace{0.2cm}
  {\small Prof. Marcelo Keese Albertini · FACOM · UFU}

  \vspace{0.2cm}
  {\footnotesize Cobre as competências do resumo S01 ·
  \href{../semana01/material_introducao_ia.html}{material de estudo} ·
  \href{../semana01/resumo_s01_introducao_ia.md}{resumo}}

  \ifdefined\GABARITO
    \ifnum\GABARITO=1
      \vspace{0.3cm}
      \fbox{\textcolor{red}{\textbf{GABARITO — Documento Exclusivo do Professor}}}
    \fi
  \fi
\end{center}

\vspace{0.3cm}
\hrule
\vspace{0.5cm}

% ════════════════ NÍVEL ● (REATIVAÇÃO) ════════════════
\section*{Nível ● \quad Reativação}

\begin{questao}{1}{●}
  Defina \textbf{agente racional} segundo Russell \& Norvig e explique,
  em no máximo 3 linhas, por que esta definição — e não o Teste de Turing —
  é adotada como definição operacional do curso.

  \sol{
    \textbf{Resposta:} Um agente racional é aquele que, para cada percepção possível,
    seleciona a ação que maximiza a métrica de desempenho esperada. O curso adota esta
    definição porque ela (1) funciona com qualquer métrica de sucesso, (2) não depende
    de replicar comportamento humano, e (3) unifica desde um termostato até o AlphaGo.

    \textbf{Resolução:} Russell \& Norvig (2013), Cap. 2. A definição de agente racional
    é independente de humanidade — um agente pode ser racional e não-humano (ex.: AlphaGo
    joga Go de forma sobre-humana).

    \errocomum{Definir agente racional como "agente que age como um humano". Isso é a visão
    "agir humanamente" (Teste de Turing), não "agir racionalmente".}

    \competencia{Explicar a diferença entre IA simbólica e aprendizado de máquina —
    e definir agente racional.}
  }
\end{questao}

\begin{questao}{2}{●}
  As quatro visões clássicas da IA organizam-se em uma matriz 2×2:
  pensamento \textit{vs.} ação × desempenho humano \textit{vs.} racionalidade.
  Classifique cada um dos sistemas abaixo em \textbf{um} dos quatro quadrantes
  e justifique em 1 linha:

  \begin{enumerate}[label=(\alph*),nosep]
    \item ELIZA (Weizenbaum, 1966)
    \item Deep Blue (IBM, 1997)
    \item O agente aspirador de pó da Aula B
    \item Um modelo de linguagem treinado para prever o próximo token
  \end{enumerate}

  \sol{
    \textbf{Respostas:}
    (a) Agir humanamente — ELIZA simulava um terapeuta humano sem raciocinar.
    (b) Agir racionalmente — Deep Blue selecionava lances que maximizavam chance de vitória,
    sem imitar o pensamento de um grande mestre humano.
    (c) Agir racionalmente — o aspirador maximiza a limpeza dada sua percepção.
    (d) Pensar racionalmente (durante o treino: otimiza uma função de perda) e
    agir humanamente (na inferência: gera texto indistinguível do humano).

    \errocomum{Classificar ELIZA como "pensar humanamente". ELIZA não modelava pensamento —
    apenas casava padrões de string.}

    \competencia{Classificar qualquer sistema de IA nas quatro visões e justificar.}
  }
\end{questao}

\begin{questao}{3}{●}
  Verdadeiro ou falso. \textbf{Justifique cada resposta em até 2 linhas.}

  \begin{enumerate}[label=(\alph*),nosep]
    \item "O primeiro inverno da IA (1969–1980) foi causado exclusivamente por falta
    de poder computacional."
    \item "Redes neurais artificiais são uma ideia dos anos 2010."
  \end{enumerate}

  \sol{
    \textbf{Respostas:}
    (a) \textbf{Falso.} O inverno foi causado principalmente pelo livro \textit{Perceptrons}
    (Minsky \& Papert, 1969), que provou limitações do perceptron simples e levou o campo
    a concluir — erroneamente — que redes neurais eram um beco sem saída. A falta de
    computação contribuiu, mas o fator decisivo foi a perda de credibilidade e financiamento.

    (b) \textbf{Falso.} O primeiro modelo de neurônio artificial é de McCulloch \& Pitts (1943).
    O perceptron é de Rosenblatt (1959). O backpropagation que as treina é de 1986
    (Rumelhart, Hinton \& Williams). As ideias fundamentais têm 40–80 anos; o que mudou
    em 2012 foi a escala (dados + GPUs).

    \errocomum{Para (a): atribuir o inverno só à falta de hardware. O fator principal foi
    sociológico — a comunidade perdeu confiança nas redes neurais.}
    \errocomum{Para (b): achar que deep learning é uma invenção recente. As ideias são antigas;
    o que é novo é a escala.}

    \competencia{Enunciar a Bitter Lesson e dar exemplos históricos.}
  }
\end{questao}

% ════════════════ NÍVEL ●● (EXECUÇÃO) ════════════════
\section*{Nível ●● \quad Execução}

\begin{questao}{4}{●●}
  Considere o seguinte ambiente de aspirador de pó com \textbf{três} cômodos
  em linha (A–B–C). O agente começa na posição B. O cômodo A está sujo; B e C
  estão limpos. O agente segue exatamente a lógica abaixo:

  \begin{verbatim}
  def agente_aspirador(percepcao):
      posicao, sujo = percepcao
      if sujo:        return "ASPIRAR"
      if posicao == "A": return "DIREITA"
      if posicao == "B": return "DIREITA"
      return "ESQUERDA"
  \end{verbatim}

  Preencha a tabela de execução para os \textbf{6 primeiros passos}, indicando:
  posição atual, sensor (sujo/limpo), ação tomada, e estado dos cômodos após a ação.
  O agente limpa todos os cômodos? Em quantos passos?

  \sol{
    \textbf{Resposta:}

    \begin{tabular}{c|c|c|c|c}
      Passo & Posição & Sensor & Ação & Cômodos [A,B,C] \\
      \hline
      0 & B & limpo & DIREITA & [sujo, limpo, limpo] \\
      1 & C & limpo & ESQUERDA & [sujo, limpo, limpo] \\
      2 & B & limpo & DIREITA & [sujo, limpo, limpo] \\
      3 & C & limpo & ESQUERDA & [sujo, limpo, limpo] \\
      \dots & \dots & \dots & \dots & \dots \\
    \end{tabular}

    O agente \textbf{nunca alcança o cômodo A} porque sua lógica só vai de A→B→C→B→C→B→C\ldots
    Ele fica preso no ciclo B↔C. Para alcançar A, seria preciso uma regra adicional
    (ex.: "se estou em C e ambos A e B estão sujos, vá para B").

    \textbf{Resolução:} O bug está na lógica de movimento: quando a posição é B ou C,
    o agente nunca se move para A. A condição \texttt{posicao == "A"} nunca é verdadeira
    se o agente não consegue chegar a A.

    \errocomum{Assumir que o agente "dá a volta". O ambiente é uma linha, não um ciclo —
    de C só se vai para B.}

    \competencia{Implementar o agente aspirador e prever seu comportamento
    com 3 cômodos.}
  }
\end{questao}

\begin{questao}{5}{●●}
  O motor de encadeamento progressivo abaixo recebe uma base de fatos inicial
  e um conjunto de regras. Execute o algoritmo \textbf{passo a passo} com os
  dados fornecidos e indique o conjunto final de fatos.

  \begin{verbatim}
  fatos = {"chovendo", "temperatura < 15"}
  regras = [
    ({"chovendo", "temperatura < 15"}, "risco de gelo"),
    ({"risco de gelo"},              "fechar pista"),
    ({"chovendo"},                   "pista molhada"),
  ]
  \end{verbatim}

  \sol{
    \textbf{Resposta:} fatos finais = \{\texttt{"chovendo"}, \texttt{"temperatura < 15"},
    \texttt{"risco de gelo"}, \texttt{"fechar pista"}, \texttt{"pista molhada"}\}

    \textbf{Resolução:}
    \begin{enumerate}[nosep]
      \item Iteração 1: chovendo + temperatura < 15 → adiciona \texttt{"risco de gelo"}.
      \item Iteração 2: risco de gelo → adiciona \texttt{"fechar pista"}.
      \item Iteração 3: chovendo → adiciona \texttt{"pista molhada"}.
      \item Iteração 4: nenhuma regra nova dispara → fim.
    \end{enumerate}

    \errocomum{Esquecer que as regras disparam em qualquer ordem — o algoritmo testa TODAS
    as regras a cada iteração, não para na primeira que dispara.}

    \competencia{Explicar o funcionamento de um motor de inferência e por que regras
    não escalam.}
  }
\end{questao}

\begin{questao}{6}{●●}
  Em 2023, um LLM foi capaz de passar no exame da ordem dos advogados (Bar Exam)
  dos EUA sem ter recebido uma única regra explícita sobre direito. Usando os
  conceitos da Semana 1, explique \textbf{por que isso é um fato tão significativo}
  na história da IA. Se não obteve uma resposta completa, use o roteiro:
  (1) como um sistema especialista resolveria o Bar Exam? (2) por que essa abordagem
  falharia? (3) o que o LLM fez de diferente?

  \sol{
    \textbf{Resposta:} Um sistema especialista exigiria que engenheiros de conhecimento
    codificassem milhares de regras jurídicas à mão — um gargalo intransponível. O LLM
    não recebeu regra nenhuma: aprendeu padrões estatísticos de centenas de gigabytes
    de texto jurídico e geral. Isso é a Bitter Lesson em ação: conhecimento aprendido
    de dados venceu conhecimento artesanal — em um domínio (direito) que se supunha
    exigir raciocínio simbólico explícito.

    \errocomum{Dizer que o LLM "decorou" as respostas. LLMs não têm memória explícita do
    corpus de treino — eles aprendem padrões estatísticos, não fazem lookup de frases.}

    \competencia{Enunciar a Bitter Lesson e dar exemplos que a corroboram.}
  }
\end{questao}

% ════════════════ NÍVEL ●●● (ANÁLISE) ════════════════
\section*{Nível ●●● \quad Análise}

\begin{questao}{7}{●●●}
  Considere a afirmação de Sutton (2019):

  \begin{quote}
  \textit{``Métodos gerais que alavancam computação são, em última instância,
  os mais efetivos — e por uma grande margem. A maioria do esforço em IA é
  gasto tentando construir conhecimento em agentes; isso ajuda no curto prazo,
  mas no longo prazo atinge um platô.''}
  \end{quote}

  \begin{enumerate}[label=(\alph*),nosep]
    \item Por que conhecimento artesanal \textit{necessariamente} atinge um platô,
    enquanto métodos que aprendem de dados não? (Dica: pense no que acontece quando
    você dobra a quantidade de dados ou de computação em cada caso.)
    \item Dê um exemplo \textbf{concreto} de um sistema que atingiu esse platô e
    explique \textit{o que especificamente} o limitou.
  \end{enumerate}

  \sol{
    \textbf{Respostas:}
    (a) Conhecimento artesanal atinge um platô porque sua melhoria depende de \textbf{humanos
    adicionando mais regras} — um processo que é linear, caro e sujeito a contradições.
    Métodos que aprendem melhoram quando você adiciona mais dados ou computação, e esse
    processo é (em princípio) ilimitado: mais dados → melhores gradientes → melhor modelo.
    Não há um humano no loop ditando o teto.

    (b) Sistemas especialistas médicos (ex.: MYCIN, anos 80). Eram melhores que médicos
    juniores em diagnósticos específicos, mas cada nova doença exigia novas regras escritas
    por especialistas. Quando o número de regras passou de centenas, as interações entre
    elas tornaram-se imprevisíveis. O sistema não melhorava com mais dados de pacientes —
    só com mais regras escritas à mão. Enquanto isso, métodos de ML atuais melhoram
    automaticamente com mais exames, prontuários e artigos.

    \errocomum{Em (a), dizer que conhecimento artesanal atinge platô "porque humanos
    são limitados". A resposta precisa identificar o mecanismo: a taxa de melhoria é
    ditada pela velocidade com que humanos escrevem regras, não pela velocidade com
    que computação cresce.}

    \competencia{Enunciar a Bitter Lesson e explicar o mecanismo por trás dela.}
  }
\end{questao}

\begin{questao}{8}{●●●}
  O efeito ELIZA — atribuir inteligência onde há apenas casamento de padrões —
  foi documentado em 1966. Ele persiste em 2025 com LLMs?

  \begin{enumerate}[label=(\alph*),nosep]
    \item Aponte \textbf{uma semelhança} entre ELIZA e um LLM moderno que torna
    o efeito ELIZA relevante hoje.
    \item Aponte \textbf{uma diferença fundamental} que torna a comparação enganosa.
    \item Em no máximo 4 linhas: um LLM é ``só casamento de padrões''? Justifique.
  \end{enumerate}

  \sol{
    \textbf{Respostas:}
    (a) Ambos produzem texto fluente sem garantia de compreensão do conteúdo semântico.
    Um LLM pode gerar uma explicação convincente de um conceito que ele não ``entende''
    no sentido humano — assim como ELIZA simulava empatia sem sentir nada. O efeito
    ELIZA ocorre quando o usuário projeta compreensão onde há apenas correlação estatística.

    (b) ELIZA usava algumas dezenas de regras de reescrita escritas à mão. Um LLM aprende
    de centenas de bilhões de parâmetros ajustados sobre trilhões de tokens. A escala e
    a origem do conhecimento (artesanal vs. aprendido) são de natureza completamente diferente.

    (c) Não. Reduzir um LLM a ``só casamento de padrões'' é como reduzir um cérebro a
    ``só sinapses disparando'' — tecnicamente verdadeiro, mas enganoso. A questão relevante
    não é se o mecanismo básico é simples (é: multiplicação de matrizes + não-linearidades),
    mas se a escala produz comportamentos emergentes que merecem ser chamados de raciocínio.
    Esta é uma pergunta em aberto na pesquisa em 2025.

    \errocomum{Responder (c) com um simples ``sim'' ou ``não''. A pergunta exige distinguir
    mecanismo de comportamento — o que o sistema \textit{faz} vs. como ele \textit{faz}.}

    \competencia{Classificar sistemas de IA e analisar criticamente alegações de inteligência.}
  }
\end{questao}

% ════════════════ NÍVEL ★ (SÍNTESE) ════════════════
\section*{Nível ★ \quad Síntese}

\begin{questao}{9}{★}
  Projete um experimento mental: você herda um sistema especialista médico com
  10.000 regras, desenvolvido ao longo de 15 anos por uma equipe de 20 engenheiros
  de conhecimento. O sistema tem 92\% de acurácia em diagnóstico.

  \begin{enumerate}[label=(\alph*),nosep]
    \item Identifique \textbf{três} cenários em que este sistema falharia de forma
    sistemática que um modelo de ML treinado sobre os mesmos dados não falharia.
    \item Para cada cenário, explique \textbf{por que} a natureza artesanal das
    regras é a causa raiz da falha.
    \item Proponha uma arquitetura híbrida (simbólico + aprendizado) que combine
    o melhor dos dois mundos para este domínio. Descreva em no máximo 6 linhas.
  \end{enumerate}

  \sol{
    \textbf{Respostas:}
    (a) e (b):
    \begin{enumerate}[nosep]
      \item \textbf{Nova doença:} o SE falha porque nenhuma regra cobre o fenômeno novo.
      O ML pode detectar que os sintomas formam um cluster distinto nos dados e sinalizar
      uma anomalia — mesmo sem rótulos. Causa raiz: SE só reconhece o que foi explicitamente
      programado.
      \item \textbf{Interação entre condições:} um paciente com diabetes \textit{e} COVID
      pode ativar regras contraditórias (uma recomenda corticoides, outra contraindica).
      O SE não tem como ponderar evidências conflitantes. O ML aprende pesos relativos
      dos features. Causa raiz: regras são discretas e não-compostas; não modelam interações.
      \item \textbf{Deriva de distribuição:} se a população de pacientes mudar (ex.:
      envelhecimento), a acurácia do SE cai porque as regras foram escritas para a
      distribuição antiga. O ML pode ser retreinado com dados novos. Causa raiz: SE
      não se adapta sem reengenharia manual.
    \end{enumerate}

    (c) \textbf{Arquitetura híbrida:} usar um modelo de ML (ex.: gradient-boosted trees
    ou rede neural) para produzir o diagnóstico base e a estimativa de confiança. As
    10.000 regras viram um sistema de \textbf{verificação}: antes de emitir o diagnóstico,
    o sistema checa se ele viola alguma regra de segurança crítica (ex.: ``nunca prescrever
    droga X para grávidas''). Se violar, escala para um médico humano. O ML lida com a
    generalização; as regras lidam com as restrições rígidas de segurança.

    \errocomum{Propor um híbrido em que as regras \textit{guiam} o aprendizado. Isso
    reintroduz o gargalo: as regras limitam o que o modelo pode aprender. O melhor
    híbrido usa regras como \textit{restrições de segurança}, não como indutores de viés.}

    \competencia{Sintetizar os conceitos de IA simbólica e aprendizado em uma arquitetura
    fundamentada.}
  }
\end{questao}

\begin{questao}{10}{★ 🌉}
  O treinamento do GPT usa a mesma estrutura de loop que o agente aspirador da Aula B:
  \textbf{percepção → decisão → ação}. A diferença está na complexidade de \texttt{decidir()}.

  \begin{enumerate}[label=(\alph*),nosep]
    \item No aspirador, \texttt{decidir()} é uma função com 4 regras \texttt{if}.
    No GPT, \texttt{decidir()} é uma rede neural com centenas de bilhões de
    parâmetros. \textbf{Por que}, conceitualmente, a mesma interface de 4 linhas
    funciona para ambos? O que isso revela sobre o paradigma de agentes?
    \item A Bitter Lesson argumenta que conhecimento artesanal atinge um platô.
    O prompt engineering — escrever instruções cuidadosas para um LLM — é uma
    forma de conhecimento artesanal? Justifique em até 5 linhas.
  \end{enumerate}

  \sol{
    \textbf{Respostas:}
    (a) A interface \texttt{obs → decidir(obs) → ação} funciona para ambos porque ela
    abstrai \textbf{o que} o agente faz de \textbf{como} ele decide. O paradigma de
    agentes é uma \textit{casca}: define a assinatura da decisão, não sua implementação.
    O que muda entre o aspirador e o GPT é apenas a função \texttt{decidir()} — e essa
    função pode ser desde 4 ifs até uma rede de 175 bilhões de parâmetros. Isso revela
    que o paradigma de agentes é \textbf{universal}: ele unifica toda a IA sob uma
    única abstração.

    (b) Sim, prompt engineering é conhecimento artesanal — mas de um tipo novo. É
    artesanal porque um humano está manualmente ajustando a entrada para melhorar
    a saída, sem que o modelo aprenda com esse ajuste. Mas difere do artesanal
    clássico porque (1) o modelo subjacente foi aprendido de dados, não programado,
    e (2) o custo de ajustar um prompt é ordens de grandeza menor que o de escrever
    regras. Dito isso, a Bitter Lesson sugere que, no longo prazo, métodos que
    \textit{aprendem} a fazer prompting (auto-prompting, RLHF) substituirão o
    prompting manual — e já há evidências disso.

    \errocomum{Em (b), responder simplesmente ``não'' porque o modelo foi treinado com ML.
    O prompt em si é conhecimento artesanal injetado \textit{na inferência}, não no treino.}

    \competencia{Reconhecer o loop percepção→decisão→ação como a interface universal
    da IA e aplicar a Bitter Lesson a desenvolvimentos recentes.}
  }
\end{questao}

\vspace{1cm}
\hrule
\vspace{0.3cm}
\begin{center}
  {\footnotesize Lista S01 · GBC063 · Prof. Marcelo Keese Albertini · FACOM/UFU}
  \ifdefined\GABARITO
    \ifnum\GABARITO=1
      \\[0.2cm] {\footnotesize\color{red}\textbf{GABARITO — USO EXCLUSIVO DO PROFESSOR}}
    \fi
  \fi
\end{center}

\end{document}
"
% ═══════════════════════════════════════════════════════════════

\documentclass[12pt,a4paper]{article}

\usepackage[utf8]{inputenc}
\usepackage[T1]{fontenc}
\usepackage[brazil]{babel}
\usepackage{geometry}
\geometry{margin=2.2cm}
\usepackage{enumitem}
\usepackage{amsmath,amssymb}
\usepackage{xcolor}
\usepackage{hyperref}
% Using standard fonts available on any TeX distribution
\usepackage{mathpazo}
\usepackage[scaled=0.9]{helvet}
\usepackage{courier}
\renewcommand{\familydefault}{\rmdefault}

% ── Conditional: gabarito on/off ──
\ifdefined\GABARITO
  \ifnum\GABARITO=1
    \newcommand{\sol}[1]{\textcolor{blue}{#1}}
    \newcommand{\errocomum}[1]{\textcolor{orange}{⚠️ Erro comum: #1}}
    \newcommand{\competencia}[1]{\textcolor{gray}{Competência: #1}}
  \else
    \newcommand{\sol}[1]{}
    \newcommand{\errocomum}[1]{}
    \newcommand{\competencia}[1]{}
  \fi
\else
  \newcommand{\sol}[1]{}
  \newcommand{\errocomum}[1]{}
  \newcommand{\competencia}[1]{}
\fi

% ── Question formatting ──
\newcommand{\nivel}[1]{\textcolor{gray}{[#1]}}
\newenvironment{questao}[2]{
  \vspace{0.5cm}
  \noindent\textbf{Q#1} \nivel{#2} \quad
}{
  \vspace{0.2cm}
}

\begin{document}

% ════════════════ HEADER ════════════════
\begin{center}
  {\sffamily\Large GBC063 · Semana 1 · Inteligência Artificial}

  \vspace{0.2cm}
  {\small Prof. Marcelo Keese Albertini · FACOM · UFU}

  \vspace{0.2cm}
  {\footnotesize Cobre as competências do resumo S01 ·
  \href{../semana01/material_introducao_ia.html}{material de estudo} ·
  \href{../semana01/resumo_s01_introducao_ia.md}{resumo}}

  \ifdefined\GABARITO
    \ifnum\GABARITO=1
      \vspace{0.3cm}
      \fbox{\textcolor{red}{\textbf{GABARITO — Documento Exclusivo do Professor}}}
    \fi
  \fi
\end{center}

\vspace{0.3cm}
\hrule
\vspace{0.5cm}

% ════════════════ NÍVEL ● (REATIVAÇÃO) ════════════════
\section*{Nível ● \quad Reativação}

\begin{questao}{1}{●}
  Defina \textbf{agente racional} segundo Russell \& Norvig e explique,
  em no máximo 3 linhas, por que esta definição — e não o Teste de Turing —
  é adotada como definição operacional do curso.

  \sol{
    \textbf{Resposta:} Um agente racional é aquele que, para cada percepção possível,
    seleciona a ação que maximiza a métrica de desempenho esperada. O curso adota esta
    definição porque ela (1) funciona com qualquer métrica de sucesso, (2) não depende
    de replicar comportamento humano, e (3) unifica desde um termostato até o AlphaGo.

    \textbf{Resolução:} Russell \& Norvig (2013), Cap. 2. A definição de agente racional
    é independente de humanidade — um agente pode ser racional e não-humano (ex.: AlphaGo
    joga Go de forma sobre-humana).

    \errocomum{Definir agente racional como "agente que age como um humano". Isso é a visão
    "agir humanamente" (Teste de Turing), não "agir racionalmente".}

    \competencia{Explicar a diferença entre IA simbólica e aprendizado de máquina —
    e definir agente racional.}
  }
\end{questao}

\begin{questao}{2}{●}
  As quatro visões clássicas da IA organizam-se em uma matriz 2×2:
  pensamento \textit{vs.} ação × desempenho humano \textit{vs.} racionalidade.
  Classifique cada um dos sistemas abaixo em \textbf{um} dos quatro quadrantes
  e justifique em 1 linha:

  \begin{enumerate}[label=(\alph*),nosep]
    \item ELIZA (Weizenbaum, 1966)
    \item Deep Blue (IBM, 1997)
    \item O agente aspirador de pó da Aula B
    \item Um modelo de linguagem treinado para prever o próximo token
  \end{enumerate}

  \sol{
    \textbf{Respostas:}
    (a) Agir humanamente — ELIZA simulava um terapeuta humano sem raciocinar.
    (b) Agir racionalmente — Deep Blue selecionava lances que maximizavam chance de vitória,
    sem imitar o pensamento de um grande mestre humano.
    (c) Agir racionalmente — o aspirador maximiza a limpeza dada sua percepção.
    (d) Pensar racionalmente (durante o treino: otimiza uma função de perda) e
    agir humanamente (na inferência: gera texto indistinguível do humano).

    \errocomum{Classificar ELIZA como "pensar humanamente". ELIZA não modelava pensamento —
    apenas casava padrões de string.}

    \competencia{Classificar qualquer sistema de IA nas quatro visões e justificar.}
  }
\end{questao}

\begin{questao}{3}{●}
  Verdadeiro ou falso. \textbf{Justifique cada resposta em até 2 linhas.}

  \begin{enumerate}[label=(\alph*),nosep]
    \item "O primeiro inverno da IA (1969–1980) foi causado exclusivamente por falta
    de poder computacional."
    \item "Redes neurais artificiais são uma ideia dos anos 2010."
  \end{enumerate}

  \sol{
    \textbf{Respostas:}
    (a) \textbf{Falso.} O inverno foi causado principalmente pelo livro \textit{Perceptrons}
    (Minsky \& Papert, 1969), que provou limitações do perceptron simples e levou o campo
    a concluir — erroneamente — que redes neurais eram um beco sem saída. A falta de
    computação contribuiu, mas o fator decisivo foi a perda de credibilidade e financiamento.

    (b) \textbf{Falso.} O primeiro modelo de neurônio artificial é de McCulloch \& Pitts (1943).
    O perceptron é de Rosenblatt (1959). O backpropagation que as treina é de 1986
    (Rumelhart, Hinton \& Williams). As ideias fundamentais têm 40–80 anos; o que mudou
    em 2012 foi a escala (dados + GPUs).

    \errocomum{Para (a): atribuir o inverno só à falta de hardware. O fator principal foi
    sociológico — a comunidade perdeu confiança nas redes neurais.}
    \errocomum{Para (b): achar que deep learning é uma invenção recente. As ideias são antigas;
    o que é novo é a escala.}

    \competencia{Enunciar a Bitter Lesson e dar exemplos históricos.}
  }
\end{questao}

% ════════════════ NÍVEL ●● (EXECUÇÃO) ════════════════
\section*{Nível ●● \quad Execução}

\begin{questao}{4}{●●}
  Considere o seguinte ambiente de aspirador de pó com \textbf{três} cômodos
  em linha (A–B–C). O agente começa na posição B. O cômodo A está sujo; B e C
  estão limpos. O agente segue exatamente a lógica abaixo:

  \begin{verbatim}
  def agente_aspirador(percepcao):
      posicao, sujo = percepcao
      if sujo:        return "ASPIRAR"
      if posicao == "A": return "DIREITA"
      if posicao == "B": return "DIREITA"
      return "ESQUERDA"
  \end{verbatim}

  Preencha a tabela de execução para os \textbf{6 primeiros passos}, indicando:
  posição atual, sensor (sujo/limpo), ação tomada, e estado dos cômodos após a ação.
  O agente limpa todos os cômodos? Em quantos passos?

  \sol{
    \textbf{Resposta:}

    \begin{tabular}{c|c|c|c|c}
      Passo & Posição & Sensor & Ação & Cômodos [A,B,C] \\
      \hline
      0 & B & limpo & DIREITA & [sujo, limpo, limpo] \\
      1 & C & limpo & ESQUERDA & [sujo, limpo, limpo] \\
      2 & B & limpo & DIREITA & [sujo, limpo, limpo] \\
      3 & C & limpo & ESQUERDA & [sujo, limpo, limpo] \\
      \dots & \dots & \dots & \dots & \dots \\
    \end{tabular}

    O agente \textbf{nunca alcança o cômodo A} porque sua lógica só vai de A→B→C→B→C→B→C\ldots
    Ele fica preso no ciclo B↔C. Para alcançar A, seria preciso uma regra adicional
    (ex.: "se estou em C e ambos A e B estão sujos, vá para B").

    \textbf{Resolução:} O bug está na lógica de movimento: quando a posição é B ou C,
    o agente nunca se move para A. A condição \texttt{posicao == "A"} nunca é verdadeira
    se o agente não consegue chegar a A.

    \errocomum{Assumir que o agente "dá a volta". O ambiente é uma linha, não um ciclo —
    de C só se vai para B.}

    \competencia{Implementar o agente aspirador e prever seu comportamento
    com 3 cômodos.}
  }
\end{questao}

\begin{questao}{5}{●●}
  O motor de encadeamento progressivo abaixo recebe uma base de fatos inicial
  e um conjunto de regras. Execute o algoritmo \textbf{passo a passo} com os
  dados fornecidos e indique o conjunto final de fatos.

  \begin{verbatim}
  fatos = {"chovendo", "temperatura < 15"}
  regras = [
    ({"chovendo", "temperatura < 15"}, "risco de gelo"),
    ({"risco de gelo"},              "fechar pista"),
    ({"chovendo"},                   "pista molhada"),
  ]
  \end{verbatim}

  \sol{
    \textbf{Resposta:} fatos finais = \{\texttt{"chovendo"}, \texttt{"temperatura < 15"},
    \texttt{"risco de gelo"}, \texttt{"fechar pista"}, \texttt{"pista molhada"}\}

    \textbf{Resolução:}
    \begin{enumerate}[nosep]
      \item Iteração 1: chovendo + temperatura < 15 → adiciona \texttt{"risco de gelo"}.
      \item Iteração 2: risco de gelo → adiciona \texttt{"fechar pista"}.
      \item Iteração 3: chovendo → adiciona \texttt{"pista molhada"}.
      \item Iteração 4: nenhuma regra nova dispara → fim.
    \end{enumerate}

    \errocomum{Esquecer que as regras disparam em qualquer ordem — o algoritmo testa TODAS
    as regras a cada iteração, não para na primeira que dispara.}

    \competencia{Explicar o funcionamento de um motor de inferência e por que regras
    não escalam.}
  }
\end{questao}

\begin{questao}{6}{●●}
  Em 2023, um LLM foi capaz de passar no exame da ordem dos advogados (Bar Exam)
  dos EUA sem ter recebido uma única regra explícita sobre direito. Usando os
  conceitos da Semana 1, explique \textbf{por que isso é um fato tão significativo}
  na história da IA. Se não obteve uma resposta completa, use o roteiro:
  (1) como um sistema especialista resolveria o Bar Exam? (2) por que essa abordagem
  falharia? (3) o que o LLM fez de diferente?

  \sol{
    \textbf{Resposta:} Um sistema especialista exigiria que engenheiros de conhecimento
    codificassem milhares de regras jurídicas à mão — um gargalo intransponível. O LLM
    não recebeu regra nenhuma: aprendeu padrões estatísticos de centenas de gigabytes
    de texto jurídico e geral. Isso é a Bitter Lesson em ação: conhecimento aprendido
    de dados venceu conhecimento artesanal — em um domínio (direito) que se supunha
    exigir raciocínio simbólico explícito.

    \errocomum{Dizer que o LLM "decorou" as respostas. LLMs não têm memória explícita do
    corpus de treino — eles aprendem padrões estatísticos, não fazem lookup de frases.}

    \competencia{Enunciar a Bitter Lesson e dar exemplos que a corroboram.}
  }
\end{questao}

% ════════════════ NÍVEL ●●● (ANÁLISE) ════════════════
\section*{Nível ●●● \quad Análise}

\begin{questao}{7}{●●●}
  Considere a afirmação de Sutton (2019):

  \begin{quote}
  \textit{``Métodos gerais que alavancam computação são, em última instância,
  os mais efetivos — e por uma grande margem. A maioria do esforço em IA é
  gasto tentando construir conhecimento em agentes; isso ajuda no curto prazo,
  mas no longo prazo atinge um platô.''}
  \end{quote}

  \begin{enumerate}[label=(\alph*),nosep]
    \item Por que conhecimento artesanal \textit{necessariamente} atinge um platô,
    enquanto métodos que aprendem de dados não? (Dica: pense no que acontece quando
    você dobra a quantidade de dados ou de computação em cada caso.)
    \item Dê um exemplo \textbf{concreto} de um sistema que atingiu esse platô e
    explique \textit{o que especificamente} o limitou.
  \end{enumerate}

  \sol{
    \textbf{Respostas:}
    (a) Conhecimento artesanal atinge um platô porque sua melhoria depende de \textbf{humanos
    adicionando mais regras} — um processo que é linear, caro e sujeito a contradições.
    Métodos que aprendem melhoram quando você adiciona mais dados ou computação, e esse
    processo é (em princípio) ilimitado: mais dados → melhores gradientes → melhor modelo.
    Não há um humano no loop ditando o teto.

    (b) Sistemas especialistas médicos (ex.: MYCIN, anos 80). Eram melhores que médicos
    juniores em diagnósticos específicos, mas cada nova doença exigia novas regras escritas
    por especialistas. Quando o número de regras passou de centenas, as interações entre
    elas tornaram-se imprevisíveis. O sistema não melhorava com mais dados de pacientes —
    só com mais regras escritas à mão. Enquanto isso, métodos de ML atuais melhoram
    automaticamente com mais exames, prontuários e artigos.

    \errocomum{Em (a), dizer que conhecimento artesanal atinge platô "porque humanos
    são limitados". A resposta precisa identificar o mecanismo: a taxa de melhoria é
    ditada pela velocidade com que humanos escrevem regras, não pela velocidade com
    que computação cresce.}

    \competencia{Enunciar a Bitter Lesson e explicar o mecanismo por trás dela.}
  }
\end{questao}

\begin{questao}{8}{●●●}
  O efeito ELIZA — atribuir inteligência onde há apenas casamento de padrões —
  foi documentado em 1966. Ele persiste em 2025 com LLMs?

  \begin{enumerate}[label=(\alph*),nosep]
    \item Aponte \textbf{uma semelhança} entre ELIZA e um LLM moderno que torna
    o efeito ELIZA relevante hoje.
    \item Aponte \textbf{uma diferença fundamental} que torna a comparação enganosa.
    \item Em no máximo 4 linhas: um LLM é ``só casamento de padrões''? Justifique.
  \end{enumerate}

  \sol{
    \textbf{Respostas:}
    (a) Ambos produzem texto fluente sem garantia de compreensão do conteúdo semântico.
    Um LLM pode gerar uma explicação convincente de um conceito que ele não ``entende''
    no sentido humano — assim como ELIZA simulava empatia sem sentir nada. O efeito
    ELIZA ocorre quando o usuário projeta compreensão onde há apenas correlação estatística.

    (b) ELIZA usava algumas dezenas de regras de reescrita escritas à mão. Um LLM aprende
    de centenas de bilhões de parâmetros ajustados sobre trilhões de tokens. A escala e
    a origem do conhecimento (artesanal vs. aprendido) são de natureza completamente diferente.

    (c) Não. Reduzir um LLM a ``só casamento de padrões'' é como reduzir um cérebro a
    ``só sinapses disparando'' — tecnicamente verdadeiro, mas enganoso. A questão relevante
    não é se o mecanismo básico é simples (é: multiplicação de matrizes + não-linearidades),
    mas se a escala produz comportamentos emergentes que merecem ser chamados de raciocínio.
    Esta é uma pergunta em aberto na pesquisa em 2025.

    \errocomum{Responder (c) com um simples ``sim'' ou ``não''. A pergunta exige distinguir
    mecanismo de comportamento — o que o sistema \textit{faz} vs. como ele \textit{faz}.}

    \competencia{Classificar sistemas de IA e analisar criticamente alegações de inteligência.}
  }
\end{questao}

% ════════════════ NÍVEL ★ (SÍNTESE) ════════════════
\section*{Nível ★ \quad Síntese}

\begin{questao}{9}{★}
  Projete um experimento mental: você herda um sistema especialista médico com
  10.000 regras, desenvolvido ao longo de 15 anos por uma equipe de 20 engenheiros
  de conhecimento. O sistema tem 92\% de acurácia em diagnóstico.

  \begin{enumerate}[label=(\alph*),nosep]
    \item Identifique \textbf{três} cenários em que este sistema falharia de forma
    sistemática que um modelo de ML treinado sobre os mesmos dados não falharia.
    \item Para cada cenário, explique \textbf{por que} a natureza artesanal das
    regras é a causa raiz da falha.
    \item Proponha uma arquitetura híbrida (simbólico + aprendizado) que combine
    o melhor dos dois mundos para este domínio. Descreva em no máximo 6 linhas.
  \end{enumerate}

  \sol{
    \textbf{Respostas:}
    (a) e (b):
    \begin{enumerate}[nosep]
      \item \textbf{Nova doença:} o SE falha porque nenhuma regra cobre o fenômeno novo.
      O ML pode detectar que os sintomas formam um cluster distinto nos dados e sinalizar
      uma anomalia — mesmo sem rótulos. Causa raiz: SE só reconhece o que foi explicitamente
      programado.
      \item \textbf{Interação entre condições:} um paciente com diabetes \textit{e} COVID
      pode ativar regras contraditórias (uma recomenda corticoides, outra contraindica).
      O SE não tem como ponderar evidências conflitantes. O ML aprende pesos relativos
      dos features. Causa raiz: regras são discretas e não-compostas; não modelam interações.
      \item \textbf{Deriva de distribuição:} se a população de pacientes mudar (ex.:
      envelhecimento), a acurácia do SE cai porque as regras foram escritas para a
      distribuição antiga. O ML pode ser retreinado com dados novos. Causa raiz: SE
      não se adapta sem reengenharia manual.
    \end{enumerate}

    (c) \textbf{Arquitetura híbrida:} usar um modelo de ML (ex.: gradient-boosted trees
    ou rede neural) para produzir o diagnóstico base e a estimativa de confiança. As
    10.000 regras viram um sistema de \textbf{verificação}: antes de emitir o diagnóstico,
    o sistema checa se ele viola alguma regra de segurança crítica (ex.: ``nunca prescrever
    droga X para grávidas''). Se violar, escala para um médico humano. O ML lida com a
    generalização; as regras lidam com as restrições rígidas de segurança.

    \errocomum{Propor um híbrido em que as regras \textit{guiam} o aprendizado. Isso
    reintroduz o gargalo: as regras limitam o que o modelo pode aprender. O melhor
    híbrido usa regras como \textit{restrições de segurança}, não como indutores de viés.}

    \competencia{Sintetizar os conceitos de IA simbólica e aprendizado em uma arquitetura
    fundamentada.}
  }
\end{questao}

\begin{questao}{10}{★ 🌉}
  O treinamento do GPT usa a mesma estrutura de loop que o agente aspirador da Aula B:
  \textbf{percepção → decisão → ação}. A diferença está na complexidade de \texttt{decidir()}.

  \begin{enumerate}[label=(\alph*),nosep]
    \item No aspirador, \texttt{decidir()} é uma função com 4 regras \texttt{if}.
    No GPT, \texttt{decidir()} é uma rede neural com centenas de bilhões de
    parâmetros. \textbf{Por que}, conceitualmente, a mesma interface de 4 linhas
    funciona para ambos? O que isso revela sobre o paradigma de agentes?
    \item A Bitter Lesson argumenta que conhecimento artesanal atinge um platô.
    O prompt engineering — escrever instruções cuidadosas para um LLM — é uma
    forma de conhecimento artesanal? Justifique em até 5 linhas.
  \end{enumerate}

  \sol{
    \textbf{Respostas:}
    (a) A interface \texttt{obs → decidir(obs) → ação} funciona para ambos porque ela
    abstrai \textbf{o que} o agente faz de \textbf{como} ele decide. O paradigma de
    agentes é uma \textit{casca}: define a assinatura da decisão, não sua implementação.
    O que muda entre o aspirador e o GPT é apenas a função \texttt{decidir()} — e essa
    função pode ser desde 4 ifs até uma rede de 175 bilhões de parâmetros. Isso revela
    que o paradigma de agentes é \textbf{universal}: ele unifica toda a IA sob uma
    única abstração.

    (b) Sim, prompt engineering é conhecimento artesanal — mas de um tipo novo. É
    artesanal porque um humano está manualmente ajustando a entrada para melhorar
    a saída, sem que o modelo aprenda com esse ajuste. Mas difere do artesanal
    clássico porque (1) o modelo subjacente foi aprendido de dados, não programado,
    e (2) o custo de ajustar um prompt é ordens de grandeza menor que o de escrever
    regras. Dito isso, a Bitter Lesson sugere que, no longo prazo, métodos que
    \textit{aprendem} a fazer prompting (auto-prompting, RLHF) substituirão o
    prompting manual — e já há evidências disso.

    \errocomum{Em (b), responder simplesmente ``não'' porque o modelo foi treinado com ML.
    O prompt em si é conhecimento artesanal injetado \textit{na inferência}, não no treino.}

    \competencia{Reconhecer o loop percepção→decisão→ação como a interface universal
    da IA e aplicar a Bitter Lesson a desenvolvimentos recentes.}
  }
\end{questao}

\vspace{1cm}
\hrule
\vspace{0.3cm}
\begin{center}
  {\footnotesize Lista S01 · GBC063 · Prof. Marcelo Keese Albertini · FACOM/UFU}
  \ifdefined\GABARITO
    \ifnum\GABARITO=1
      \\[0.2cm] {\footnotesize\color{red}\textbf{GABARITO — USO EXCLUSIVO DO PROFESSOR}}
    \fi
  \fi
\end{center}

\end{document}
"
%   Gabarito:  pdflatex -jobname=gabarito_lista_s01_introducao_ia "\def\GABARITO{1}% ═══════════════════════════════════════════════════════════════
% Lista S01 — O que é Inteligência Artificial?
% Prof. Marcelo Keese Albertini · FACOM · UFU · GBC063
%
% Compilação:
%   Questões:  pdflatex -jobname=lista_s01_introducao_ia "\def\GABARITO{0}% ═══════════════════════════════════════════════════════════════
% Lista S01 — O que é Inteligência Artificial?
% Prof. Marcelo Keese Albertini · FACOM · UFU · GBC063
%
% Compilação:
%   Questões:  pdflatex -jobname=lista_s01_introducao_ia "\def\GABARITO{0}\input{lista_s01_introducao_ia.tex}"
%   Gabarito:  pdflatex -jobname=gabarito_lista_s01_introducao_ia "\def\GABARITO{1}\input{lista_s01_introducao_ia.tex}"
% ═══════════════════════════════════════════════════════════════

\documentclass[12pt,a4paper]{article}

\usepackage[utf8]{inputenc}
\usepackage[T1]{fontenc}
\usepackage[brazil]{babel}
\usepackage{geometry}
\geometry{margin=2.2cm}
\usepackage{enumitem}
\usepackage{amsmath,amssymb}
\usepackage{xcolor}
\usepackage{hyperref}
% Using standard fonts available on any TeX distribution
\usepackage{mathpazo}
\usepackage[scaled=0.9]{helvet}
\usepackage{courier}
\renewcommand{\familydefault}{\rmdefault}

% ── Conditional: gabarito on/off ──
\ifdefined\GABARITO
  \ifnum\GABARITO=1
    \newcommand{\sol}[1]{\textcolor{blue}{#1}}
    \newcommand{\errocomum}[1]{\textcolor{orange}{⚠️ Erro comum: #1}}
    \newcommand{\competencia}[1]{\textcolor{gray}{Competência: #1}}
  \else
    \newcommand{\sol}[1]{}
    \newcommand{\errocomum}[1]{}
    \newcommand{\competencia}[1]{}
  \fi
\else
  \newcommand{\sol}[1]{}
  \newcommand{\errocomum}[1]{}
  \newcommand{\competencia}[1]{}
\fi

% ── Question formatting ──
\newcommand{\nivel}[1]{\textcolor{gray}{[#1]}}
\newenvironment{questao}[2]{
  \vspace{0.5cm}
  \noindent\textbf{Q#1} \nivel{#2} \quad
}{
  \vspace{0.2cm}
}

\begin{document}

% ════════════════ HEADER ════════════════
\begin{center}
  {\sffamily\Large GBC063 · Semana 1 · Inteligência Artificial}

  \vspace{0.2cm}
  {\small Prof. Marcelo Keese Albertini · FACOM · UFU}

  \vspace{0.2cm}
  {\footnotesize Cobre as competências do resumo S01 ·
  \href{../semana01/material_introducao_ia.html}{material de estudo} ·
  \href{../semana01/resumo_s01_introducao_ia.md}{resumo}}

  \ifdefined\GABARITO
    \ifnum\GABARITO=1
      \vspace{0.3cm}
      \fbox{\textcolor{red}{\textbf{GABARITO — Documento Exclusivo do Professor}}}
    \fi
  \fi
\end{center}

\vspace{0.3cm}
\hrule
\vspace{0.5cm}

% ════════════════ NÍVEL ● (REATIVAÇÃO) ════════════════
\section*{Nível ● \quad Reativação}

\begin{questao}{1}{●}
  Defina \textbf{agente racional} segundo Russell \& Norvig e explique,
  em no máximo 3 linhas, por que esta definição — e não o Teste de Turing —
  é adotada como definição operacional do curso.

  \sol{
    \textbf{Resposta:} Um agente racional é aquele que, para cada percepção possível,
    seleciona a ação que maximiza a métrica de desempenho esperada. O curso adota esta
    definição porque ela (1) funciona com qualquer métrica de sucesso, (2) não depende
    de replicar comportamento humano, e (3) unifica desde um termostato até o AlphaGo.

    \textbf{Resolução:} Russell \& Norvig (2013), Cap. 2. A definição de agente racional
    é independente de humanidade — um agente pode ser racional e não-humano (ex.: AlphaGo
    joga Go de forma sobre-humana).

    \errocomum{Definir agente racional como "agente que age como um humano". Isso é a visão
    "agir humanamente" (Teste de Turing), não "agir racionalmente".}

    \competencia{Explicar a diferença entre IA simbólica e aprendizado de máquina —
    e definir agente racional.}
  }
\end{questao}

\begin{questao}{2}{●}
  As quatro visões clássicas da IA organizam-se em uma matriz 2×2:
  pensamento \textit{vs.} ação × desempenho humano \textit{vs.} racionalidade.
  Classifique cada um dos sistemas abaixo em \textbf{um} dos quatro quadrantes
  e justifique em 1 linha:

  \begin{enumerate}[label=(\alph*),nosep]
    \item ELIZA (Weizenbaum, 1966)
    \item Deep Blue (IBM, 1997)
    \item O agente aspirador de pó da Aula B
    \item Um modelo de linguagem treinado para prever o próximo token
  \end{enumerate}

  \sol{
    \textbf{Respostas:}
    (a) Agir humanamente — ELIZA simulava um terapeuta humano sem raciocinar.
    (b) Agir racionalmente — Deep Blue selecionava lances que maximizavam chance de vitória,
    sem imitar o pensamento de um grande mestre humano.
    (c) Agir racionalmente — o aspirador maximiza a limpeza dada sua percepção.
    (d) Pensar racionalmente (durante o treino: otimiza uma função de perda) e
    agir humanamente (na inferência: gera texto indistinguível do humano).

    \errocomum{Classificar ELIZA como "pensar humanamente". ELIZA não modelava pensamento —
    apenas casava padrões de string.}

    \competencia{Classificar qualquer sistema de IA nas quatro visões e justificar.}
  }
\end{questao}

\begin{questao}{3}{●}
  Verdadeiro ou falso. \textbf{Justifique cada resposta em até 2 linhas.}

  \begin{enumerate}[label=(\alph*),nosep]
    \item "O primeiro inverno da IA (1969–1980) foi causado exclusivamente por falta
    de poder computacional."
    \item "Redes neurais artificiais são uma ideia dos anos 2010."
  \end{enumerate}

  \sol{
    \textbf{Respostas:}
    (a) \textbf{Falso.} O inverno foi causado principalmente pelo livro \textit{Perceptrons}
    (Minsky \& Papert, 1969), que provou limitações do perceptron simples e levou o campo
    a concluir — erroneamente — que redes neurais eram um beco sem saída. A falta de
    computação contribuiu, mas o fator decisivo foi a perda de credibilidade e financiamento.

    (b) \textbf{Falso.} O primeiro modelo de neurônio artificial é de McCulloch \& Pitts (1943).
    O perceptron é de Rosenblatt (1959). O backpropagation que as treina é de 1986
    (Rumelhart, Hinton \& Williams). As ideias fundamentais têm 40–80 anos; o que mudou
    em 2012 foi a escala (dados + GPUs).

    \errocomum{Para (a): atribuir o inverno só à falta de hardware. O fator principal foi
    sociológico — a comunidade perdeu confiança nas redes neurais.}
    \errocomum{Para (b): achar que deep learning é uma invenção recente. As ideias são antigas;
    o que é novo é a escala.}

    \competencia{Enunciar a Bitter Lesson e dar exemplos históricos.}
  }
\end{questao}

% ════════════════ NÍVEL ●● (EXECUÇÃO) ════════════════
\section*{Nível ●● \quad Execução}

\begin{questao}{4}{●●}
  Considere o seguinte ambiente de aspirador de pó com \textbf{três} cômodos
  em linha (A–B–C). O agente começa na posição B. O cômodo A está sujo; B e C
  estão limpos. O agente segue exatamente a lógica abaixo:

  \begin{verbatim}
  def agente_aspirador(percepcao):
      posicao, sujo = percepcao
      if sujo:        return "ASPIRAR"
      if posicao == "A": return "DIREITA"
      if posicao == "B": return "DIREITA"
      return "ESQUERDA"
  \end{verbatim}

  Preencha a tabela de execução para os \textbf{6 primeiros passos}, indicando:
  posição atual, sensor (sujo/limpo), ação tomada, e estado dos cômodos após a ação.
  O agente limpa todos os cômodos? Em quantos passos?

  \sol{
    \textbf{Resposta:}

    \begin{tabular}{c|c|c|c|c}
      Passo & Posição & Sensor & Ação & Cômodos [A,B,C] \\
      \hline
      0 & B & limpo & DIREITA & [sujo, limpo, limpo] \\
      1 & C & limpo & ESQUERDA & [sujo, limpo, limpo] \\
      2 & B & limpo & DIREITA & [sujo, limpo, limpo] \\
      3 & C & limpo & ESQUERDA & [sujo, limpo, limpo] \\
      \dots & \dots & \dots & \dots & \dots \\
    \end{tabular}

    O agente \textbf{nunca alcança o cômodo A} porque sua lógica só vai de A→B→C→B→C→B→C\ldots
    Ele fica preso no ciclo B↔C. Para alcançar A, seria preciso uma regra adicional
    (ex.: "se estou em C e ambos A e B estão sujos, vá para B").

    \textbf{Resolução:} O bug está na lógica de movimento: quando a posição é B ou C,
    o agente nunca se move para A. A condição \texttt{posicao == "A"} nunca é verdadeira
    se o agente não consegue chegar a A.

    \errocomum{Assumir que o agente "dá a volta". O ambiente é uma linha, não um ciclo —
    de C só se vai para B.}

    \competencia{Implementar o agente aspirador e prever seu comportamento
    com 3 cômodos.}
  }
\end{questao}

\begin{questao}{5}{●●}
  O motor de encadeamento progressivo abaixo recebe uma base de fatos inicial
  e um conjunto de regras. Execute o algoritmo \textbf{passo a passo} com os
  dados fornecidos e indique o conjunto final de fatos.

  \begin{verbatim}
  fatos = {"chovendo", "temperatura < 15"}
  regras = [
    ({"chovendo", "temperatura < 15"}, "risco de gelo"),
    ({"risco de gelo"},              "fechar pista"),
    ({"chovendo"},                   "pista molhada"),
  ]
  \end{verbatim}

  \sol{
    \textbf{Resposta:} fatos finais = \{\texttt{"chovendo"}, \texttt{"temperatura < 15"},
    \texttt{"risco de gelo"}, \texttt{"fechar pista"}, \texttt{"pista molhada"}\}

    \textbf{Resolução:}
    \begin{enumerate}[nosep]
      \item Iteração 1: chovendo + temperatura < 15 → adiciona \texttt{"risco de gelo"}.
      \item Iteração 2: risco de gelo → adiciona \texttt{"fechar pista"}.
      \item Iteração 3: chovendo → adiciona \texttt{"pista molhada"}.
      \item Iteração 4: nenhuma regra nova dispara → fim.
    \end{enumerate}

    \errocomum{Esquecer que as regras disparam em qualquer ordem — o algoritmo testa TODAS
    as regras a cada iteração, não para na primeira que dispara.}

    \competencia{Explicar o funcionamento de um motor de inferência e por que regras
    não escalam.}
  }
\end{questao}

\begin{questao}{6}{●●}
  Em 2023, um LLM foi capaz de passar no exame da ordem dos advogados (Bar Exam)
  dos EUA sem ter recebido uma única regra explícita sobre direito. Usando os
  conceitos da Semana 1, explique \textbf{por que isso é um fato tão significativo}
  na história da IA. Se não obteve uma resposta completa, use o roteiro:
  (1) como um sistema especialista resolveria o Bar Exam? (2) por que essa abordagem
  falharia? (3) o que o LLM fez de diferente?

  \sol{
    \textbf{Resposta:} Um sistema especialista exigiria que engenheiros de conhecimento
    codificassem milhares de regras jurídicas à mão — um gargalo intransponível. O LLM
    não recebeu regra nenhuma: aprendeu padrões estatísticos de centenas de gigabytes
    de texto jurídico e geral. Isso é a Bitter Lesson em ação: conhecimento aprendido
    de dados venceu conhecimento artesanal — em um domínio (direito) que se supunha
    exigir raciocínio simbólico explícito.

    \errocomum{Dizer que o LLM "decorou" as respostas. LLMs não têm memória explícita do
    corpus de treino — eles aprendem padrões estatísticos, não fazem lookup de frases.}

    \competencia{Enunciar a Bitter Lesson e dar exemplos que a corroboram.}
  }
\end{questao}

% ════════════════ NÍVEL ●●● (ANÁLISE) ════════════════
\section*{Nível ●●● \quad Análise}

\begin{questao}{7}{●●●}
  Considere a afirmação de Sutton (2019):

  \begin{quote}
  \textit{``Métodos gerais que alavancam computação são, em última instância,
  os mais efetivos — e por uma grande margem. A maioria do esforço em IA é
  gasto tentando construir conhecimento em agentes; isso ajuda no curto prazo,
  mas no longo prazo atinge um platô.''}
  \end{quote}

  \begin{enumerate}[label=(\alph*),nosep]
    \item Por que conhecimento artesanal \textit{necessariamente} atinge um platô,
    enquanto métodos que aprendem de dados não? (Dica: pense no que acontece quando
    você dobra a quantidade de dados ou de computação em cada caso.)
    \item Dê um exemplo \textbf{concreto} de um sistema que atingiu esse platô e
    explique \textit{o que especificamente} o limitou.
  \end{enumerate}

  \sol{
    \textbf{Respostas:}
    (a) Conhecimento artesanal atinge um platô porque sua melhoria depende de \textbf{humanos
    adicionando mais regras} — um processo que é linear, caro e sujeito a contradições.
    Métodos que aprendem melhoram quando você adiciona mais dados ou computação, e esse
    processo é (em princípio) ilimitado: mais dados → melhores gradientes → melhor modelo.
    Não há um humano no loop ditando o teto.

    (b) Sistemas especialistas médicos (ex.: MYCIN, anos 80). Eram melhores que médicos
    juniores em diagnósticos específicos, mas cada nova doença exigia novas regras escritas
    por especialistas. Quando o número de regras passou de centenas, as interações entre
    elas tornaram-se imprevisíveis. O sistema não melhorava com mais dados de pacientes —
    só com mais regras escritas à mão. Enquanto isso, métodos de ML atuais melhoram
    automaticamente com mais exames, prontuários e artigos.

    \errocomum{Em (a), dizer que conhecimento artesanal atinge platô "porque humanos
    são limitados". A resposta precisa identificar o mecanismo: a taxa de melhoria é
    ditada pela velocidade com que humanos escrevem regras, não pela velocidade com
    que computação cresce.}

    \competencia{Enunciar a Bitter Lesson e explicar o mecanismo por trás dela.}
  }
\end{questao}

\begin{questao}{8}{●●●}
  O efeito ELIZA — atribuir inteligência onde há apenas casamento de padrões —
  foi documentado em 1966. Ele persiste em 2025 com LLMs?

  \begin{enumerate}[label=(\alph*),nosep]
    \item Aponte \textbf{uma semelhança} entre ELIZA e um LLM moderno que torna
    o efeito ELIZA relevante hoje.
    \item Aponte \textbf{uma diferença fundamental} que torna a comparação enganosa.
    \item Em no máximo 4 linhas: um LLM é ``só casamento de padrões''? Justifique.
  \end{enumerate}

  \sol{
    \textbf{Respostas:}
    (a) Ambos produzem texto fluente sem garantia de compreensão do conteúdo semântico.
    Um LLM pode gerar uma explicação convincente de um conceito que ele não ``entende''
    no sentido humano — assim como ELIZA simulava empatia sem sentir nada. O efeito
    ELIZA ocorre quando o usuário projeta compreensão onde há apenas correlação estatística.

    (b) ELIZA usava algumas dezenas de regras de reescrita escritas à mão. Um LLM aprende
    de centenas de bilhões de parâmetros ajustados sobre trilhões de tokens. A escala e
    a origem do conhecimento (artesanal vs. aprendido) são de natureza completamente diferente.

    (c) Não. Reduzir um LLM a ``só casamento de padrões'' é como reduzir um cérebro a
    ``só sinapses disparando'' — tecnicamente verdadeiro, mas enganoso. A questão relevante
    não é se o mecanismo básico é simples (é: multiplicação de matrizes + não-linearidades),
    mas se a escala produz comportamentos emergentes que merecem ser chamados de raciocínio.
    Esta é uma pergunta em aberto na pesquisa em 2025.

    \errocomum{Responder (c) com um simples ``sim'' ou ``não''. A pergunta exige distinguir
    mecanismo de comportamento — o que o sistema \textit{faz} vs. como ele \textit{faz}.}

    \competencia{Classificar sistemas de IA e analisar criticamente alegações de inteligência.}
  }
\end{questao}

% ════════════════ NÍVEL ★ (SÍNTESE) ════════════════
\section*{Nível ★ \quad Síntese}

\begin{questao}{9}{★}
  Projete um experimento mental: você herda um sistema especialista médico com
  10.000 regras, desenvolvido ao longo de 15 anos por uma equipe de 20 engenheiros
  de conhecimento. O sistema tem 92\% de acurácia em diagnóstico.

  \begin{enumerate}[label=(\alph*),nosep]
    \item Identifique \textbf{três} cenários em que este sistema falharia de forma
    sistemática que um modelo de ML treinado sobre os mesmos dados não falharia.
    \item Para cada cenário, explique \textbf{por que} a natureza artesanal das
    regras é a causa raiz da falha.
    \item Proponha uma arquitetura híbrida (simbólico + aprendizado) que combine
    o melhor dos dois mundos para este domínio. Descreva em no máximo 6 linhas.
  \end{enumerate}

  \sol{
    \textbf{Respostas:}
    (a) e (b):
    \begin{enumerate}[nosep]
      \item \textbf{Nova doença:} o SE falha porque nenhuma regra cobre o fenômeno novo.
      O ML pode detectar que os sintomas formam um cluster distinto nos dados e sinalizar
      uma anomalia — mesmo sem rótulos. Causa raiz: SE só reconhece o que foi explicitamente
      programado.
      \item \textbf{Interação entre condições:} um paciente com diabetes \textit{e} COVID
      pode ativar regras contraditórias (uma recomenda corticoides, outra contraindica).
      O SE não tem como ponderar evidências conflitantes. O ML aprende pesos relativos
      dos features. Causa raiz: regras são discretas e não-compostas; não modelam interações.
      \item \textbf{Deriva de distribuição:} se a população de pacientes mudar (ex.:
      envelhecimento), a acurácia do SE cai porque as regras foram escritas para a
      distribuição antiga. O ML pode ser retreinado com dados novos. Causa raiz: SE
      não se adapta sem reengenharia manual.
    \end{enumerate}

    (c) \textbf{Arquitetura híbrida:} usar um modelo de ML (ex.: gradient-boosted trees
    ou rede neural) para produzir o diagnóstico base e a estimativa de confiança. As
    10.000 regras viram um sistema de \textbf{verificação}: antes de emitir o diagnóstico,
    o sistema checa se ele viola alguma regra de segurança crítica (ex.: ``nunca prescrever
    droga X para grávidas''). Se violar, escala para um médico humano. O ML lida com a
    generalização; as regras lidam com as restrições rígidas de segurança.

    \errocomum{Propor um híbrido em que as regras \textit{guiam} o aprendizado. Isso
    reintroduz o gargalo: as regras limitam o que o modelo pode aprender. O melhor
    híbrido usa regras como \textit{restrições de segurança}, não como indutores de viés.}

    \competencia{Sintetizar os conceitos de IA simbólica e aprendizado em uma arquitetura
    fundamentada.}
  }
\end{questao}

\begin{questao}{10}{★ 🌉}
  O treinamento do GPT usa a mesma estrutura de loop que o agente aspirador da Aula B:
  \textbf{percepção → decisão → ação}. A diferença está na complexidade de \texttt{decidir()}.

  \begin{enumerate}[label=(\alph*),nosep]
    \item No aspirador, \texttt{decidir()} é uma função com 4 regras \texttt{if}.
    No GPT, \texttt{decidir()} é uma rede neural com centenas de bilhões de
    parâmetros. \textbf{Por que}, conceitualmente, a mesma interface de 4 linhas
    funciona para ambos? O que isso revela sobre o paradigma de agentes?
    \item A Bitter Lesson argumenta que conhecimento artesanal atinge um platô.
    O prompt engineering — escrever instruções cuidadosas para um LLM — é uma
    forma de conhecimento artesanal? Justifique em até 5 linhas.
  \end{enumerate}

  \sol{
    \textbf{Respostas:}
    (a) A interface \texttt{obs → decidir(obs) → ação} funciona para ambos porque ela
    abstrai \textbf{o que} o agente faz de \textbf{como} ele decide. O paradigma de
    agentes é uma \textit{casca}: define a assinatura da decisão, não sua implementação.
    O que muda entre o aspirador e o GPT é apenas a função \texttt{decidir()} — e essa
    função pode ser desde 4 ifs até uma rede de 175 bilhões de parâmetros. Isso revela
    que o paradigma de agentes é \textbf{universal}: ele unifica toda a IA sob uma
    única abstração.

    (b) Sim, prompt engineering é conhecimento artesanal — mas de um tipo novo. É
    artesanal porque um humano está manualmente ajustando a entrada para melhorar
    a saída, sem que o modelo aprenda com esse ajuste. Mas difere do artesanal
    clássico porque (1) o modelo subjacente foi aprendido de dados, não programado,
    e (2) o custo de ajustar um prompt é ordens de grandeza menor que o de escrever
    regras. Dito isso, a Bitter Lesson sugere que, no longo prazo, métodos que
    \textit{aprendem} a fazer prompting (auto-prompting, RLHF) substituirão o
    prompting manual — e já há evidências disso.

    \errocomum{Em (b), responder simplesmente ``não'' porque o modelo foi treinado com ML.
    O prompt em si é conhecimento artesanal injetado \textit{na inferência}, não no treino.}

    \competencia{Reconhecer o loop percepção→decisão→ação como a interface universal
    da IA e aplicar a Bitter Lesson a desenvolvimentos recentes.}
  }
\end{questao}

\vspace{1cm}
\hrule
\vspace{0.3cm}
\begin{center}
  {\footnotesize Lista S01 · GBC063 · Prof. Marcelo Keese Albertini · FACOM/UFU}
  \ifdefined\GABARITO
    \ifnum\GABARITO=1
      \\[0.2cm] {\footnotesize\color{red}\textbf{GABARITO — USO EXCLUSIVO DO PROFESSOR}}
    \fi
  \fi
\end{center}

\end{document}
"
%   Gabarito:  pdflatex -jobname=gabarito_lista_s01_introducao_ia "\def\GABARITO{1}% ═══════════════════════════════════════════════════════════════
% Lista S01 — O que é Inteligência Artificial?
% Prof. Marcelo Keese Albertini · FACOM · UFU · GBC063
%
% Compilação:
%   Questões:  pdflatex -jobname=lista_s01_introducao_ia "\def\GABARITO{0}\input{lista_s01_introducao_ia.tex}"
%   Gabarito:  pdflatex -jobname=gabarito_lista_s01_introducao_ia "\def\GABARITO{1}\input{lista_s01_introducao_ia.tex}"
% ═══════════════════════════════════════════════════════════════

\documentclass[12pt,a4paper]{article}

\usepackage[utf8]{inputenc}
\usepackage[T1]{fontenc}
\usepackage[brazil]{babel}
\usepackage{geometry}
\geometry{margin=2.2cm}
\usepackage{enumitem}
\usepackage{amsmath,amssymb}
\usepackage{xcolor}
\usepackage{hyperref}
% Using standard fonts available on any TeX distribution
\usepackage{mathpazo}
\usepackage[scaled=0.9]{helvet}
\usepackage{courier}
\renewcommand{\familydefault}{\rmdefault}

% ── Conditional: gabarito on/off ──
\ifdefined\GABARITO
  \ifnum\GABARITO=1
    \newcommand{\sol}[1]{\textcolor{blue}{#1}}
    \newcommand{\errocomum}[1]{\textcolor{orange}{⚠️ Erro comum: #1}}
    \newcommand{\competencia}[1]{\textcolor{gray}{Competência: #1}}
  \else
    \newcommand{\sol}[1]{}
    \newcommand{\errocomum}[1]{}
    \newcommand{\competencia}[1]{}
  \fi
\else
  \newcommand{\sol}[1]{}
  \newcommand{\errocomum}[1]{}
  \newcommand{\competencia}[1]{}
\fi

% ── Question formatting ──
\newcommand{\nivel}[1]{\textcolor{gray}{[#1]}}
\newenvironment{questao}[2]{
  \vspace{0.5cm}
  \noindent\textbf{Q#1} \nivel{#2} \quad
}{
  \vspace{0.2cm}
}

\begin{document}

% ════════════════ HEADER ════════════════
\begin{center}
  {\sffamily\Large GBC063 · Semana 1 · Inteligência Artificial}

  \vspace{0.2cm}
  {\small Prof. Marcelo Keese Albertini · FACOM · UFU}

  \vspace{0.2cm}
  {\footnotesize Cobre as competências do resumo S01 ·
  \href{../semana01/material_introducao_ia.html}{material de estudo} ·
  \href{../semana01/resumo_s01_introducao_ia.md}{resumo}}

  \ifdefined\GABARITO
    \ifnum\GABARITO=1
      \vspace{0.3cm}
      \fbox{\textcolor{red}{\textbf{GABARITO — Documento Exclusivo do Professor}}}
    \fi
  \fi
\end{center}

\vspace{0.3cm}
\hrule
\vspace{0.5cm}

% ════════════════ NÍVEL ● (REATIVAÇÃO) ════════════════
\section*{Nível ● \quad Reativação}

\begin{questao}{1}{●}
  Defina \textbf{agente racional} segundo Russell \& Norvig e explique,
  em no máximo 3 linhas, por que esta definição — e não o Teste de Turing —
  é adotada como definição operacional do curso.

  \sol{
    \textbf{Resposta:} Um agente racional é aquele que, para cada percepção possível,
    seleciona a ação que maximiza a métrica de desempenho esperada. O curso adota esta
    definição porque ela (1) funciona com qualquer métrica de sucesso, (2) não depende
    de replicar comportamento humano, e (3) unifica desde um termostato até o AlphaGo.

    \textbf{Resolução:} Russell \& Norvig (2013), Cap. 2. A definição de agente racional
    é independente de humanidade — um agente pode ser racional e não-humano (ex.: AlphaGo
    joga Go de forma sobre-humana).

    \errocomum{Definir agente racional como "agente que age como um humano". Isso é a visão
    "agir humanamente" (Teste de Turing), não "agir racionalmente".}

    \competencia{Explicar a diferença entre IA simbólica e aprendizado de máquina —
    e definir agente racional.}
  }
\end{questao}

\begin{questao}{2}{●}
  As quatro visões clássicas da IA organizam-se em uma matriz 2×2:
  pensamento \textit{vs.} ação × desempenho humano \textit{vs.} racionalidade.
  Classifique cada um dos sistemas abaixo em \textbf{um} dos quatro quadrantes
  e justifique em 1 linha:

  \begin{enumerate}[label=(\alph*),nosep]
    \item ELIZA (Weizenbaum, 1966)
    \item Deep Blue (IBM, 1997)
    \item O agente aspirador de pó da Aula B
    \item Um modelo de linguagem treinado para prever o próximo token
  \end{enumerate}

  \sol{
    \textbf{Respostas:}
    (a) Agir humanamente — ELIZA simulava um terapeuta humano sem raciocinar.
    (b) Agir racionalmente — Deep Blue selecionava lances que maximizavam chance de vitória,
    sem imitar o pensamento de um grande mestre humano.
    (c) Agir racionalmente — o aspirador maximiza a limpeza dada sua percepção.
    (d) Pensar racionalmente (durante o treino: otimiza uma função de perda) e
    agir humanamente (na inferência: gera texto indistinguível do humano).

    \errocomum{Classificar ELIZA como "pensar humanamente". ELIZA não modelava pensamento —
    apenas casava padrões de string.}

    \competencia{Classificar qualquer sistema de IA nas quatro visões e justificar.}
  }
\end{questao}

\begin{questao}{3}{●}
  Verdadeiro ou falso. \textbf{Justifique cada resposta em até 2 linhas.}

  \begin{enumerate}[label=(\alph*),nosep]
    \item "O primeiro inverno da IA (1969–1980) foi causado exclusivamente por falta
    de poder computacional."
    \item "Redes neurais artificiais são uma ideia dos anos 2010."
  \end{enumerate}

  \sol{
    \textbf{Respostas:}
    (a) \textbf{Falso.} O inverno foi causado principalmente pelo livro \textit{Perceptrons}
    (Minsky \& Papert, 1969), que provou limitações do perceptron simples e levou o campo
    a concluir — erroneamente — que redes neurais eram um beco sem saída. A falta de
    computação contribuiu, mas o fator decisivo foi a perda de credibilidade e financiamento.

    (b) \textbf{Falso.} O primeiro modelo de neurônio artificial é de McCulloch \& Pitts (1943).
    O perceptron é de Rosenblatt (1959). O backpropagation que as treina é de 1986
    (Rumelhart, Hinton \& Williams). As ideias fundamentais têm 40–80 anos; o que mudou
    em 2012 foi a escala (dados + GPUs).

    \errocomum{Para (a): atribuir o inverno só à falta de hardware. O fator principal foi
    sociológico — a comunidade perdeu confiança nas redes neurais.}
    \errocomum{Para (b): achar que deep learning é uma invenção recente. As ideias são antigas;
    o que é novo é a escala.}

    \competencia{Enunciar a Bitter Lesson e dar exemplos históricos.}
  }
\end{questao}

% ════════════════ NÍVEL ●● (EXECUÇÃO) ════════════════
\section*{Nível ●● \quad Execução}

\begin{questao}{4}{●●}
  Considere o seguinte ambiente de aspirador de pó com \textbf{três} cômodos
  em linha (A–B–C). O agente começa na posição B. O cômodo A está sujo; B e C
  estão limpos. O agente segue exatamente a lógica abaixo:

  \begin{verbatim}
  def agente_aspirador(percepcao):
      posicao, sujo = percepcao
      if sujo:        return "ASPIRAR"
      if posicao == "A": return "DIREITA"
      if posicao == "B": return "DIREITA"
      return "ESQUERDA"
  \end{verbatim}

  Preencha a tabela de execução para os \textbf{6 primeiros passos}, indicando:
  posição atual, sensor (sujo/limpo), ação tomada, e estado dos cômodos após a ação.
  O agente limpa todos os cômodos? Em quantos passos?

  \sol{
    \textbf{Resposta:}

    \begin{tabular}{c|c|c|c|c}
      Passo & Posição & Sensor & Ação & Cômodos [A,B,C] \\
      \hline
      0 & B & limpo & DIREITA & [sujo, limpo, limpo] \\
      1 & C & limpo & ESQUERDA & [sujo, limpo, limpo] \\
      2 & B & limpo & DIREITA & [sujo, limpo, limpo] \\
      3 & C & limpo & ESQUERDA & [sujo, limpo, limpo] \\
      \dots & \dots & \dots & \dots & \dots \\
    \end{tabular}

    O agente \textbf{nunca alcança o cômodo A} porque sua lógica só vai de A→B→C→B→C→B→C\ldots
    Ele fica preso no ciclo B↔C. Para alcançar A, seria preciso uma regra adicional
    (ex.: "se estou em C e ambos A e B estão sujos, vá para B").

    \textbf{Resolução:} O bug está na lógica de movimento: quando a posição é B ou C,
    o agente nunca se move para A. A condição \texttt{posicao == "A"} nunca é verdadeira
    se o agente não consegue chegar a A.

    \errocomum{Assumir que o agente "dá a volta". O ambiente é uma linha, não um ciclo —
    de C só se vai para B.}

    \competencia{Implementar o agente aspirador e prever seu comportamento
    com 3 cômodos.}
  }
\end{questao}

\begin{questao}{5}{●●}
  O motor de encadeamento progressivo abaixo recebe uma base de fatos inicial
  e um conjunto de regras. Execute o algoritmo \textbf{passo a passo} com os
  dados fornecidos e indique o conjunto final de fatos.

  \begin{verbatim}
  fatos = {"chovendo", "temperatura < 15"}
  regras = [
    ({"chovendo", "temperatura < 15"}, "risco de gelo"),
    ({"risco de gelo"},              "fechar pista"),
    ({"chovendo"},                   "pista molhada"),
  ]
  \end{verbatim}

  \sol{
    \textbf{Resposta:} fatos finais = \{\texttt{"chovendo"}, \texttt{"temperatura < 15"},
    \texttt{"risco de gelo"}, \texttt{"fechar pista"}, \texttt{"pista molhada"}\}

    \textbf{Resolução:}
    \begin{enumerate}[nosep]
      \item Iteração 1: chovendo + temperatura < 15 → adiciona \texttt{"risco de gelo"}.
      \item Iteração 2: risco de gelo → adiciona \texttt{"fechar pista"}.
      \item Iteração 3: chovendo → adiciona \texttt{"pista molhada"}.
      \item Iteração 4: nenhuma regra nova dispara → fim.
    \end{enumerate}

    \errocomum{Esquecer que as regras disparam em qualquer ordem — o algoritmo testa TODAS
    as regras a cada iteração, não para na primeira que dispara.}

    \competencia{Explicar o funcionamento de um motor de inferência e por que regras
    não escalam.}
  }
\end{questao}

\begin{questao}{6}{●●}
  Em 2023, um LLM foi capaz de passar no exame da ordem dos advogados (Bar Exam)
  dos EUA sem ter recebido uma única regra explícita sobre direito. Usando os
  conceitos da Semana 1, explique \textbf{por que isso é um fato tão significativo}
  na história da IA. Se não obteve uma resposta completa, use o roteiro:
  (1) como um sistema especialista resolveria o Bar Exam? (2) por que essa abordagem
  falharia? (3) o que o LLM fez de diferente?

  \sol{
    \textbf{Resposta:} Um sistema especialista exigiria que engenheiros de conhecimento
    codificassem milhares de regras jurídicas à mão — um gargalo intransponível. O LLM
    não recebeu regra nenhuma: aprendeu padrões estatísticos de centenas de gigabytes
    de texto jurídico e geral. Isso é a Bitter Lesson em ação: conhecimento aprendido
    de dados venceu conhecimento artesanal — em um domínio (direito) que se supunha
    exigir raciocínio simbólico explícito.

    \errocomum{Dizer que o LLM "decorou" as respostas. LLMs não têm memória explícita do
    corpus de treino — eles aprendem padrões estatísticos, não fazem lookup de frases.}

    \competencia{Enunciar a Bitter Lesson e dar exemplos que a corroboram.}
  }
\end{questao}

% ════════════════ NÍVEL ●●● (ANÁLISE) ════════════════
\section*{Nível ●●● \quad Análise}

\begin{questao}{7}{●●●}
  Considere a afirmação de Sutton (2019):

  \begin{quote}
  \textit{``Métodos gerais que alavancam computação são, em última instância,
  os mais efetivos — e por uma grande margem. A maioria do esforço em IA é
  gasto tentando construir conhecimento em agentes; isso ajuda no curto prazo,
  mas no longo prazo atinge um platô.''}
  \end{quote}

  \begin{enumerate}[label=(\alph*),nosep]
    \item Por que conhecimento artesanal \textit{necessariamente} atinge um platô,
    enquanto métodos que aprendem de dados não? (Dica: pense no que acontece quando
    você dobra a quantidade de dados ou de computação em cada caso.)
    \item Dê um exemplo \textbf{concreto} de um sistema que atingiu esse platô e
    explique \textit{o que especificamente} o limitou.
  \end{enumerate}

  \sol{
    \textbf{Respostas:}
    (a) Conhecimento artesanal atinge um platô porque sua melhoria depende de \textbf{humanos
    adicionando mais regras} — um processo que é linear, caro e sujeito a contradições.
    Métodos que aprendem melhoram quando você adiciona mais dados ou computação, e esse
    processo é (em princípio) ilimitado: mais dados → melhores gradientes → melhor modelo.
    Não há um humano no loop ditando o teto.

    (b) Sistemas especialistas médicos (ex.: MYCIN, anos 80). Eram melhores que médicos
    juniores em diagnósticos específicos, mas cada nova doença exigia novas regras escritas
    por especialistas. Quando o número de regras passou de centenas, as interações entre
    elas tornaram-se imprevisíveis. O sistema não melhorava com mais dados de pacientes —
    só com mais regras escritas à mão. Enquanto isso, métodos de ML atuais melhoram
    automaticamente com mais exames, prontuários e artigos.

    \errocomum{Em (a), dizer que conhecimento artesanal atinge platô "porque humanos
    são limitados". A resposta precisa identificar o mecanismo: a taxa de melhoria é
    ditada pela velocidade com que humanos escrevem regras, não pela velocidade com
    que computação cresce.}

    \competencia{Enunciar a Bitter Lesson e explicar o mecanismo por trás dela.}
  }
\end{questao}

\begin{questao}{8}{●●●}
  O efeito ELIZA — atribuir inteligência onde há apenas casamento de padrões —
  foi documentado em 1966. Ele persiste em 2025 com LLMs?

  \begin{enumerate}[label=(\alph*),nosep]
    \item Aponte \textbf{uma semelhança} entre ELIZA e um LLM moderno que torna
    o efeito ELIZA relevante hoje.
    \item Aponte \textbf{uma diferença fundamental} que torna a comparação enganosa.
    \item Em no máximo 4 linhas: um LLM é ``só casamento de padrões''? Justifique.
  \end{enumerate}

  \sol{
    \textbf{Respostas:}
    (a) Ambos produzem texto fluente sem garantia de compreensão do conteúdo semântico.
    Um LLM pode gerar uma explicação convincente de um conceito que ele não ``entende''
    no sentido humano — assim como ELIZA simulava empatia sem sentir nada. O efeito
    ELIZA ocorre quando o usuário projeta compreensão onde há apenas correlação estatística.

    (b) ELIZA usava algumas dezenas de regras de reescrita escritas à mão. Um LLM aprende
    de centenas de bilhões de parâmetros ajustados sobre trilhões de tokens. A escala e
    a origem do conhecimento (artesanal vs. aprendido) são de natureza completamente diferente.

    (c) Não. Reduzir um LLM a ``só casamento de padrões'' é como reduzir um cérebro a
    ``só sinapses disparando'' — tecnicamente verdadeiro, mas enganoso. A questão relevante
    não é se o mecanismo básico é simples (é: multiplicação de matrizes + não-linearidades),
    mas se a escala produz comportamentos emergentes que merecem ser chamados de raciocínio.
    Esta é uma pergunta em aberto na pesquisa em 2025.

    \errocomum{Responder (c) com um simples ``sim'' ou ``não''. A pergunta exige distinguir
    mecanismo de comportamento — o que o sistema \textit{faz} vs. como ele \textit{faz}.}

    \competencia{Classificar sistemas de IA e analisar criticamente alegações de inteligência.}
  }
\end{questao}

% ════════════════ NÍVEL ★ (SÍNTESE) ════════════════
\section*{Nível ★ \quad Síntese}

\begin{questao}{9}{★}
  Projete um experimento mental: você herda um sistema especialista médico com
  10.000 regras, desenvolvido ao longo de 15 anos por uma equipe de 20 engenheiros
  de conhecimento. O sistema tem 92\% de acurácia em diagnóstico.

  \begin{enumerate}[label=(\alph*),nosep]
    \item Identifique \textbf{três} cenários em que este sistema falharia de forma
    sistemática que um modelo de ML treinado sobre os mesmos dados não falharia.
    \item Para cada cenário, explique \textbf{por que} a natureza artesanal das
    regras é a causa raiz da falha.
    \item Proponha uma arquitetura híbrida (simbólico + aprendizado) que combine
    o melhor dos dois mundos para este domínio. Descreva em no máximo 6 linhas.
  \end{enumerate}

  \sol{
    \textbf{Respostas:}
    (a) e (b):
    \begin{enumerate}[nosep]
      \item \textbf{Nova doença:} o SE falha porque nenhuma regra cobre o fenômeno novo.
      O ML pode detectar que os sintomas formam um cluster distinto nos dados e sinalizar
      uma anomalia — mesmo sem rótulos. Causa raiz: SE só reconhece o que foi explicitamente
      programado.
      \item \textbf{Interação entre condições:} um paciente com diabetes \textit{e} COVID
      pode ativar regras contraditórias (uma recomenda corticoides, outra contraindica).
      O SE não tem como ponderar evidências conflitantes. O ML aprende pesos relativos
      dos features. Causa raiz: regras são discretas e não-compostas; não modelam interações.
      \item \textbf{Deriva de distribuição:} se a população de pacientes mudar (ex.:
      envelhecimento), a acurácia do SE cai porque as regras foram escritas para a
      distribuição antiga. O ML pode ser retreinado com dados novos. Causa raiz: SE
      não se adapta sem reengenharia manual.
    \end{enumerate}

    (c) \textbf{Arquitetura híbrida:} usar um modelo de ML (ex.: gradient-boosted trees
    ou rede neural) para produzir o diagnóstico base e a estimativa de confiança. As
    10.000 regras viram um sistema de \textbf{verificação}: antes de emitir o diagnóstico,
    o sistema checa se ele viola alguma regra de segurança crítica (ex.: ``nunca prescrever
    droga X para grávidas''). Se violar, escala para um médico humano. O ML lida com a
    generalização; as regras lidam com as restrições rígidas de segurança.

    \errocomum{Propor um híbrido em que as regras \textit{guiam} o aprendizado. Isso
    reintroduz o gargalo: as regras limitam o que o modelo pode aprender. O melhor
    híbrido usa regras como \textit{restrições de segurança}, não como indutores de viés.}

    \competencia{Sintetizar os conceitos de IA simbólica e aprendizado em uma arquitetura
    fundamentada.}
  }
\end{questao}

\begin{questao}{10}{★ 🌉}
  O treinamento do GPT usa a mesma estrutura de loop que o agente aspirador da Aula B:
  \textbf{percepção → decisão → ação}. A diferença está na complexidade de \texttt{decidir()}.

  \begin{enumerate}[label=(\alph*),nosep]
    \item No aspirador, \texttt{decidir()} é uma função com 4 regras \texttt{if}.
    No GPT, \texttt{decidir()} é uma rede neural com centenas de bilhões de
    parâmetros. \textbf{Por que}, conceitualmente, a mesma interface de 4 linhas
    funciona para ambos? O que isso revela sobre o paradigma de agentes?
    \item A Bitter Lesson argumenta que conhecimento artesanal atinge um platô.
    O prompt engineering — escrever instruções cuidadosas para um LLM — é uma
    forma de conhecimento artesanal? Justifique em até 5 linhas.
  \end{enumerate}

  \sol{
    \textbf{Respostas:}
    (a) A interface \texttt{obs → decidir(obs) → ação} funciona para ambos porque ela
    abstrai \textbf{o que} o agente faz de \textbf{como} ele decide. O paradigma de
    agentes é uma \textit{casca}: define a assinatura da decisão, não sua implementação.
    O que muda entre o aspirador e o GPT é apenas a função \texttt{decidir()} — e essa
    função pode ser desde 4 ifs até uma rede de 175 bilhões de parâmetros. Isso revela
    que o paradigma de agentes é \textbf{universal}: ele unifica toda a IA sob uma
    única abstração.

    (b) Sim, prompt engineering é conhecimento artesanal — mas de um tipo novo. É
    artesanal porque um humano está manualmente ajustando a entrada para melhorar
    a saída, sem que o modelo aprenda com esse ajuste. Mas difere do artesanal
    clássico porque (1) o modelo subjacente foi aprendido de dados, não programado,
    e (2) o custo de ajustar um prompt é ordens de grandeza menor que o de escrever
    regras. Dito isso, a Bitter Lesson sugere que, no longo prazo, métodos que
    \textit{aprendem} a fazer prompting (auto-prompting, RLHF) substituirão o
    prompting manual — e já há evidências disso.

    \errocomum{Em (b), responder simplesmente ``não'' porque o modelo foi treinado com ML.
    O prompt em si é conhecimento artesanal injetado \textit{na inferência}, não no treino.}

    \competencia{Reconhecer o loop percepção→decisão→ação como a interface universal
    da IA e aplicar a Bitter Lesson a desenvolvimentos recentes.}
  }
\end{questao}

\vspace{1cm}
\hrule
\vspace{0.3cm}
\begin{center}
  {\footnotesize Lista S01 · GBC063 · Prof. Marcelo Keese Albertini · FACOM/UFU}
  \ifdefined\GABARITO
    \ifnum\GABARITO=1
      \\[0.2cm] {\footnotesize\color{red}\textbf{GABARITO — USO EXCLUSIVO DO PROFESSOR}}
    \fi
  \fi
\end{center}

\end{document}
"
% ═══════════════════════════════════════════════════════════════

\documentclass[12pt,a4paper]{article}

\usepackage[utf8]{inputenc}
\usepackage[T1]{fontenc}
\usepackage[brazil]{babel}
\usepackage{geometry}
\geometry{margin=2.2cm}
\usepackage{enumitem}
\usepackage{amsmath,amssymb}
\usepackage{xcolor}
\usepackage{hyperref}
% Using standard fonts available on any TeX distribution
\usepackage{mathpazo}
\usepackage[scaled=0.9]{helvet}
\usepackage{courier}
\renewcommand{\familydefault}{\rmdefault}

% ── Conditional: gabarito on/off ──
\ifdefined\GABARITO
  \ifnum\GABARITO=1
    \newcommand{\sol}[1]{\textcolor{blue}{#1}}
    \newcommand{\errocomum}[1]{\textcolor{orange}{⚠️ Erro comum: #1}}
    \newcommand{\competencia}[1]{\textcolor{gray}{Competência: #1}}
  \else
    \newcommand{\sol}[1]{}
    \newcommand{\errocomum}[1]{}
    \newcommand{\competencia}[1]{}
  \fi
\else
  \newcommand{\sol}[1]{}
  \newcommand{\errocomum}[1]{}
  \newcommand{\competencia}[1]{}
\fi

% ── Question formatting ──
\newcommand{\nivel}[1]{\textcolor{gray}{[#1]}}
\newenvironment{questao}[2]{
  \vspace{0.5cm}
  \noindent\textbf{Q#1} \nivel{#2} \quad
}{
  \vspace{0.2cm}
}

\begin{document}

% ════════════════ HEADER ════════════════
\begin{center}
  {\sffamily\Large GBC063 · Semana 1 · Inteligência Artificial}

  \vspace{0.2cm}
  {\small Prof. Marcelo Keese Albertini · FACOM · UFU}

  \vspace{0.2cm}
  {\footnotesize Cobre as competências do resumo S01 ·
  \href{../semana01/material_introducao_ia.html}{material de estudo} ·
  \href{../semana01/resumo_s01_introducao_ia.md}{resumo}}

  \ifdefined\GABARITO
    \ifnum\GABARITO=1
      \vspace{0.3cm}
      \fbox{\textcolor{red}{\textbf{GABARITO — Documento Exclusivo do Professor}}}
    \fi
  \fi
\end{center}

\vspace{0.3cm}
\hrule
\vspace{0.5cm}

% ════════════════ NÍVEL ● (REATIVAÇÃO) ════════════════
\section*{Nível ● \quad Reativação}

\begin{questao}{1}{●}
  Defina \textbf{agente racional} segundo Russell \& Norvig e explique,
  em no máximo 3 linhas, por que esta definição — e não o Teste de Turing —
  é adotada como definição operacional do curso.

  \sol{
    \textbf{Resposta:} Um agente racional é aquele que, para cada percepção possível,
    seleciona a ação que maximiza a métrica de desempenho esperada. O curso adota esta
    definição porque ela (1) funciona com qualquer métrica de sucesso, (2) não depende
    de replicar comportamento humano, e (3) unifica desde um termostato até o AlphaGo.

    \textbf{Resolução:} Russell \& Norvig (2013), Cap. 2. A definição de agente racional
    é independente de humanidade — um agente pode ser racional e não-humano (ex.: AlphaGo
    joga Go de forma sobre-humana).

    \errocomum{Definir agente racional como "agente que age como um humano". Isso é a visão
    "agir humanamente" (Teste de Turing), não "agir racionalmente".}

    \competencia{Explicar a diferença entre IA simbólica e aprendizado de máquina —
    e definir agente racional.}
  }
\end{questao}

\begin{questao}{2}{●}
  As quatro visões clássicas da IA organizam-se em uma matriz 2×2:
  pensamento \textit{vs.} ação × desempenho humano \textit{vs.} racionalidade.
  Classifique cada um dos sistemas abaixo em \textbf{um} dos quatro quadrantes
  e justifique em 1 linha:

  \begin{enumerate}[label=(\alph*),nosep]
    \item ELIZA (Weizenbaum, 1966)
    \item Deep Blue (IBM, 1997)
    \item O agente aspirador de pó da Aula B
    \item Um modelo de linguagem treinado para prever o próximo token
  \end{enumerate}

  \sol{
    \textbf{Respostas:}
    (a) Agir humanamente — ELIZA simulava um terapeuta humano sem raciocinar.
    (b) Agir racionalmente — Deep Blue selecionava lances que maximizavam chance de vitória,
    sem imitar o pensamento de um grande mestre humano.
    (c) Agir racionalmente — o aspirador maximiza a limpeza dada sua percepção.
    (d) Pensar racionalmente (durante o treino: otimiza uma função de perda) e
    agir humanamente (na inferência: gera texto indistinguível do humano).

    \errocomum{Classificar ELIZA como "pensar humanamente". ELIZA não modelava pensamento —
    apenas casava padrões de string.}

    \competencia{Classificar qualquer sistema de IA nas quatro visões e justificar.}
  }
\end{questao}

\begin{questao}{3}{●}
  Verdadeiro ou falso. \textbf{Justifique cada resposta em até 2 linhas.}

  \begin{enumerate}[label=(\alph*),nosep]
    \item "O primeiro inverno da IA (1969–1980) foi causado exclusivamente por falta
    de poder computacional."
    \item "Redes neurais artificiais são uma ideia dos anos 2010."
  \end{enumerate}

  \sol{
    \textbf{Respostas:}
    (a) \textbf{Falso.} O inverno foi causado principalmente pelo livro \textit{Perceptrons}
    (Minsky \& Papert, 1969), que provou limitações do perceptron simples e levou o campo
    a concluir — erroneamente — que redes neurais eram um beco sem saída. A falta de
    computação contribuiu, mas o fator decisivo foi a perda de credibilidade e financiamento.

    (b) \textbf{Falso.} O primeiro modelo de neurônio artificial é de McCulloch \& Pitts (1943).
    O perceptron é de Rosenblatt (1959). O backpropagation que as treina é de 1986
    (Rumelhart, Hinton \& Williams). As ideias fundamentais têm 40–80 anos; o que mudou
    em 2012 foi a escala (dados + GPUs).

    \errocomum{Para (a): atribuir o inverno só à falta de hardware. O fator principal foi
    sociológico — a comunidade perdeu confiança nas redes neurais.}
    \errocomum{Para (b): achar que deep learning é uma invenção recente. As ideias são antigas;
    o que é novo é a escala.}

    \competencia{Enunciar a Bitter Lesson e dar exemplos históricos.}
  }
\end{questao}

% ════════════════ NÍVEL ●● (EXECUÇÃO) ════════════════
\section*{Nível ●● \quad Execução}

\begin{questao}{4}{●●}
  Considere o seguinte ambiente de aspirador de pó com \textbf{três} cômodos
  em linha (A–B–C). O agente começa na posição B. O cômodo A está sujo; B e C
  estão limpos. O agente segue exatamente a lógica abaixo:

  \begin{verbatim}
  def agente_aspirador(percepcao):
      posicao, sujo = percepcao
      if sujo:        return "ASPIRAR"
      if posicao == "A": return "DIREITA"
      if posicao == "B": return "DIREITA"
      return "ESQUERDA"
  \end{verbatim}

  Preencha a tabela de execução para os \textbf{6 primeiros passos}, indicando:
  posição atual, sensor (sujo/limpo), ação tomada, e estado dos cômodos após a ação.
  O agente limpa todos os cômodos? Em quantos passos?

  \sol{
    \textbf{Resposta:}

    \begin{tabular}{c|c|c|c|c}
      Passo & Posição & Sensor & Ação & Cômodos [A,B,C] \\
      \hline
      0 & B & limpo & DIREITA & [sujo, limpo, limpo] \\
      1 & C & limpo & ESQUERDA & [sujo, limpo, limpo] \\
      2 & B & limpo & DIREITA & [sujo, limpo, limpo] \\
      3 & C & limpo & ESQUERDA & [sujo, limpo, limpo] \\
      \dots & \dots & \dots & \dots & \dots \\
    \end{tabular}

    O agente \textbf{nunca alcança o cômodo A} porque sua lógica só vai de A→B→C→B→C→B→C\ldots
    Ele fica preso no ciclo B↔C. Para alcançar A, seria preciso uma regra adicional
    (ex.: "se estou em C e ambos A e B estão sujos, vá para B").

    \textbf{Resolução:} O bug está na lógica de movimento: quando a posição é B ou C,
    o agente nunca se move para A. A condição \texttt{posicao == "A"} nunca é verdadeira
    se o agente não consegue chegar a A.

    \errocomum{Assumir que o agente "dá a volta". O ambiente é uma linha, não um ciclo —
    de C só se vai para B.}

    \competencia{Implementar o agente aspirador e prever seu comportamento
    com 3 cômodos.}
  }
\end{questao}

\begin{questao}{5}{●●}
  O motor de encadeamento progressivo abaixo recebe uma base de fatos inicial
  e um conjunto de regras. Execute o algoritmo \textbf{passo a passo} com os
  dados fornecidos e indique o conjunto final de fatos.

  \begin{verbatim}
  fatos = {"chovendo", "temperatura < 15"}
  regras = [
    ({"chovendo", "temperatura < 15"}, "risco de gelo"),
    ({"risco de gelo"},              "fechar pista"),
    ({"chovendo"},                   "pista molhada"),
  ]
  \end{verbatim}

  \sol{
    \textbf{Resposta:} fatos finais = \{\texttt{"chovendo"}, \texttt{"temperatura < 15"},
    \texttt{"risco de gelo"}, \texttt{"fechar pista"}, \texttt{"pista molhada"}\}

    \textbf{Resolução:}
    \begin{enumerate}[nosep]
      \item Iteração 1: chovendo + temperatura < 15 → adiciona \texttt{"risco de gelo"}.
      \item Iteração 2: risco de gelo → adiciona \texttt{"fechar pista"}.
      \item Iteração 3: chovendo → adiciona \texttt{"pista molhada"}.
      \item Iteração 4: nenhuma regra nova dispara → fim.
    \end{enumerate}

    \errocomum{Esquecer que as regras disparam em qualquer ordem — o algoritmo testa TODAS
    as regras a cada iteração, não para na primeira que dispara.}

    \competencia{Explicar o funcionamento de um motor de inferência e por que regras
    não escalam.}
  }
\end{questao}

\begin{questao}{6}{●●}
  Em 2023, um LLM foi capaz de passar no exame da ordem dos advogados (Bar Exam)
  dos EUA sem ter recebido uma única regra explícita sobre direito. Usando os
  conceitos da Semana 1, explique \textbf{por que isso é um fato tão significativo}
  na história da IA. Se não obteve uma resposta completa, use o roteiro:
  (1) como um sistema especialista resolveria o Bar Exam? (2) por que essa abordagem
  falharia? (3) o que o LLM fez de diferente?

  \sol{
    \textbf{Resposta:} Um sistema especialista exigiria que engenheiros de conhecimento
    codificassem milhares de regras jurídicas à mão — um gargalo intransponível. O LLM
    não recebeu regra nenhuma: aprendeu padrões estatísticos de centenas de gigabytes
    de texto jurídico e geral. Isso é a Bitter Lesson em ação: conhecimento aprendido
    de dados venceu conhecimento artesanal — em um domínio (direito) que se supunha
    exigir raciocínio simbólico explícito.

    \errocomum{Dizer que o LLM "decorou" as respostas. LLMs não têm memória explícita do
    corpus de treino — eles aprendem padrões estatísticos, não fazem lookup de frases.}

    \competencia{Enunciar a Bitter Lesson e dar exemplos que a corroboram.}
  }
\end{questao}

% ════════════════ NÍVEL ●●● (ANÁLISE) ════════════════
\section*{Nível ●●● \quad Análise}

\begin{questao}{7}{●●●}
  Considere a afirmação de Sutton (2019):

  \begin{quote}
  \textit{``Métodos gerais que alavancam computação são, em última instância,
  os mais efetivos — e por uma grande margem. A maioria do esforço em IA é
  gasto tentando construir conhecimento em agentes; isso ajuda no curto prazo,
  mas no longo prazo atinge um platô.''}
  \end{quote}

  \begin{enumerate}[label=(\alph*),nosep]
    \item Por que conhecimento artesanal \textit{necessariamente} atinge um platô,
    enquanto métodos que aprendem de dados não? (Dica: pense no que acontece quando
    você dobra a quantidade de dados ou de computação em cada caso.)
    \item Dê um exemplo \textbf{concreto} de um sistema que atingiu esse platô e
    explique \textit{o que especificamente} o limitou.
  \end{enumerate}

  \sol{
    \textbf{Respostas:}
    (a) Conhecimento artesanal atinge um platô porque sua melhoria depende de \textbf{humanos
    adicionando mais regras} — um processo que é linear, caro e sujeito a contradições.
    Métodos que aprendem melhoram quando você adiciona mais dados ou computação, e esse
    processo é (em princípio) ilimitado: mais dados → melhores gradientes → melhor modelo.
    Não há um humano no loop ditando o teto.

    (b) Sistemas especialistas médicos (ex.: MYCIN, anos 80). Eram melhores que médicos
    juniores em diagnósticos específicos, mas cada nova doença exigia novas regras escritas
    por especialistas. Quando o número de regras passou de centenas, as interações entre
    elas tornaram-se imprevisíveis. O sistema não melhorava com mais dados de pacientes —
    só com mais regras escritas à mão. Enquanto isso, métodos de ML atuais melhoram
    automaticamente com mais exames, prontuários e artigos.

    \errocomum{Em (a), dizer que conhecimento artesanal atinge platô "porque humanos
    são limitados". A resposta precisa identificar o mecanismo: a taxa de melhoria é
    ditada pela velocidade com que humanos escrevem regras, não pela velocidade com
    que computação cresce.}

    \competencia{Enunciar a Bitter Lesson e explicar o mecanismo por trás dela.}
  }
\end{questao}

\begin{questao}{8}{●●●}
  O efeito ELIZA — atribuir inteligência onde há apenas casamento de padrões —
  foi documentado em 1966. Ele persiste em 2025 com LLMs?

  \begin{enumerate}[label=(\alph*),nosep]
    \item Aponte \textbf{uma semelhança} entre ELIZA e um LLM moderno que torna
    o efeito ELIZA relevante hoje.
    \item Aponte \textbf{uma diferença fundamental} que torna a comparação enganosa.
    \item Em no máximo 4 linhas: um LLM é ``só casamento de padrões''? Justifique.
  \end{enumerate}

  \sol{
    \textbf{Respostas:}
    (a) Ambos produzem texto fluente sem garantia de compreensão do conteúdo semântico.
    Um LLM pode gerar uma explicação convincente de um conceito que ele não ``entende''
    no sentido humano — assim como ELIZA simulava empatia sem sentir nada. O efeito
    ELIZA ocorre quando o usuário projeta compreensão onde há apenas correlação estatística.

    (b) ELIZA usava algumas dezenas de regras de reescrita escritas à mão. Um LLM aprende
    de centenas de bilhões de parâmetros ajustados sobre trilhões de tokens. A escala e
    a origem do conhecimento (artesanal vs. aprendido) são de natureza completamente diferente.

    (c) Não. Reduzir um LLM a ``só casamento de padrões'' é como reduzir um cérebro a
    ``só sinapses disparando'' — tecnicamente verdadeiro, mas enganoso. A questão relevante
    não é se o mecanismo básico é simples (é: multiplicação de matrizes + não-linearidades),
    mas se a escala produz comportamentos emergentes que merecem ser chamados de raciocínio.
    Esta é uma pergunta em aberto na pesquisa em 2025.

    \errocomum{Responder (c) com um simples ``sim'' ou ``não''. A pergunta exige distinguir
    mecanismo de comportamento — o que o sistema \textit{faz} vs. como ele \textit{faz}.}

    \competencia{Classificar sistemas de IA e analisar criticamente alegações de inteligência.}
  }
\end{questao}

% ════════════════ NÍVEL ★ (SÍNTESE) ════════════════
\section*{Nível ★ \quad Síntese}

\begin{questao}{9}{★}
  Projete um experimento mental: você herda um sistema especialista médico com
  10.000 regras, desenvolvido ao longo de 15 anos por uma equipe de 20 engenheiros
  de conhecimento. O sistema tem 92\% de acurácia em diagnóstico.

  \begin{enumerate}[label=(\alph*),nosep]
    \item Identifique \textbf{três} cenários em que este sistema falharia de forma
    sistemática que um modelo de ML treinado sobre os mesmos dados não falharia.
    \item Para cada cenário, explique \textbf{por que} a natureza artesanal das
    regras é a causa raiz da falha.
    \item Proponha uma arquitetura híbrida (simbólico + aprendizado) que combine
    o melhor dos dois mundos para este domínio. Descreva em no máximo 6 linhas.
  \end{enumerate}

  \sol{
    \textbf{Respostas:}
    (a) e (b):
    \begin{enumerate}[nosep]
      \item \textbf{Nova doença:} o SE falha porque nenhuma regra cobre o fenômeno novo.
      O ML pode detectar que os sintomas formam um cluster distinto nos dados e sinalizar
      uma anomalia — mesmo sem rótulos. Causa raiz: SE só reconhece o que foi explicitamente
      programado.
      \item \textbf{Interação entre condições:} um paciente com diabetes \textit{e} COVID
      pode ativar regras contraditórias (uma recomenda corticoides, outra contraindica).
      O SE não tem como ponderar evidências conflitantes. O ML aprende pesos relativos
      dos features. Causa raiz: regras são discretas e não-compostas; não modelam interações.
      \item \textbf{Deriva de distribuição:} se a população de pacientes mudar (ex.:
      envelhecimento), a acurácia do SE cai porque as regras foram escritas para a
      distribuição antiga. O ML pode ser retreinado com dados novos. Causa raiz: SE
      não se adapta sem reengenharia manual.
    \end{enumerate}

    (c) \textbf{Arquitetura híbrida:} usar um modelo de ML (ex.: gradient-boosted trees
    ou rede neural) para produzir o diagnóstico base e a estimativa de confiança. As
    10.000 regras viram um sistema de \textbf{verificação}: antes de emitir o diagnóstico,
    o sistema checa se ele viola alguma regra de segurança crítica (ex.: ``nunca prescrever
    droga X para grávidas''). Se violar, escala para um médico humano. O ML lida com a
    generalização; as regras lidam com as restrições rígidas de segurança.

    \errocomum{Propor um híbrido em que as regras \textit{guiam} o aprendizado. Isso
    reintroduz o gargalo: as regras limitam o que o modelo pode aprender. O melhor
    híbrido usa regras como \textit{restrições de segurança}, não como indutores de viés.}

    \competencia{Sintetizar os conceitos de IA simbólica e aprendizado em uma arquitetura
    fundamentada.}
  }
\end{questao}

\begin{questao}{10}{★ 🌉}
  O treinamento do GPT usa a mesma estrutura de loop que o agente aspirador da Aula B:
  \textbf{percepção → decisão → ação}. A diferença está na complexidade de \texttt{decidir()}.

  \begin{enumerate}[label=(\alph*),nosep]
    \item No aspirador, \texttt{decidir()} é uma função com 4 regras \texttt{if}.
    No GPT, \texttt{decidir()} é uma rede neural com centenas de bilhões de
    parâmetros. \textbf{Por que}, conceitualmente, a mesma interface de 4 linhas
    funciona para ambos? O que isso revela sobre o paradigma de agentes?
    \item A Bitter Lesson argumenta que conhecimento artesanal atinge um platô.
    O prompt engineering — escrever instruções cuidadosas para um LLM — é uma
    forma de conhecimento artesanal? Justifique em até 5 linhas.
  \end{enumerate}

  \sol{
    \textbf{Respostas:}
    (a) A interface \texttt{obs → decidir(obs) → ação} funciona para ambos porque ela
    abstrai \textbf{o que} o agente faz de \textbf{como} ele decide. O paradigma de
    agentes é uma \textit{casca}: define a assinatura da decisão, não sua implementação.
    O que muda entre o aspirador e o GPT é apenas a função \texttt{decidir()} — e essa
    função pode ser desde 4 ifs até uma rede de 175 bilhões de parâmetros. Isso revela
    que o paradigma de agentes é \textbf{universal}: ele unifica toda a IA sob uma
    única abstração.

    (b) Sim, prompt engineering é conhecimento artesanal — mas de um tipo novo. É
    artesanal porque um humano está manualmente ajustando a entrada para melhorar
    a saída, sem que o modelo aprenda com esse ajuste. Mas difere do artesanal
    clássico porque (1) o modelo subjacente foi aprendido de dados, não programado,
    e (2) o custo de ajustar um prompt é ordens de grandeza menor que o de escrever
    regras. Dito isso, a Bitter Lesson sugere que, no longo prazo, métodos que
    \textit{aprendem} a fazer prompting (auto-prompting, RLHF) substituirão o
    prompting manual — e já há evidências disso.

    \errocomum{Em (b), responder simplesmente ``não'' porque o modelo foi treinado com ML.
    O prompt em si é conhecimento artesanal injetado \textit{na inferência}, não no treino.}

    \competencia{Reconhecer o loop percepção→decisão→ação como a interface universal
    da IA e aplicar a Bitter Lesson a desenvolvimentos recentes.}
  }
\end{questao}

\vspace{1cm}
\hrule
\vspace{0.3cm}
\begin{center}
  {\footnotesize Lista S01 · GBC063 · Prof. Marcelo Keese Albertini · FACOM/UFU}
  \ifdefined\GABARITO
    \ifnum\GABARITO=1
      \\[0.2cm] {\footnotesize\color{red}\textbf{GABARITO — USO EXCLUSIVO DO PROFESSOR}}
    \fi
  \fi
\end{center}

\end{document}
"
% ═══════════════════════════════════════════════════════════════

\documentclass[12pt,a4paper]{article}

\usepackage[utf8]{inputenc}
\usepackage[T1]{fontenc}
\usepackage[brazil]{babel}
\usepackage{geometry}
\geometry{margin=2.2cm}
\usepackage{enumitem}
\usepackage{amsmath,amssymb}
\usepackage{xcolor}
\usepackage{hyperref}
% Using standard fonts available on any TeX distribution
\usepackage{mathpazo}
\usepackage[scaled=0.9]{helvet}
\usepackage{courier}
\renewcommand{\familydefault}{\rmdefault}

% ── Conditional: gabarito on/off ──
\ifdefined\GABARITO
  \ifnum\GABARITO=1
    \newcommand{\sol}[1]{\textcolor{blue}{#1}}
    \newcommand{\errocomum}[1]{\textcolor{orange}{⚠️ Erro comum: #1}}
    \newcommand{\competencia}[1]{\textcolor{gray}{Competência: #1}}
  \else
    \newcommand{\sol}[1]{}
    \newcommand{\errocomum}[1]{}
    \newcommand{\competencia}[1]{}
  \fi
\else
  \newcommand{\sol}[1]{}
  \newcommand{\errocomum}[1]{}
  \newcommand{\competencia}[1]{}
\fi

% ── Question formatting ──
\newcommand{\nivel}[1]{\textcolor{gray}{[#1]}}
\newenvironment{questao}[2]{
  \vspace{0.5cm}
  \noindent\textbf{Q#1} \nivel{#2} \quad
}{
  \vspace{0.2cm}
}

\begin{document}

% ════════════════ HEADER ════════════════
\begin{center}
  {\sffamily\Large GBC063 · Semana 1 · Inteligência Artificial}

  \vspace{0.2cm}
  {\small Prof. Marcelo Keese Albertini · FACOM · UFU}

  \vspace{0.2cm}
  {\footnotesize Cobre as competências do resumo S01 ·
  \href{../semana01/material_introducao_ia.html}{material de estudo} ·
  \href{../semana01/resumo_s01_introducao_ia.md}{resumo}}

  \ifdefined\GABARITO
    \ifnum\GABARITO=1
      \vspace{0.3cm}
      \fbox{\textcolor{red}{\textbf{GABARITO — Documento Exclusivo do Professor}}}
    \fi
  \fi
\end{center}

\vspace{0.3cm}
\hrule
\vspace{0.5cm}

% ════════════════ NÍVEL ● (REATIVAÇÃO) ════════════════
\section*{Nível ● \quad Reativação}

\begin{questao}{1}{●}
  Defina \textbf{agente racional} segundo Russell \& Norvig e explique,
  em no máximo 3 linhas, por que esta definição — e não o Teste de Turing —
  é adotada como definição operacional do curso.

  \sol{
    \textbf{Resposta:} Um agente racional é aquele que, para cada percepção possível,
    seleciona a ação que maximiza a métrica de desempenho esperada. O curso adota esta
    definição porque ela (1) funciona com qualquer métrica de sucesso, (2) não depende
    de replicar comportamento humano, e (3) unifica desde um termostato até o AlphaGo.

    \textbf{Resolução:} Russell \& Norvig (2013), Cap. 2. A definição de agente racional
    é independente de humanidade — um agente pode ser racional e não-humano (ex.: AlphaGo
    joga Go de forma sobre-humana).

    \errocomum{Definir agente racional como "agente que age como um humano". Isso é a visão
    "agir humanamente" (Teste de Turing), não "agir racionalmente".}

    \competencia{Explicar a diferença entre IA simbólica e aprendizado de máquina —
    e definir agente racional.}
  }
\end{questao}

\begin{questao}{2}{●}
  As quatro visões clássicas da IA organizam-se em uma matriz 2×2:
  pensamento \textit{vs.} ação × desempenho humano \textit{vs.} racionalidade.
  Classifique cada um dos sistemas abaixo em \textbf{um} dos quatro quadrantes
  e justifique em 1 linha:

  \begin{enumerate}[label=(\alph*),nosep]
    \item ELIZA (Weizenbaum, 1966)
    \item Deep Blue (IBM, 1997)
    \item O agente aspirador de pó da Aula B
    \item Um modelo de linguagem treinado para prever o próximo token
  \end{enumerate}

  \sol{
    \textbf{Respostas:}
    (a) Agir humanamente — ELIZA simulava um terapeuta humano sem raciocinar.
    (b) Agir racionalmente — Deep Blue selecionava lances que maximizavam chance de vitória,
    sem imitar o pensamento de um grande mestre humano.
    (c) Agir racionalmente — o aspirador maximiza a limpeza dada sua percepção.
    (d) Pensar racionalmente (durante o treino: otimiza uma função de perda) e
    agir humanamente (na inferência: gera texto indistinguível do humano).

    \errocomum{Classificar ELIZA como "pensar humanamente". ELIZA não modelava pensamento —
    apenas casava padrões de string.}

    \competencia{Classificar qualquer sistema de IA nas quatro visões e justificar.}
  }
\end{questao}

\begin{questao}{3}{●}
  Verdadeiro ou falso. \textbf{Justifique cada resposta em até 2 linhas.}

  \begin{enumerate}[label=(\alph*),nosep]
    \item "O primeiro inverno da IA (1969–1980) foi causado exclusivamente por falta
    de poder computacional."
    \item "Redes neurais artificiais são uma ideia dos anos 2010."
  \end{enumerate}

  \sol{
    \textbf{Respostas:}
    (a) \textbf{Falso.} O inverno foi causado principalmente pelo livro \textit{Perceptrons}
    (Minsky \& Papert, 1969), que provou limitações do perceptron simples e levou o campo
    a concluir — erroneamente — que redes neurais eram um beco sem saída. A falta de
    computação contribuiu, mas o fator decisivo foi a perda de credibilidade e financiamento.

    (b) \textbf{Falso.} O primeiro modelo de neurônio artificial é de McCulloch \& Pitts (1943).
    O perceptron é de Rosenblatt (1959). O backpropagation que as treina é de 1986
    (Rumelhart, Hinton \& Williams). As ideias fundamentais têm 40–80 anos; o que mudou
    em 2012 foi a escala (dados + GPUs).

    \errocomum{Para (a): atribuir o inverno só à falta de hardware. O fator principal foi
    sociológico — a comunidade perdeu confiança nas redes neurais.}
    \errocomum{Para (b): achar que deep learning é uma invenção recente. As ideias são antigas;
    o que é novo é a escala.}

    \competencia{Enunciar a Bitter Lesson e dar exemplos históricos.}
  }
\end{questao}

% ════════════════ NÍVEL ●● (EXECUÇÃO) ════════════════
\section*{Nível ●● \quad Execução}

\begin{questao}{4}{●●}
  Considere o seguinte ambiente de aspirador de pó com \textbf{três} cômodos
  em linha (A–B–C). O agente começa na posição B. O cômodo A está sujo; B e C
  estão limpos. O agente segue exatamente a lógica abaixo:

  \begin{verbatim}
  def agente_aspirador(percepcao):
      posicao, sujo = percepcao
      if sujo:        return "ASPIRAR"
      if posicao == "A": return "DIREITA"
      if posicao == "B": return "DIREITA"
      return "ESQUERDA"
  \end{verbatim}

  Preencha a tabela de execução para os \textbf{6 primeiros passos}, indicando:
  posição atual, sensor (sujo/limpo), ação tomada, e estado dos cômodos após a ação.
  O agente limpa todos os cômodos? Em quantos passos?

  \sol{
    \textbf{Resposta:}

    \begin{tabular}{c|c|c|c|c}
      Passo & Posição & Sensor & Ação & Cômodos [A,B,C] \\
      \hline
      0 & B & limpo & DIREITA & [sujo, limpo, limpo] \\
      1 & C & limpo & ESQUERDA & [sujo, limpo, limpo] \\
      2 & B & limpo & DIREITA & [sujo, limpo, limpo] \\
      3 & C & limpo & ESQUERDA & [sujo, limpo, limpo] \\
      \dots & \dots & \dots & \dots & \dots \\
    \end{tabular}

    O agente \textbf{nunca alcança o cômodo A} porque sua lógica só vai de A→B→C→B→C→B→C\ldots
    Ele fica preso no ciclo B↔C. Para alcançar A, seria preciso uma regra adicional
    (ex.: "se estou em C e ambos A e B estão sujos, vá para B").

    \textbf{Resolução:} O bug está na lógica de movimento: quando a posição é B ou C,
    o agente nunca se move para A. A condição \texttt{posicao == "A"} nunca é verdadeira
    se o agente não consegue chegar a A.

    \errocomum{Assumir que o agente "dá a volta". O ambiente é uma linha, não um ciclo —
    de C só se vai para B.}

    \competencia{Implementar o agente aspirador e prever seu comportamento
    com 3 cômodos.}
  }
\end{questao}

\begin{questao}{5}{●●}
  O motor de encadeamento progressivo abaixo recebe uma base de fatos inicial
  e um conjunto de regras. Execute o algoritmo \textbf{passo a passo} com os
  dados fornecidos e indique o conjunto final de fatos.

  \begin{verbatim}
  fatos = {"chovendo", "temperatura < 15"}
  regras = [
    ({"chovendo", "temperatura < 15"}, "risco de gelo"),
    ({"risco de gelo"},              "fechar pista"),
    ({"chovendo"},                   "pista molhada"),
  ]
  \end{verbatim}

  \sol{
    \textbf{Resposta:} fatos finais = \{\texttt{"chovendo"}, \texttt{"temperatura < 15"},
    \texttt{"risco de gelo"}, \texttt{"fechar pista"}, \texttt{"pista molhada"}\}

    \textbf{Resolução:}
    \begin{enumerate}[nosep]
      \item Iteração 1: chovendo + temperatura < 15 → adiciona \texttt{"risco de gelo"}.
      \item Iteração 2: risco de gelo → adiciona \texttt{"fechar pista"}.
      \item Iteração 3: chovendo → adiciona \texttt{"pista molhada"}.
      \item Iteração 4: nenhuma regra nova dispara → fim.
    \end{enumerate}

    \errocomum{Esquecer que as regras disparam em qualquer ordem — o algoritmo testa TODAS
    as regras a cada iteração, não para na primeira que dispara.}

    \competencia{Explicar o funcionamento de um motor de inferência e por que regras
    não escalam.}
  }
\end{questao}

\begin{questao}{6}{●●}
  Em 2023, um LLM foi capaz de passar no exame da ordem dos advogados (Bar Exam)
  dos EUA sem ter recebido uma única regra explícita sobre direito. Usando os
  conceitos da Semana 1, explique \textbf{por que isso é um fato tão significativo}
  na história da IA. Se não obteve uma resposta completa, use o roteiro:
  (1) como um sistema especialista resolveria o Bar Exam? (2) por que essa abordagem
  falharia? (3) o que o LLM fez de diferente?

  \sol{
    \textbf{Resposta:} Um sistema especialista exigiria que engenheiros de conhecimento
    codificassem milhares de regras jurídicas à mão — um gargalo intransponível. O LLM
    não recebeu regra nenhuma: aprendeu padrões estatísticos de centenas de gigabytes
    de texto jurídico e geral. Isso é a Bitter Lesson em ação: conhecimento aprendido
    de dados venceu conhecimento artesanal — em um domínio (direito) que se supunha
    exigir raciocínio simbólico explícito.

    \errocomum{Dizer que o LLM "decorou" as respostas. LLMs não têm memória explícita do
    corpus de treino — eles aprendem padrões estatísticos, não fazem lookup de frases.}

    \competencia{Enunciar a Bitter Lesson e dar exemplos que a corroboram.}
  }
\end{questao}

% ════════════════ NÍVEL ●●● (ANÁLISE) ════════════════
\section*{Nível ●●● \quad Análise}

\begin{questao}{7}{●●●}
  Considere a afirmação de Sutton (2019):

  \begin{quote}
  \textit{``Métodos gerais que alavancam computação são, em última instância,
  os mais efetivos — e por uma grande margem. A maioria do esforço em IA é
  gasto tentando construir conhecimento em agentes; isso ajuda no curto prazo,
  mas no longo prazo atinge um platô.''}
  \end{quote}

  \begin{enumerate}[label=(\alph*),nosep]
    \item Por que conhecimento artesanal \textit{necessariamente} atinge um platô,
    enquanto métodos que aprendem de dados não? (Dica: pense no que acontece quando
    você dobra a quantidade de dados ou de computação em cada caso.)
    \item Dê um exemplo \textbf{concreto} de um sistema que atingiu esse platô e
    explique \textit{o que especificamente} o limitou.
  \end{enumerate}

  \sol{
    \textbf{Respostas:}
    (a) Conhecimento artesanal atinge um platô porque sua melhoria depende de \textbf{humanos
    adicionando mais regras} — um processo que é linear, caro e sujeito a contradições.
    Métodos que aprendem melhoram quando você adiciona mais dados ou computação, e esse
    processo é (em princípio) ilimitado: mais dados → melhores gradientes → melhor modelo.
    Não há um humano no loop ditando o teto.

    (b) Sistemas especialistas médicos (ex.: MYCIN, anos 80). Eram melhores que médicos
    juniores em diagnósticos específicos, mas cada nova doença exigia novas regras escritas
    por especialistas. Quando o número de regras passou de centenas, as interações entre
    elas tornaram-se imprevisíveis. O sistema não melhorava com mais dados de pacientes —
    só com mais regras escritas à mão. Enquanto isso, métodos de ML atuais melhoram
    automaticamente com mais exames, prontuários e artigos.

    \errocomum{Em (a), dizer que conhecimento artesanal atinge platô "porque humanos
    são limitados". A resposta precisa identificar o mecanismo: a taxa de melhoria é
    ditada pela velocidade com que humanos escrevem regras, não pela velocidade com
    que computação cresce.}

    \competencia{Enunciar a Bitter Lesson e explicar o mecanismo por trás dela.}
  }
\end{questao}

\begin{questao}{8}{●●●}
  O efeito ELIZA — atribuir inteligência onde há apenas casamento de padrões —
  foi documentado em 1966. Ele persiste em 2025 com LLMs?

  \begin{enumerate}[label=(\alph*),nosep]
    \item Aponte \textbf{uma semelhança} entre ELIZA e um LLM moderno que torna
    o efeito ELIZA relevante hoje.
    \item Aponte \textbf{uma diferença fundamental} que torna a comparação enganosa.
    \item Em no máximo 4 linhas: um LLM é ``só casamento de padrões''? Justifique.
  \end{enumerate}

  \sol{
    \textbf{Respostas:}
    (a) Ambos produzem texto fluente sem garantia de compreensão do conteúdo semântico.
    Um LLM pode gerar uma explicação convincente de um conceito que ele não ``entende''
    no sentido humano — assim como ELIZA simulava empatia sem sentir nada. O efeito
    ELIZA ocorre quando o usuário projeta compreensão onde há apenas correlação estatística.

    (b) ELIZA usava algumas dezenas de regras de reescrita escritas à mão. Um LLM aprende
    de centenas de bilhões de parâmetros ajustados sobre trilhões de tokens. A escala e
    a origem do conhecimento (artesanal vs. aprendido) são de natureza completamente diferente.

    (c) Não. Reduzir um LLM a ``só casamento de padrões'' é como reduzir um cérebro a
    ``só sinapses disparando'' — tecnicamente verdadeiro, mas enganoso. A questão relevante
    não é se o mecanismo básico é simples (é: multiplicação de matrizes + não-linearidades),
    mas se a escala produz comportamentos emergentes que merecem ser chamados de raciocínio.
    Esta é uma pergunta em aberto na pesquisa em 2025.

    \errocomum{Responder (c) com um simples ``sim'' ou ``não''. A pergunta exige distinguir
    mecanismo de comportamento — o que o sistema \textit{faz} vs. como ele \textit{faz}.}

    \competencia{Classificar sistemas de IA e analisar criticamente alegações de inteligência.}
  }
\end{questao}

% ════════════════ NÍVEL ★ (SÍNTESE) ════════════════
\section*{Nível ★ \quad Síntese}

\begin{questao}{9}{★}
  Projete um experimento mental: você herda um sistema especialista médico com
  10.000 regras, desenvolvido ao longo de 15 anos por uma equipe de 20 engenheiros
  de conhecimento. O sistema tem 92\% de acurácia em diagnóstico.

  \begin{enumerate}[label=(\alph*),nosep]
    \item Identifique \textbf{três} cenários em que este sistema falharia de forma
    sistemática que um modelo de ML treinado sobre os mesmos dados não falharia.
    \item Para cada cenário, explique \textbf{por que} a natureza artesanal das
    regras é a causa raiz da falha.
    \item Proponha uma arquitetura híbrida (simbólico + aprendizado) que combine
    o melhor dos dois mundos para este domínio. Descreva em no máximo 6 linhas.
  \end{enumerate}

  \sol{
    \textbf{Respostas:}
    (a) e (b):
    \begin{enumerate}[nosep]
      \item \textbf{Nova doença:} o SE falha porque nenhuma regra cobre o fenômeno novo.
      O ML pode detectar que os sintomas formam um cluster distinto nos dados e sinalizar
      uma anomalia — mesmo sem rótulos. Causa raiz: SE só reconhece o que foi explicitamente
      programado.
      \item \textbf{Interação entre condições:} um paciente com diabetes \textit{e} COVID
      pode ativar regras contraditórias (uma recomenda corticoides, outra contraindica).
      O SE não tem como ponderar evidências conflitantes. O ML aprende pesos relativos
      dos features. Causa raiz: regras são discretas e não-compostas; não modelam interações.
      \item \textbf{Deriva de distribuição:} se a população de pacientes mudar (ex.:
      envelhecimento), a acurácia do SE cai porque as regras foram escritas para a
      distribuição antiga. O ML pode ser retreinado com dados novos. Causa raiz: SE
      não se adapta sem reengenharia manual.
    \end{enumerate}

    (c) \textbf{Arquitetura híbrida:} usar um modelo de ML (ex.: gradient-boosted trees
    ou rede neural) para produzir o diagnóstico base e a estimativa de confiança. As
    10.000 regras viram um sistema de \textbf{verificação}: antes de emitir o diagnóstico,
    o sistema checa se ele viola alguma regra de segurança crítica (ex.: ``nunca prescrever
    droga X para grávidas''). Se violar, escala para um médico humano. O ML lida com a
    generalização; as regras lidam com as restrições rígidas de segurança.

    \errocomum{Propor um híbrido em que as regras \textit{guiam} o aprendizado. Isso
    reintroduz o gargalo: as regras limitam o que o modelo pode aprender. O melhor
    híbrido usa regras como \textit{restrições de segurança}, não como indutores de viés.}

    \competencia{Sintetizar os conceitos de IA simbólica e aprendizado em uma arquitetura
    fundamentada.}
  }
\end{questao}

\begin{questao}{10}{★ 🌉}
  O treinamento do GPT usa a mesma estrutura de loop que o agente aspirador da Aula B:
  \textbf{percepção → decisão → ação}. A diferença está na complexidade de \texttt{decidir()}.

  \begin{enumerate}[label=(\alph*),nosep]
    \item No aspirador, \texttt{decidir()} é uma função com 4 regras \texttt{if}.
    No GPT, \texttt{decidir()} é uma rede neural com centenas de bilhões de
    parâmetros. \textbf{Por que}, conceitualmente, a mesma interface de 4 linhas
    funciona para ambos? O que isso revela sobre o paradigma de agentes?
    \item A Bitter Lesson argumenta que conhecimento artesanal atinge um platô.
    O prompt engineering — escrever instruções cuidadosas para um LLM — é uma
    forma de conhecimento artesanal? Justifique em até 5 linhas.
  \end{enumerate}

  \sol{
    \textbf{Respostas:}
    (a) A interface \texttt{obs → decidir(obs) → ação} funciona para ambos porque ela
    abstrai \textbf{o que} o agente faz de \textbf{como} ele decide. O paradigma de
    agentes é uma \textit{casca}: define a assinatura da decisão, não sua implementação.
    O que muda entre o aspirador e o GPT é apenas a função \texttt{decidir()} — e essa
    função pode ser desde 4 ifs até uma rede de 175 bilhões de parâmetros. Isso revela
    que o paradigma de agentes é \textbf{universal}: ele unifica toda a IA sob uma
    única abstração.

    (b) Sim, prompt engineering é conhecimento artesanal — mas de um tipo novo. É
    artesanal porque um humano está manualmente ajustando a entrada para melhorar
    a saída, sem que o modelo aprenda com esse ajuste. Mas difere do artesanal
    clássico porque (1) o modelo subjacente foi aprendido de dados, não programado,
    e (2) o custo de ajustar um prompt é ordens de grandeza menor que o de escrever
    regras. Dito isso, a Bitter Lesson sugere que, no longo prazo, métodos que
    \textit{aprendem} a fazer prompting (auto-prompting, RLHF) substituirão o
    prompting manual — e já há evidências disso.

    \errocomum{Em (b), responder simplesmente ``não'' porque o modelo foi treinado com ML.
    O prompt em si é conhecimento artesanal injetado \textit{na inferência}, não no treino.}

    \competencia{Reconhecer o loop percepção→decisão→ação como a interface universal
    da IA e aplicar a Bitter Lesson a desenvolvimentos recentes.}
  }
\end{questao}

\vspace{1cm}
\hrule
\vspace{0.3cm}
\begin{center}
  {\footnotesize Lista S01 · GBC063 · Prof. Marcelo Keese Albertini · FACOM/UFU}
  \ifdefined\GABARITO
    \ifnum\GABARITO=1
      \\[0.2cm] {\footnotesize\color{red}\textbf{GABARITO — USO EXCLUSIVO DO PROFESSOR}}
    \fi
  \fi
\end{center}

\end{document}
"
%   Gabarito:  pdflatex -jobname=gabarito_lista_s01_introducao_ia "\def\GABARITO{1}% ═══════════════════════════════════════════════════════════════
% Lista S01 — O que é Inteligência Artificial?
% Prof. Marcelo Keese Albertini · FACOM · UFU · GBC063
%
% Compilação:
%   Questões:  pdflatex -jobname=lista_s01_introducao_ia "\def\GABARITO{0}% ═══════════════════════════════════════════════════════════════
% Lista S01 — O que é Inteligência Artificial?
% Prof. Marcelo Keese Albertini · FACOM · UFU · GBC063
%
% Compilação:
%   Questões:  pdflatex -jobname=lista_s01_introducao_ia "\def\GABARITO{0}% ═══════════════════════════════════════════════════════════════
% Lista S01 — O que é Inteligência Artificial?
% Prof. Marcelo Keese Albertini · FACOM · UFU · GBC063
%
% Compilação:
%   Questões:  pdflatex -jobname=lista_s01_introducao_ia "\def\GABARITO{0}\input{lista_s01_introducao_ia.tex}"
%   Gabarito:  pdflatex -jobname=gabarito_lista_s01_introducao_ia "\def\GABARITO{1}\input{lista_s01_introducao_ia.tex}"
% ═══════════════════════════════════════════════════════════════

\documentclass[12pt,a4paper]{article}

\usepackage[utf8]{inputenc}
\usepackage[T1]{fontenc}
\usepackage[brazil]{babel}
\usepackage{geometry}
\geometry{margin=2.2cm}
\usepackage{enumitem}
\usepackage{amsmath,amssymb}
\usepackage{xcolor}
\usepackage{hyperref}
% Using standard fonts available on any TeX distribution
\usepackage{mathpazo}
\usepackage[scaled=0.9]{helvet}
\usepackage{courier}
\renewcommand{\familydefault}{\rmdefault}

% ── Conditional: gabarito on/off ──
\ifdefined\GABARITO
  \ifnum\GABARITO=1
    \newcommand{\sol}[1]{\textcolor{blue}{#1}}
    \newcommand{\errocomum}[1]{\textcolor{orange}{⚠️ Erro comum: #1}}
    \newcommand{\competencia}[1]{\textcolor{gray}{Competência: #1}}
  \else
    \newcommand{\sol}[1]{}
    \newcommand{\errocomum}[1]{}
    \newcommand{\competencia}[1]{}
  \fi
\else
  \newcommand{\sol}[1]{}
  \newcommand{\errocomum}[1]{}
  \newcommand{\competencia}[1]{}
\fi

% ── Question formatting ──
\newcommand{\nivel}[1]{\textcolor{gray}{[#1]}}
\newenvironment{questao}[2]{
  \vspace{0.5cm}
  \noindent\textbf{Q#1} \nivel{#2} \quad
}{
  \vspace{0.2cm}
}

\begin{document}

% ════════════════ HEADER ════════════════
\begin{center}
  {\sffamily\Large GBC063 · Semana 1 · Inteligência Artificial}

  \vspace{0.2cm}
  {\small Prof. Marcelo Keese Albertini · FACOM · UFU}

  \vspace{0.2cm}
  {\footnotesize Cobre as competências do resumo S01 ·
  \href{../semana01/material_introducao_ia.html}{material de estudo} ·
  \href{../semana01/resumo_s01_introducao_ia.md}{resumo}}

  \ifdefined\GABARITO
    \ifnum\GABARITO=1
      \vspace{0.3cm}
      \fbox{\textcolor{red}{\textbf{GABARITO — Documento Exclusivo do Professor}}}
    \fi
  \fi
\end{center}

\vspace{0.3cm}
\hrule
\vspace{0.5cm}

% ════════════════ NÍVEL ● (REATIVAÇÃO) ════════════════
\section*{Nível ● \quad Reativação}

\begin{questao}{1}{●}
  Defina \textbf{agente racional} segundo Russell \& Norvig e explique,
  em no máximo 3 linhas, por que esta definição — e não o Teste de Turing —
  é adotada como definição operacional do curso.

  \sol{
    \textbf{Resposta:} Um agente racional é aquele que, para cada percepção possível,
    seleciona a ação que maximiza a métrica de desempenho esperada. O curso adota esta
    definição porque ela (1) funciona com qualquer métrica de sucesso, (2) não depende
    de replicar comportamento humano, e (3) unifica desde um termostato até o AlphaGo.

    \textbf{Resolução:} Russell \& Norvig (2013), Cap. 2. A definição de agente racional
    é independente de humanidade — um agente pode ser racional e não-humano (ex.: AlphaGo
    joga Go de forma sobre-humana).

    \errocomum{Definir agente racional como "agente que age como um humano". Isso é a visão
    "agir humanamente" (Teste de Turing), não "agir racionalmente".}

    \competencia{Explicar a diferença entre IA simbólica e aprendizado de máquina —
    e definir agente racional.}
  }
\end{questao}

\begin{questao}{2}{●}
  As quatro visões clássicas da IA organizam-se em uma matriz 2×2:
  pensamento \textit{vs.} ação × desempenho humano \textit{vs.} racionalidade.
  Classifique cada um dos sistemas abaixo em \textbf{um} dos quatro quadrantes
  e justifique em 1 linha:

  \begin{enumerate}[label=(\alph*),nosep]
    \item ELIZA (Weizenbaum, 1966)
    \item Deep Blue (IBM, 1997)
    \item O agente aspirador de pó da Aula B
    \item Um modelo de linguagem treinado para prever o próximo token
  \end{enumerate}

  \sol{
    \textbf{Respostas:}
    (a) Agir humanamente — ELIZA simulava um terapeuta humano sem raciocinar.
    (b) Agir racionalmente — Deep Blue selecionava lances que maximizavam chance de vitória,
    sem imitar o pensamento de um grande mestre humano.
    (c) Agir racionalmente — o aspirador maximiza a limpeza dada sua percepção.
    (d) Pensar racionalmente (durante o treino: otimiza uma função de perda) e
    agir humanamente (na inferência: gera texto indistinguível do humano).

    \errocomum{Classificar ELIZA como "pensar humanamente". ELIZA não modelava pensamento —
    apenas casava padrões de string.}

    \competencia{Classificar qualquer sistema de IA nas quatro visões e justificar.}
  }
\end{questao}

\begin{questao}{3}{●}
  Verdadeiro ou falso. \textbf{Justifique cada resposta em até 2 linhas.}

  \begin{enumerate}[label=(\alph*),nosep]
    \item "O primeiro inverno da IA (1969–1980) foi causado exclusivamente por falta
    de poder computacional."
    \item "Redes neurais artificiais são uma ideia dos anos 2010."
  \end{enumerate}

  \sol{
    \textbf{Respostas:}
    (a) \textbf{Falso.} O inverno foi causado principalmente pelo livro \textit{Perceptrons}
    (Minsky \& Papert, 1969), que provou limitações do perceptron simples e levou o campo
    a concluir — erroneamente — que redes neurais eram um beco sem saída. A falta de
    computação contribuiu, mas o fator decisivo foi a perda de credibilidade e financiamento.

    (b) \textbf{Falso.} O primeiro modelo de neurônio artificial é de McCulloch \& Pitts (1943).
    O perceptron é de Rosenblatt (1959). O backpropagation que as treina é de 1986
    (Rumelhart, Hinton \& Williams). As ideias fundamentais têm 40–80 anos; o que mudou
    em 2012 foi a escala (dados + GPUs).

    \errocomum{Para (a): atribuir o inverno só à falta de hardware. O fator principal foi
    sociológico — a comunidade perdeu confiança nas redes neurais.}
    \errocomum{Para (b): achar que deep learning é uma invenção recente. As ideias são antigas;
    o que é novo é a escala.}

    \competencia{Enunciar a Bitter Lesson e dar exemplos históricos.}
  }
\end{questao}

% ════════════════ NÍVEL ●● (EXECUÇÃO) ════════════════
\section*{Nível ●● \quad Execução}

\begin{questao}{4}{●●}
  Considere o seguinte ambiente de aspirador de pó com \textbf{três} cômodos
  em linha (A–B–C). O agente começa na posição B. O cômodo A está sujo; B e C
  estão limpos. O agente segue exatamente a lógica abaixo:

  \begin{verbatim}
  def agente_aspirador(percepcao):
      posicao, sujo = percepcao
      if sujo:        return "ASPIRAR"
      if posicao == "A": return "DIREITA"
      if posicao == "B": return "DIREITA"
      return "ESQUERDA"
  \end{verbatim}

  Preencha a tabela de execução para os \textbf{6 primeiros passos}, indicando:
  posição atual, sensor (sujo/limpo), ação tomada, e estado dos cômodos após a ação.
  O agente limpa todos os cômodos? Em quantos passos?

  \sol{
    \textbf{Resposta:}

    \begin{tabular}{c|c|c|c|c}
      Passo & Posição & Sensor & Ação & Cômodos [A,B,C] \\
      \hline
      0 & B & limpo & DIREITA & [sujo, limpo, limpo] \\
      1 & C & limpo & ESQUERDA & [sujo, limpo, limpo] \\
      2 & B & limpo & DIREITA & [sujo, limpo, limpo] \\
      3 & C & limpo & ESQUERDA & [sujo, limpo, limpo] \\
      \dots & \dots & \dots & \dots & \dots \\
    \end{tabular}

    O agente \textbf{nunca alcança o cômodo A} porque sua lógica só vai de A→B→C→B→C→B→C\ldots
    Ele fica preso no ciclo B↔C. Para alcançar A, seria preciso uma regra adicional
    (ex.: "se estou em C e ambos A e B estão sujos, vá para B").

    \textbf{Resolução:} O bug está na lógica de movimento: quando a posição é B ou C,
    o agente nunca se move para A. A condição \texttt{posicao == "A"} nunca é verdadeira
    se o agente não consegue chegar a A.

    \errocomum{Assumir que o agente "dá a volta". O ambiente é uma linha, não um ciclo —
    de C só se vai para B.}

    \competencia{Implementar o agente aspirador e prever seu comportamento
    com 3 cômodos.}
  }
\end{questao}

\begin{questao}{5}{●●}
  O motor de encadeamento progressivo abaixo recebe uma base de fatos inicial
  e um conjunto de regras. Execute o algoritmo \textbf{passo a passo} com os
  dados fornecidos e indique o conjunto final de fatos.

  \begin{verbatim}
  fatos = {"chovendo", "temperatura < 15"}
  regras = [
    ({"chovendo", "temperatura < 15"}, "risco de gelo"),
    ({"risco de gelo"},              "fechar pista"),
    ({"chovendo"},                   "pista molhada"),
  ]
  \end{verbatim}

  \sol{
    \textbf{Resposta:} fatos finais = \{\texttt{"chovendo"}, \texttt{"temperatura < 15"},
    \texttt{"risco de gelo"}, \texttt{"fechar pista"}, \texttt{"pista molhada"}\}

    \textbf{Resolução:}
    \begin{enumerate}[nosep]
      \item Iteração 1: chovendo + temperatura < 15 → adiciona \texttt{"risco de gelo"}.
      \item Iteração 2: risco de gelo → adiciona \texttt{"fechar pista"}.
      \item Iteração 3: chovendo → adiciona \texttt{"pista molhada"}.
      \item Iteração 4: nenhuma regra nova dispara → fim.
    \end{enumerate}

    \errocomum{Esquecer que as regras disparam em qualquer ordem — o algoritmo testa TODAS
    as regras a cada iteração, não para na primeira que dispara.}

    \competencia{Explicar o funcionamento de um motor de inferência e por que regras
    não escalam.}
  }
\end{questao}

\begin{questao}{6}{●●}
  Em 2023, um LLM foi capaz de passar no exame da ordem dos advogados (Bar Exam)
  dos EUA sem ter recebido uma única regra explícita sobre direito. Usando os
  conceitos da Semana 1, explique \textbf{por que isso é um fato tão significativo}
  na história da IA. Se não obteve uma resposta completa, use o roteiro:
  (1) como um sistema especialista resolveria o Bar Exam? (2) por que essa abordagem
  falharia? (3) o que o LLM fez de diferente?

  \sol{
    \textbf{Resposta:} Um sistema especialista exigiria que engenheiros de conhecimento
    codificassem milhares de regras jurídicas à mão — um gargalo intransponível. O LLM
    não recebeu regra nenhuma: aprendeu padrões estatísticos de centenas de gigabytes
    de texto jurídico e geral. Isso é a Bitter Lesson em ação: conhecimento aprendido
    de dados venceu conhecimento artesanal — em um domínio (direito) que se supunha
    exigir raciocínio simbólico explícito.

    \errocomum{Dizer que o LLM "decorou" as respostas. LLMs não têm memória explícita do
    corpus de treino — eles aprendem padrões estatísticos, não fazem lookup de frases.}

    \competencia{Enunciar a Bitter Lesson e dar exemplos que a corroboram.}
  }
\end{questao}

% ════════════════ NÍVEL ●●● (ANÁLISE) ════════════════
\section*{Nível ●●● \quad Análise}

\begin{questao}{7}{●●●}
  Considere a afirmação de Sutton (2019):

  \begin{quote}
  \textit{``Métodos gerais que alavancam computação são, em última instância,
  os mais efetivos — e por uma grande margem. A maioria do esforço em IA é
  gasto tentando construir conhecimento em agentes; isso ajuda no curto prazo,
  mas no longo prazo atinge um platô.''}
  \end{quote}

  \begin{enumerate}[label=(\alph*),nosep]
    \item Por que conhecimento artesanal \textit{necessariamente} atinge um platô,
    enquanto métodos que aprendem de dados não? (Dica: pense no que acontece quando
    você dobra a quantidade de dados ou de computação em cada caso.)
    \item Dê um exemplo \textbf{concreto} de um sistema que atingiu esse platô e
    explique \textit{o que especificamente} o limitou.
  \end{enumerate}

  \sol{
    \textbf{Respostas:}
    (a) Conhecimento artesanal atinge um platô porque sua melhoria depende de \textbf{humanos
    adicionando mais regras} — um processo que é linear, caro e sujeito a contradições.
    Métodos que aprendem melhoram quando você adiciona mais dados ou computação, e esse
    processo é (em princípio) ilimitado: mais dados → melhores gradientes → melhor modelo.
    Não há um humano no loop ditando o teto.

    (b) Sistemas especialistas médicos (ex.: MYCIN, anos 80). Eram melhores que médicos
    juniores em diagnósticos específicos, mas cada nova doença exigia novas regras escritas
    por especialistas. Quando o número de regras passou de centenas, as interações entre
    elas tornaram-se imprevisíveis. O sistema não melhorava com mais dados de pacientes —
    só com mais regras escritas à mão. Enquanto isso, métodos de ML atuais melhoram
    automaticamente com mais exames, prontuários e artigos.

    \errocomum{Em (a), dizer que conhecimento artesanal atinge platô "porque humanos
    são limitados". A resposta precisa identificar o mecanismo: a taxa de melhoria é
    ditada pela velocidade com que humanos escrevem regras, não pela velocidade com
    que computação cresce.}

    \competencia{Enunciar a Bitter Lesson e explicar o mecanismo por trás dela.}
  }
\end{questao}

\begin{questao}{8}{●●●}
  O efeito ELIZA — atribuir inteligência onde há apenas casamento de padrões —
  foi documentado em 1966. Ele persiste em 2025 com LLMs?

  \begin{enumerate}[label=(\alph*),nosep]
    \item Aponte \textbf{uma semelhança} entre ELIZA e um LLM moderno que torna
    o efeito ELIZA relevante hoje.
    \item Aponte \textbf{uma diferença fundamental} que torna a comparação enganosa.
    \item Em no máximo 4 linhas: um LLM é ``só casamento de padrões''? Justifique.
  \end{enumerate}

  \sol{
    \textbf{Respostas:}
    (a) Ambos produzem texto fluente sem garantia de compreensão do conteúdo semântico.
    Um LLM pode gerar uma explicação convincente de um conceito que ele não ``entende''
    no sentido humano — assim como ELIZA simulava empatia sem sentir nada. O efeito
    ELIZA ocorre quando o usuário projeta compreensão onde há apenas correlação estatística.

    (b) ELIZA usava algumas dezenas de regras de reescrita escritas à mão. Um LLM aprende
    de centenas de bilhões de parâmetros ajustados sobre trilhões de tokens. A escala e
    a origem do conhecimento (artesanal vs. aprendido) são de natureza completamente diferente.

    (c) Não. Reduzir um LLM a ``só casamento de padrões'' é como reduzir um cérebro a
    ``só sinapses disparando'' — tecnicamente verdadeiro, mas enganoso. A questão relevante
    não é se o mecanismo básico é simples (é: multiplicação de matrizes + não-linearidades),
    mas se a escala produz comportamentos emergentes que merecem ser chamados de raciocínio.
    Esta é uma pergunta em aberto na pesquisa em 2025.

    \errocomum{Responder (c) com um simples ``sim'' ou ``não''. A pergunta exige distinguir
    mecanismo de comportamento — o que o sistema \textit{faz} vs. como ele \textit{faz}.}

    \competencia{Classificar sistemas de IA e analisar criticamente alegações de inteligência.}
  }
\end{questao}

% ════════════════ NÍVEL ★ (SÍNTESE) ════════════════
\section*{Nível ★ \quad Síntese}

\begin{questao}{9}{★}
  Projete um experimento mental: você herda um sistema especialista médico com
  10.000 regras, desenvolvido ao longo de 15 anos por uma equipe de 20 engenheiros
  de conhecimento. O sistema tem 92\% de acurácia em diagnóstico.

  \begin{enumerate}[label=(\alph*),nosep]
    \item Identifique \textbf{três} cenários em que este sistema falharia de forma
    sistemática que um modelo de ML treinado sobre os mesmos dados não falharia.
    \item Para cada cenário, explique \textbf{por que} a natureza artesanal das
    regras é a causa raiz da falha.
    \item Proponha uma arquitetura híbrida (simbólico + aprendizado) que combine
    o melhor dos dois mundos para este domínio. Descreva em no máximo 6 linhas.
  \end{enumerate}

  \sol{
    \textbf{Respostas:}
    (a) e (b):
    \begin{enumerate}[nosep]
      \item \textbf{Nova doença:} o SE falha porque nenhuma regra cobre o fenômeno novo.
      O ML pode detectar que os sintomas formam um cluster distinto nos dados e sinalizar
      uma anomalia — mesmo sem rótulos. Causa raiz: SE só reconhece o que foi explicitamente
      programado.
      \item \textbf{Interação entre condições:} um paciente com diabetes \textit{e} COVID
      pode ativar regras contraditórias (uma recomenda corticoides, outra contraindica).
      O SE não tem como ponderar evidências conflitantes. O ML aprende pesos relativos
      dos features. Causa raiz: regras são discretas e não-compostas; não modelam interações.
      \item \textbf{Deriva de distribuição:} se a população de pacientes mudar (ex.:
      envelhecimento), a acurácia do SE cai porque as regras foram escritas para a
      distribuição antiga. O ML pode ser retreinado com dados novos. Causa raiz: SE
      não se adapta sem reengenharia manual.
    \end{enumerate}

    (c) \textbf{Arquitetura híbrida:} usar um modelo de ML (ex.: gradient-boosted trees
    ou rede neural) para produzir o diagnóstico base e a estimativa de confiança. As
    10.000 regras viram um sistema de \textbf{verificação}: antes de emitir o diagnóstico,
    o sistema checa se ele viola alguma regra de segurança crítica (ex.: ``nunca prescrever
    droga X para grávidas''). Se violar, escala para um médico humano. O ML lida com a
    generalização; as regras lidam com as restrições rígidas de segurança.

    \errocomum{Propor um híbrido em que as regras \textit{guiam} o aprendizado. Isso
    reintroduz o gargalo: as regras limitam o que o modelo pode aprender. O melhor
    híbrido usa regras como \textit{restrições de segurança}, não como indutores de viés.}

    \competencia{Sintetizar os conceitos de IA simbólica e aprendizado em uma arquitetura
    fundamentada.}
  }
\end{questao}

\begin{questao}{10}{★ 🌉}
  O treinamento do GPT usa a mesma estrutura de loop que o agente aspirador da Aula B:
  \textbf{percepção → decisão → ação}. A diferença está na complexidade de \texttt{decidir()}.

  \begin{enumerate}[label=(\alph*),nosep]
    \item No aspirador, \texttt{decidir()} é uma função com 4 regras \texttt{if}.
    No GPT, \texttt{decidir()} é uma rede neural com centenas de bilhões de
    parâmetros. \textbf{Por que}, conceitualmente, a mesma interface de 4 linhas
    funciona para ambos? O que isso revela sobre o paradigma de agentes?
    \item A Bitter Lesson argumenta que conhecimento artesanal atinge um platô.
    O prompt engineering — escrever instruções cuidadosas para um LLM — é uma
    forma de conhecimento artesanal? Justifique em até 5 linhas.
  \end{enumerate}

  \sol{
    \textbf{Respostas:}
    (a) A interface \texttt{obs → decidir(obs) → ação} funciona para ambos porque ela
    abstrai \textbf{o que} o agente faz de \textbf{como} ele decide. O paradigma de
    agentes é uma \textit{casca}: define a assinatura da decisão, não sua implementação.
    O que muda entre o aspirador e o GPT é apenas a função \texttt{decidir()} — e essa
    função pode ser desde 4 ifs até uma rede de 175 bilhões de parâmetros. Isso revela
    que o paradigma de agentes é \textbf{universal}: ele unifica toda a IA sob uma
    única abstração.

    (b) Sim, prompt engineering é conhecimento artesanal — mas de um tipo novo. É
    artesanal porque um humano está manualmente ajustando a entrada para melhorar
    a saída, sem que o modelo aprenda com esse ajuste. Mas difere do artesanal
    clássico porque (1) o modelo subjacente foi aprendido de dados, não programado,
    e (2) o custo de ajustar um prompt é ordens de grandeza menor que o de escrever
    regras. Dito isso, a Bitter Lesson sugere que, no longo prazo, métodos que
    \textit{aprendem} a fazer prompting (auto-prompting, RLHF) substituirão o
    prompting manual — e já há evidências disso.

    \errocomum{Em (b), responder simplesmente ``não'' porque o modelo foi treinado com ML.
    O prompt em si é conhecimento artesanal injetado \textit{na inferência}, não no treino.}

    \competencia{Reconhecer o loop percepção→decisão→ação como a interface universal
    da IA e aplicar a Bitter Lesson a desenvolvimentos recentes.}
  }
\end{questao}

\vspace{1cm}
\hrule
\vspace{0.3cm}
\begin{center}
  {\footnotesize Lista S01 · GBC063 · Prof. Marcelo Keese Albertini · FACOM/UFU}
  \ifdefined\GABARITO
    \ifnum\GABARITO=1
      \\[0.2cm] {\footnotesize\color{red}\textbf{GABARITO — USO EXCLUSIVO DO PROFESSOR}}
    \fi
  \fi
\end{center}

\end{document}
"
%   Gabarito:  pdflatex -jobname=gabarito_lista_s01_introducao_ia "\def\GABARITO{1}% ═══════════════════════════════════════════════════════════════
% Lista S01 — O que é Inteligência Artificial?
% Prof. Marcelo Keese Albertini · FACOM · UFU · GBC063
%
% Compilação:
%   Questões:  pdflatex -jobname=lista_s01_introducao_ia "\def\GABARITO{0}\input{lista_s01_introducao_ia.tex}"
%   Gabarito:  pdflatex -jobname=gabarito_lista_s01_introducao_ia "\def\GABARITO{1}\input{lista_s01_introducao_ia.tex}"
% ═══════════════════════════════════════════════════════════════

\documentclass[12pt,a4paper]{article}

\usepackage[utf8]{inputenc}
\usepackage[T1]{fontenc}
\usepackage[brazil]{babel}
\usepackage{geometry}
\geometry{margin=2.2cm}
\usepackage{enumitem}
\usepackage{amsmath,amssymb}
\usepackage{xcolor}
\usepackage{hyperref}
% Using standard fonts available on any TeX distribution
\usepackage{mathpazo}
\usepackage[scaled=0.9]{helvet}
\usepackage{courier}
\renewcommand{\familydefault}{\rmdefault}

% ── Conditional: gabarito on/off ──
\ifdefined\GABARITO
  \ifnum\GABARITO=1
    \newcommand{\sol}[1]{\textcolor{blue}{#1}}
    \newcommand{\errocomum}[1]{\textcolor{orange}{⚠️ Erro comum: #1}}
    \newcommand{\competencia}[1]{\textcolor{gray}{Competência: #1}}
  \else
    \newcommand{\sol}[1]{}
    \newcommand{\errocomum}[1]{}
    \newcommand{\competencia}[1]{}
  \fi
\else
  \newcommand{\sol}[1]{}
  \newcommand{\errocomum}[1]{}
  \newcommand{\competencia}[1]{}
\fi

% ── Question formatting ──
\newcommand{\nivel}[1]{\textcolor{gray}{[#1]}}
\newenvironment{questao}[2]{
  \vspace{0.5cm}
  \noindent\textbf{Q#1} \nivel{#2} \quad
}{
  \vspace{0.2cm}
}

\begin{document}

% ════════════════ HEADER ════════════════
\begin{center}
  {\sffamily\Large GBC063 · Semana 1 · Inteligência Artificial}

  \vspace{0.2cm}
  {\small Prof. Marcelo Keese Albertini · FACOM · UFU}

  \vspace{0.2cm}
  {\footnotesize Cobre as competências do resumo S01 ·
  \href{../semana01/material_introducao_ia.html}{material de estudo} ·
  \href{../semana01/resumo_s01_introducao_ia.md}{resumo}}

  \ifdefined\GABARITO
    \ifnum\GABARITO=1
      \vspace{0.3cm}
      \fbox{\textcolor{red}{\textbf{GABARITO — Documento Exclusivo do Professor}}}
    \fi
  \fi
\end{center}

\vspace{0.3cm}
\hrule
\vspace{0.5cm}

% ════════════════ NÍVEL ● (REATIVAÇÃO) ════════════════
\section*{Nível ● \quad Reativação}

\begin{questao}{1}{●}
  Defina \textbf{agente racional} segundo Russell \& Norvig e explique,
  em no máximo 3 linhas, por que esta definição — e não o Teste de Turing —
  é adotada como definição operacional do curso.

  \sol{
    \textbf{Resposta:} Um agente racional é aquele que, para cada percepção possível,
    seleciona a ação que maximiza a métrica de desempenho esperada. O curso adota esta
    definição porque ela (1) funciona com qualquer métrica de sucesso, (2) não depende
    de replicar comportamento humano, e (3) unifica desde um termostato até o AlphaGo.

    \textbf{Resolução:} Russell \& Norvig (2013), Cap. 2. A definição de agente racional
    é independente de humanidade — um agente pode ser racional e não-humano (ex.: AlphaGo
    joga Go de forma sobre-humana).

    \errocomum{Definir agente racional como "agente que age como um humano". Isso é a visão
    "agir humanamente" (Teste de Turing), não "agir racionalmente".}

    \competencia{Explicar a diferença entre IA simbólica e aprendizado de máquina —
    e definir agente racional.}
  }
\end{questao}

\begin{questao}{2}{●}
  As quatro visões clássicas da IA organizam-se em uma matriz 2×2:
  pensamento \textit{vs.} ação × desempenho humano \textit{vs.} racionalidade.
  Classifique cada um dos sistemas abaixo em \textbf{um} dos quatro quadrantes
  e justifique em 1 linha:

  \begin{enumerate}[label=(\alph*),nosep]
    \item ELIZA (Weizenbaum, 1966)
    \item Deep Blue (IBM, 1997)
    \item O agente aspirador de pó da Aula B
    \item Um modelo de linguagem treinado para prever o próximo token
  \end{enumerate}

  \sol{
    \textbf{Respostas:}
    (a) Agir humanamente — ELIZA simulava um terapeuta humano sem raciocinar.
    (b) Agir racionalmente — Deep Blue selecionava lances que maximizavam chance de vitória,
    sem imitar o pensamento de um grande mestre humano.
    (c) Agir racionalmente — o aspirador maximiza a limpeza dada sua percepção.
    (d) Pensar racionalmente (durante o treino: otimiza uma função de perda) e
    agir humanamente (na inferência: gera texto indistinguível do humano).

    \errocomum{Classificar ELIZA como "pensar humanamente". ELIZA não modelava pensamento —
    apenas casava padrões de string.}

    \competencia{Classificar qualquer sistema de IA nas quatro visões e justificar.}
  }
\end{questao}

\begin{questao}{3}{●}
  Verdadeiro ou falso. \textbf{Justifique cada resposta em até 2 linhas.}

  \begin{enumerate}[label=(\alph*),nosep]
    \item "O primeiro inverno da IA (1969–1980) foi causado exclusivamente por falta
    de poder computacional."
    \item "Redes neurais artificiais são uma ideia dos anos 2010."
  \end{enumerate}

  \sol{
    \textbf{Respostas:}
    (a) \textbf{Falso.} O inverno foi causado principalmente pelo livro \textit{Perceptrons}
    (Minsky \& Papert, 1969), que provou limitações do perceptron simples e levou o campo
    a concluir — erroneamente — que redes neurais eram um beco sem saída. A falta de
    computação contribuiu, mas o fator decisivo foi a perda de credibilidade e financiamento.

    (b) \textbf{Falso.} O primeiro modelo de neurônio artificial é de McCulloch \& Pitts (1943).
    O perceptron é de Rosenblatt (1959). O backpropagation que as treina é de 1986
    (Rumelhart, Hinton \& Williams). As ideias fundamentais têm 40–80 anos; o que mudou
    em 2012 foi a escala (dados + GPUs).

    \errocomum{Para (a): atribuir o inverno só à falta de hardware. O fator principal foi
    sociológico — a comunidade perdeu confiança nas redes neurais.}
    \errocomum{Para (b): achar que deep learning é uma invenção recente. As ideias são antigas;
    o que é novo é a escala.}

    \competencia{Enunciar a Bitter Lesson e dar exemplos históricos.}
  }
\end{questao}

% ════════════════ NÍVEL ●● (EXECUÇÃO) ════════════════
\section*{Nível ●● \quad Execução}

\begin{questao}{4}{●●}
  Considere o seguinte ambiente de aspirador de pó com \textbf{três} cômodos
  em linha (A–B–C). O agente começa na posição B. O cômodo A está sujo; B e C
  estão limpos. O agente segue exatamente a lógica abaixo:

  \begin{verbatim}
  def agente_aspirador(percepcao):
      posicao, sujo = percepcao
      if sujo:        return "ASPIRAR"
      if posicao == "A": return "DIREITA"
      if posicao == "B": return "DIREITA"
      return "ESQUERDA"
  \end{verbatim}

  Preencha a tabela de execução para os \textbf{6 primeiros passos}, indicando:
  posição atual, sensor (sujo/limpo), ação tomada, e estado dos cômodos após a ação.
  O agente limpa todos os cômodos? Em quantos passos?

  \sol{
    \textbf{Resposta:}

    \begin{tabular}{c|c|c|c|c}
      Passo & Posição & Sensor & Ação & Cômodos [A,B,C] \\
      \hline
      0 & B & limpo & DIREITA & [sujo, limpo, limpo] \\
      1 & C & limpo & ESQUERDA & [sujo, limpo, limpo] \\
      2 & B & limpo & DIREITA & [sujo, limpo, limpo] \\
      3 & C & limpo & ESQUERDA & [sujo, limpo, limpo] \\
      \dots & \dots & \dots & \dots & \dots \\
    \end{tabular}

    O agente \textbf{nunca alcança o cômodo A} porque sua lógica só vai de A→B→C→B→C→B→C\ldots
    Ele fica preso no ciclo B↔C. Para alcançar A, seria preciso uma regra adicional
    (ex.: "se estou em C e ambos A e B estão sujos, vá para B").

    \textbf{Resolução:} O bug está na lógica de movimento: quando a posição é B ou C,
    o agente nunca se move para A. A condição \texttt{posicao == "A"} nunca é verdadeira
    se o agente não consegue chegar a A.

    \errocomum{Assumir que o agente "dá a volta". O ambiente é uma linha, não um ciclo —
    de C só se vai para B.}

    \competencia{Implementar o agente aspirador e prever seu comportamento
    com 3 cômodos.}
  }
\end{questao}

\begin{questao}{5}{●●}
  O motor de encadeamento progressivo abaixo recebe uma base de fatos inicial
  e um conjunto de regras. Execute o algoritmo \textbf{passo a passo} com os
  dados fornecidos e indique o conjunto final de fatos.

  \begin{verbatim}
  fatos = {"chovendo", "temperatura < 15"}
  regras = [
    ({"chovendo", "temperatura < 15"}, "risco de gelo"),
    ({"risco de gelo"},              "fechar pista"),
    ({"chovendo"},                   "pista molhada"),
  ]
  \end{verbatim}

  \sol{
    \textbf{Resposta:} fatos finais = \{\texttt{"chovendo"}, \texttt{"temperatura < 15"},
    \texttt{"risco de gelo"}, \texttt{"fechar pista"}, \texttt{"pista molhada"}\}

    \textbf{Resolução:}
    \begin{enumerate}[nosep]
      \item Iteração 1: chovendo + temperatura < 15 → adiciona \texttt{"risco de gelo"}.
      \item Iteração 2: risco de gelo → adiciona \texttt{"fechar pista"}.
      \item Iteração 3: chovendo → adiciona \texttt{"pista molhada"}.
      \item Iteração 4: nenhuma regra nova dispara → fim.
    \end{enumerate}

    \errocomum{Esquecer que as regras disparam em qualquer ordem — o algoritmo testa TODAS
    as regras a cada iteração, não para na primeira que dispara.}

    \competencia{Explicar o funcionamento de um motor de inferência e por que regras
    não escalam.}
  }
\end{questao}

\begin{questao}{6}{●●}
  Em 2023, um LLM foi capaz de passar no exame da ordem dos advogados (Bar Exam)
  dos EUA sem ter recebido uma única regra explícita sobre direito. Usando os
  conceitos da Semana 1, explique \textbf{por que isso é um fato tão significativo}
  na história da IA. Se não obteve uma resposta completa, use o roteiro:
  (1) como um sistema especialista resolveria o Bar Exam? (2) por que essa abordagem
  falharia? (3) o que o LLM fez de diferente?

  \sol{
    \textbf{Resposta:} Um sistema especialista exigiria que engenheiros de conhecimento
    codificassem milhares de regras jurídicas à mão — um gargalo intransponível. O LLM
    não recebeu regra nenhuma: aprendeu padrões estatísticos de centenas de gigabytes
    de texto jurídico e geral. Isso é a Bitter Lesson em ação: conhecimento aprendido
    de dados venceu conhecimento artesanal — em um domínio (direito) que se supunha
    exigir raciocínio simbólico explícito.

    \errocomum{Dizer que o LLM "decorou" as respostas. LLMs não têm memória explícita do
    corpus de treino — eles aprendem padrões estatísticos, não fazem lookup de frases.}

    \competencia{Enunciar a Bitter Lesson e dar exemplos que a corroboram.}
  }
\end{questao}

% ════════════════ NÍVEL ●●● (ANÁLISE) ════════════════
\section*{Nível ●●● \quad Análise}

\begin{questao}{7}{●●●}
  Considere a afirmação de Sutton (2019):

  \begin{quote}
  \textit{``Métodos gerais que alavancam computação são, em última instância,
  os mais efetivos — e por uma grande margem. A maioria do esforço em IA é
  gasto tentando construir conhecimento em agentes; isso ajuda no curto prazo,
  mas no longo prazo atinge um platô.''}
  \end{quote}

  \begin{enumerate}[label=(\alph*),nosep]
    \item Por que conhecimento artesanal \textit{necessariamente} atinge um platô,
    enquanto métodos que aprendem de dados não? (Dica: pense no que acontece quando
    você dobra a quantidade de dados ou de computação em cada caso.)
    \item Dê um exemplo \textbf{concreto} de um sistema que atingiu esse platô e
    explique \textit{o que especificamente} o limitou.
  \end{enumerate}

  \sol{
    \textbf{Respostas:}
    (a) Conhecimento artesanal atinge um platô porque sua melhoria depende de \textbf{humanos
    adicionando mais regras} — um processo que é linear, caro e sujeito a contradições.
    Métodos que aprendem melhoram quando você adiciona mais dados ou computação, e esse
    processo é (em princípio) ilimitado: mais dados → melhores gradientes → melhor modelo.
    Não há um humano no loop ditando o teto.

    (b) Sistemas especialistas médicos (ex.: MYCIN, anos 80). Eram melhores que médicos
    juniores em diagnósticos específicos, mas cada nova doença exigia novas regras escritas
    por especialistas. Quando o número de regras passou de centenas, as interações entre
    elas tornaram-se imprevisíveis. O sistema não melhorava com mais dados de pacientes —
    só com mais regras escritas à mão. Enquanto isso, métodos de ML atuais melhoram
    automaticamente com mais exames, prontuários e artigos.

    \errocomum{Em (a), dizer que conhecimento artesanal atinge platô "porque humanos
    são limitados". A resposta precisa identificar o mecanismo: a taxa de melhoria é
    ditada pela velocidade com que humanos escrevem regras, não pela velocidade com
    que computação cresce.}

    \competencia{Enunciar a Bitter Lesson e explicar o mecanismo por trás dela.}
  }
\end{questao}

\begin{questao}{8}{●●●}
  O efeito ELIZA — atribuir inteligência onde há apenas casamento de padrões —
  foi documentado em 1966. Ele persiste em 2025 com LLMs?

  \begin{enumerate}[label=(\alph*),nosep]
    \item Aponte \textbf{uma semelhança} entre ELIZA e um LLM moderno que torna
    o efeito ELIZA relevante hoje.
    \item Aponte \textbf{uma diferença fundamental} que torna a comparação enganosa.
    \item Em no máximo 4 linhas: um LLM é ``só casamento de padrões''? Justifique.
  \end{enumerate}

  \sol{
    \textbf{Respostas:}
    (a) Ambos produzem texto fluente sem garantia de compreensão do conteúdo semântico.
    Um LLM pode gerar uma explicação convincente de um conceito que ele não ``entende''
    no sentido humano — assim como ELIZA simulava empatia sem sentir nada. O efeito
    ELIZA ocorre quando o usuário projeta compreensão onde há apenas correlação estatística.

    (b) ELIZA usava algumas dezenas de regras de reescrita escritas à mão. Um LLM aprende
    de centenas de bilhões de parâmetros ajustados sobre trilhões de tokens. A escala e
    a origem do conhecimento (artesanal vs. aprendido) são de natureza completamente diferente.

    (c) Não. Reduzir um LLM a ``só casamento de padrões'' é como reduzir um cérebro a
    ``só sinapses disparando'' — tecnicamente verdadeiro, mas enganoso. A questão relevante
    não é se o mecanismo básico é simples (é: multiplicação de matrizes + não-linearidades),
    mas se a escala produz comportamentos emergentes que merecem ser chamados de raciocínio.
    Esta é uma pergunta em aberto na pesquisa em 2025.

    \errocomum{Responder (c) com um simples ``sim'' ou ``não''. A pergunta exige distinguir
    mecanismo de comportamento — o que o sistema \textit{faz} vs. como ele \textit{faz}.}

    \competencia{Classificar sistemas de IA e analisar criticamente alegações de inteligência.}
  }
\end{questao}

% ════════════════ NÍVEL ★ (SÍNTESE) ════════════════
\section*{Nível ★ \quad Síntese}

\begin{questao}{9}{★}
  Projete um experimento mental: você herda um sistema especialista médico com
  10.000 regras, desenvolvido ao longo de 15 anos por uma equipe de 20 engenheiros
  de conhecimento. O sistema tem 92\% de acurácia em diagnóstico.

  \begin{enumerate}[label=(\alph*),nosep]
    \item Identifique \textbf{três} cenários em que este sistema falharia de forma
    sistemática que um modelo de ML treinado sobre os mesmos dados não falharia.
    \item Para cada cenário, explique \textbf{por que} a natureza artesanal das
    regras é a causa raiz da falha.
    \item Proponha uma arquitetura híbrida (simbólico + aprendizado) que combine
    o melhor dos dois mundos para este domínio. Descreva em no máximo 6 linhas.
  \end{enumerate}

  \sol{
    \textbf{Respostas:}
    (a) e (b):
    \begin{enumerate}[nosep]
      \item \textbf{Nova doença:} o SE falha porque nenhuma regra cobre o fenômeno novo.
      O ML pode detectar que os sintomas formam um cluster distinto nos dados e sinalizar
      uma anomalia — mesmo sem rótulos. Causa raiz: SE só reconhece o que foi explicitamente
      programado.
      \item \textbf{Interação entre condições:} um paciente com diabetes \textit{e} COVID
      pode ativar regras contraditórias (uma recomenda corticoides, outra contraindica).
      O SE não tem como ponderar evidências conflitantes. O ML aprende pesos relativos
      dos features. Causa raiz: regras são discretas e não-compostas; não modelam interações.
      \item \textbf{Deriva de distribuição:} se a população de pacientes mudar (ex.:
      envelhecimento), a acurácia do SE cai porque as regras foram escritas para a
      distribuição antiga. O ML pode ser retreinado com dados novos. Causa raiz: SE
      não se adapta sem reengenharia manual.
    \end{enumerate}

    (c) \textbf{Arquitetura híbrida:} usar um modelo de ML (ex.: gradient-boosted trees
    ou rede neural) para produzir o diagnóstico base e a estimativa de confiança. As
    10.000 regras viram um sistema de \textbf{verificação}: antes de emitir o diagnóstico,
    o sistema checa se ele viola alguma regra de segurança crítica (ex.: ``nunca prescrever
    droga X para grávidas''). Se violar, escala para um médico humano. O ML lida com a
    generalização; as regras lidam com as restrições rígidas de segurança.

    \errocomum{Propor um híbrido em que as regras \textit{guiam} o aprendizado. Isso
    reintroduz o gargalo: as regras limitam o que o modelo pode aprender. O melhor
    híbrido usa regras como \textit{restrições de segurança}, não como indutores de viés.}

    \competencia{Sintetizar os conceitos de IA simbólica e aprendizado em uma arquitetura
    fundamentada.}
  }
\end{questao}

\begin{questao}{10}{★ 🌉}
  O treinamento do GPT usa a mesma estrutura de loop que o agente aspirador da Aula B:
  \textbf{percepção → decisão → ação}. A diferença está na complexidade de \texttt{decidir()}.

  \begin{enumerate}[label=(\alph*),nosep]
    \item No aspirador, \texttt{decidir()} é uma função com 4 regras \texttt{if}.
    No GPT, \texttt{decidir()} é uma rede neural com centenas de bilhões de
    parâmetros. \textbf{Por que}, conceitualmente, a mesma interface de 4 linhas
    funciona para ambos? O que isso revela sobre o paradigma de agentes?
    \item A Bitter Lesson argumenta que conhecimento artesanal atinge um platô.
    O prompt engineering — escrever instruções cuidadosas para um LLM — é uma
    forma de conhecimento artesanal? Justifique em até 5 linhas.
  \end{enumerate}

  \sol{
    \textbf{Respostas:}
    (a) A interface \texttt{obs → decidir(obs) → ação} funciona para ambos porque ela
    abstrai \textbf{o que} o agente faz de \textbf{como} ele decide. O paradigma de
    agentes é uma \textit{casca}: define a assinatura da decisão, não sua implementação.
    O que muda entre o aspirador e o GPT é apenas a função \texttt{decidir()} — e essa
    função pode ser desde 4 ifs até uma rede de 175 bilhões de parâmetros. Isso revela
    que o paradigma de agentes é \textbf{universal}: ele unifica toda a IA sob uma
    única abstração.

    (b) Sim, prompt engineering é conhecimento artesanal — mas de um tipo novo. É
    artesanal porque um humano está manualmente ajustando a entrada para melhorar
    a saída, sem que o modelo aprenda com esse ajuste. Mas difere do artesanal
    clássico porque (1) o modelo subjacente foi aprendido de dados, não programado,
    e (2) o custo de ajustar um prompt é ordens de grandeza menor que o de escrever
    regras. Dito isso, a Bitter Lesson sugere que, no longo prazo, métodos que
    \textit{aprendem} a fazer prompting (auto-prompting, RLHF) substituirão o
    prompting manual — e já há evidências disso.

    \errocomum{Em (b), responder simplesmente ``não'' porque o modelo foi treinado com ML.
    O prompt em si é conhecimento artesanal injetado \textit{na inferência}, não no treino.}

    \competencia{Reconhecer o loop percepção→decisão→ação como a interface universal
    da IA e aplicar a Bitter Lesson a desenvolvimentos recentes.}
  }
\end{questao}

\vspace{1cm}
\hrule
\vspace{0.3cm}
\begin{center}
  {\footnotesize Lista S01 · GBC063 · Prof. Marcelo Keese Albertini · FACOM/UFU}
  \ifdefined\GABARITO
    \ifnum\GABARITO=1
      \\[0.2cm] {\footnotesize\color{red}\textbf{GABARITO — USO EXCLUSIVO DO PROFESSOR}}
    \fi
  \fi
\end{center}

\end{document}
"
% ═══════════════════════════════════════════════════════════════

\documentclass[12pt,a4paper]{article}

\usepackage[utf8]{inputenc}
\usepackage[T1]{fontenc}
\usepackage[brazil]{babel}
\usepackage{geometry}
\geometry{margin=2.2cm}
\usepackage{enumitem}
\usepackage{amsmath,amssymb}
\usepackage{xcolor}
\usepackage{hyperref}
% Using standard fonts available on any TeX distribution
\usepackage{mathpazo}
\usepackage[scaled=0.9]{helvet}
\usepackage{courier}
\renewcommand{\familydefault}{\rmdefault}

% ── Conditional: gabarito on/off ──
\ifdefined\GABARITO
  \ifnum\GABARITO=1
    \newcommand{\sol}[1]{\textcolor{blue}{#1}}
    \newcommand{\errocomum}[1]{\textcolor{orange}{⚠️ Erro comum: #1}}
    \newcommand{\competencia}[1]{\textcolor{gray}{Competência: #1}}
  \else
    \newcommand{\sol}[1]{}
    \newcommand{\errocomum}[1]{}
    \newcommand{\competencia}[1]{}
  \fi
\else
  \newcommand{\sol}[1]{}
  \newcommand{\errocomum}[1]{}
  \newcommand{\competencia}[1]{}
\fi

% ── Question formatting ──
\newcommand{\nivel}[1]{\textcolor{gray}{[#1]}}
\newenvironment{questao}[2]{
  \vspace{0.5cm}
  \noindent\textbf{Q#1} \nivel{#2} \quad
}{
  \vspace{0.2cm}
}

\begin{document}

% ════════════════ HEADER ════════════════
\begin{center}
  {\sffamily\Large GBC063 · Semana 1 · Inteligência Artificial}

  \vspace{0.2cm}
  {\small Prof. Marcelo Keese Albertini · FACOM · UFU}

  \vspace{0.2cm}
  {\footnotesize Cobre as competências do resumo S01 ·
  \href{../semana01/material_introducao_ia.html}{material de estudo} ·
  \href{../semana01/resumo_s01_introducao_ia.md}{resumo}}

  \ifdefined\GABARITO
    \ifnum\GABARITO=1
      \vspace{0.3cm}
      \fbox{\textcolor{red}{\textbf{GABARITO — Documento Exclusivo do Professor}}}
    \fi
  \fi
\end{center}

\vspace{0.3cm}
\hrule
\vspace{0.5cm}

% ════════════════ NÍVEL ● (REATIVAÇÃO) ════════════════
\section*{Nível ● \quad Reativação}

\begin{questao}{1}{●}
  Defina \textbf{agente racional} segundo Russell \& Norvig e explique,
  em no máximo 3 linhas, por que esta definição — e não o Teste de Turing —
  é adotada como definição operacional do curso.

  \sol{
    \textbf{Resposta:} Um agente racional é aquele que, para cada percepção possível,
    seleciona a ação que maximiza a métrica de desempenho esperada. O curso adota esta
    definição porque ela (1) funciona com qualquer métrica de sucesso, (2) não depende
    de replicar comportamento humano, e (3) unifica desde um termostato até o AlphaGo.

    \textbf{Resolução:} Russell \& Norvig (2013), Cap. 2. A definição de agente racional
    é independente de humanidade — um agente pode ser racional e não-humano (ex.: AlphaGo
    joga Go de forma sobre-humana).

    \errocomum{Definir agente racional como "agente que age como um humano". Isso é a visão
    "agir humanamente" (Teste de Turing), não "agir racionalmente".}

    \competencia{Explicar a diferença entre IA simbólica e aprendizado de máquina —
    e definir agente racional.}
  }
\end{questao}

\begin{questao}{2}{●}
  As quatro visões clássicas da IA organizam-se em uma matriz 2×2:
  pensamento \textit{vs.} ação × desempenho humano \textit{vs.} racionalidade.
  Classifique cada um dos sistemas abaixo em \textbf{um} dos quatro quadrantes
  e justifique em 1 linha:

  \begin{enumerate}[label=(\alph*),nosep]
    \item ELIZA (Weizenbaum, 1966)
    \item Deep Blue (IBM, 1997)
    \item O agente aspirador de pó da Aula B
    \item Um modelo de linguagem treinado para prever o próximo token
  \end{enumerate}

  \sol{
    \textbf{Respostas:}
    (a) Agir humanamente — ELIZA simulava um terapeuta humano sem raciocinar.
    (b) Agir racionalmente — Deep Blue selecionava lances que maximizavam chance de vitória,
    sem imitar o pensamento de um grande mestre humano.
    (c) Agir racionalmente — o aspirador maximiza a limpeza dada sua percepção.
    (d) Pensar racionalmente (durante o treino: otimiza uma função de perda) e
    agir humanamente (na inferência: gera texto indistinguível do humano).

    \errocomum{Classificar ELIZA como "pensar humanamente". ELIZA não modelava pensamento —
    apenas casava padrões de string.}

    \competencia{Classificar qualquer sistema de IA nas quatro visões e justificar.}
  }
\end{questao}

\begin{questao}{3}{●}
  Verdadeiro ou falso. \textbf{Justifique cada resposta em até 2 linhas.}

  \begin{enumerate}[label=(\alph*),nosep]
    \item "O primeiro inverno da IA (1969–1980) foi causado exclusivamente por falta
    de poder computacional."
    \item "Redes neurais artificiais são uma ideia dos anos 2010."
  \end{enumerate}

  \sol{
    \textbf{Respostas:}
    (a) \textbf{Falso.} O inverno foi causado principalmente pelo livro \textit{Perceptrons}
    (Minsky \& Papert, 1969), que provou limitações do perceptron simples e levou o campo
    a concluir — erroneamente — que redes neurais eram um beco sem saída. A falta de
    computação contribuiu, mas o fator decisivo foi a perda de credibilidade e financiamento.

    (b) \textbf{Falso.} O primeiro modelo de neurônio artificial é de McCulloch \& Pitts (1943).
    O perceptron é de Rosenblatt (1959). O backpropagation que as treina é de 1986
    (Rumelhart, Hinton \& Williams). As ideias fundamentais têm 40–80 anos; o que mudou
    em 2012 foi a escala (dados + GPUs).

    \errocomum{Para (a): atribuir o inverno só à falta de hardware. O fator principal foi
    sociológico — a comunidade perdeu confiança nas redes neurais.}
    \errocomum{Para (b): achar que deep learning é uma invenção recente. As ideias são antigas;
    o que é novo é a escala.}

    \competencia{Enunciar a Bitter Lesson e dar exemplos históricos.}
  }
\end{questao}

% ════════════════ NÍVEL ●● (EXECUÇÃO) ════════════════
\section*{Nível ●● \quad Execução}

\begin{questao}{4}{●●}
  Considere o seguinte ambiente de aspirador de pó com \textbf{três} cômodos
  em linha (A–B–C). O agente começa na posição B. O cômodo A está sujo; B e C
  estão limpos. O agente segue exatamente a lógica abaixo:

  \begin{verbatim}
  def agente_aspirador(percepcao):
      posicao, sujo = percepcao
      if sujo:        return "ASPIRAR"
      if posicao == "A": return "DIREITA"
      if posicao == "B": return "DIREITA"
      return "ESQUERDA"
  \end{verbatim}

  Preencha a tabela de execução para os \textbf{6 primeiros passos}, indicando:
  posição atual, sensor (sujo/limpo), ação tomada, e estado dos cômodos após a ação.
  O agente limpa todos os cômodos? Em quantos passos?

  \sol{
    \textbf{Resposta:}

    \begin{tabular}{c|c|c|c|c}
      Passo & Posição & Sensor & Ação & Cômodos [A,B,C] \\
      \hline
      0 & B & limpo & DIREITA & [sujo, limpo, limpo] \\
      1 & C & limpo & ESQUERDA & [sujo, limpo, limpo] \\
      2 & B & limpo & DIREITA & [sujo, limpo, limpo] \\
      3 & C & limpo & ESQUERDA & [sujo, limpo, limpo] \\
      \dots & \dots & \dots & \dots & \dots \\
    \end{tabular}

    O agente \textbf{nunca alcança o cômodo A} porque sua lógica só vai de A→B→C→B→C→B→C\ldots
    Ele fica preso no ciclo B↔C. Para alcançar A, seria preciso uma regra adicional
    (ex.: "se estou em C e ambos A e B estão sujos, vá para B").

    \textbf{Resolução:} O bug está na lógica de movimento: quando a posição é B ou C,
    o agente nunca se move para A. A condição \texttt{posicao == "A"} nunca é verdadeira
    se o agente não consegue chegar a A.

    \errocomum{Assumir que o agente "dá a volta". O ambiente é uma linha, não um ciclo —
    de C só se vai para B.}

    \competencia{Implementar o agente aspirador e prever seu comportamento
    com 3 cômodos.}
  }
\end{questao}

\begin{questao}{5}{●●}
  O motor de encadeamento progressivo abaixo recebe uma base de fatos inicial
  e um conjunto de regras. Execute o algoritmo \textbf{passo a passo} com os
  dados fornecidos e indique o conjunto final de fatos.

  \begin{verbatim}
  fatos = {"chovendo", "temperatura < 15"}
  regras = [
    ({"chovendo", "temperatura < 15"}, "risco de gelo"),
    ({"risco de gelo"},              "fechar pista"),
    ({"chovendo"},                   "pista molhada"),
  ]
  \end{verbatim}

  \sol{
    \textbf{Resposta:} fatos finais = \{\texttt{"chovendo"}, \texttt{"temperatura < 15"},
    \texttt{"risco de gelo"}, \texttt{"fechar pista"}, \texttt{"pista molhada"}\}

    \textbf{Resolução:}
    \begin{enumerate}[nosep]
      \item Iteração 1: chovendo + temperatura < 15 → adiciona \texttt{"risco de gelo"}.
      \item Iteração 2: risco de gelo → adiciona \texttt{"fechar pista"}.
      \item Iteração 3: chovendo → adiciona \texttt{"pista molhada"}.
      \item Iteração 4: nenhuma regra nova dispara → fim.
    \end{enumerate}

    \errocomum{Esquecer que as regras disparam em qualquer ordem — o algoritmo testa TODAS
    as regras a cada iteração, não para na primeira que dispara.}

    \competencia{Explicar o funcionamento de um motor de inferência e por que regras
    não escalam.}
  }
\end{questao}

\begin{questao}{6}{●●}
  Em 2023, um LLM foi capaz de passar no exame da ordem dos advogados (Bar Exam)
  dos EUA sem ter recebido uma única regra explícita sobre direito. Usando os
  conceitos da Semana 1, explique \textbf{por que isso é um fato tão significativo}
  na história da IA. Se não obteve uma resposta completa, use o roteiro:
  (1) como um sistema especialista resolveria o Bar Exam? (2) por que essa abordagem
  falharia? (3) o que o LLM fez de diferente?

  \sol{
    \textbf{Resposta:} Um sistema especialista exigiria que engenheiros de conhecimento
    codificassem milhares de regras jurídicas à mão — um gargalo intransponível. O LLM
    não recebeu regra nenhuma: aprendeu padrões estatísticos de centenas de gigabytes
    de texto jurídico e geral. Isso é a Bitter Lesson em ação: conhecimento aprendido
    de dados venceu conhecimento artesanal — em um domínio (direito) que se supunha
    exigir raciocínio simbólico explícito.

    \errocomum{Dizer que o LLM "decorou" as respostas. LLMs não têm memória explícita do
    corpus de treino — eles aprendem padrões estatísticos, não fazem lookup de frases.}

    \competencia{Enunciar a Bitter Lesson e dar exemplos que a corroboram.}
  }
\end{questao}

% ════════════════ NÍVEL ●●● (ANÁLISE) ════════════════
\section*{Nível ●●● \quad Análise}

\begin{questao}{7}{●●●}
  Considere a afirmação de Sutton (2019):

  \begin{quote}
  \textit{``Métodos gerais que alavancam computação são, em última instância,
  os mais efetivos — e por uma grande margem. A maioria do esforço em IA é
  gasto tentando construir conhecimento em agentes; isso ajuda no curto prazo,
  mas no longo prazo atinge um platô.''}
  \end{quote}

  \begin{enumerate}[label=(\alph*),nosep]
    \item Por que conhecimento artesanal \textit{necessariamente} atinge um platô,
    enquanto métodos que aprendem de dados não? (Dica: pense no que acontece quando
    você dobra a quantidade de dados ou de computação em cada caso.)
    \item Dê um exemplo \textbf{concreto} de um sistema que atingiu esse platô e
    explique \textit{o que especificamente} o limitou.
  \end{enumerate}

  \sol{
    \textbf{Respostas:}
    (a) Conhecimento artesanal atinge um platô porque sua melhoria depende de \textbf{humanos
    adicionando mais regras} — um processo que é linear, caro e sujeito a contradições.
    Métodos que aprendem melhoram quando você adiciona mais dados ou computação, e esse
    processo é (em princípio) ilimitado: mais dados → melhores gradientes → melhor modelo.
    Não há um humano no loop ditando o teto.

    (b) Sistemas especialistas médicos (ex.: MYCIN, anos 80). Eram melhores que médicos
    juniores em diagnósticos específicos, mas cada nova doença exigia novas regras escritas
    por especialistas. Quando o número de regras passou de centenas, as interações entre
    elas tornaram-se imprevisíveis. O sistema não melhorava com mais dados de pacientes —
    só com mais regras escritas à mão. Enquanto isso, métodos de ML atuais melhoram
    automaticamente com mais exames, prontuários e artigos.

    \errocomum{Em (a), dizer que conhecimento artesanal atinge platô "porque humanos
    são limitados". A resposta precisa identificar o mecanismo: a taxa de melhoria é
    ditada pela velocidade com que humanos escrevem regras, não pela velocidade com
    que computação cresce.}

    \competencia{Enunciar a Bitter Lesson e explicar o mecanismo por trás dela.}
  }
\end{questao}

\begin{questao}{8}{●●●}
  O efeito ELIZA — atribuir inteligência onde há apenas casamento de padrões —
  foi documentado em 1966. Ele persiste em 2025 com LLMs?

  \begin{enumerate}[label=(\alph*),nosep]
    \item Aponte \textbf{uma semelhança} entre ELIZA e um LLM moderno que torna
    o efeito ELIZA relevante hoje.
    \item Aponte \textbf{uma diferença fundamental} que torna a comparação enganosa.
    \item Em no máximo 4 linhas: um LLM é ``só casamento de padrões''? Justifique.
  \end{enumerate}

  \sol{
    \textbf{Respostas:}
    (a) Ambos produzem texto fluente sem garantia de compreensão do conteúdo semântico.
    Um LLM pode gerar uma explicação convincente de um conceito que ele não ``entende''
    no sentido humano — assim como ELIZA simulava empatia sem sentir nada. O efeito
    ELIZA ocorre quando o usuário projeta compreensão onde há apenas correlação estatística.

    (b) ELIZA usava algumas dezenas de regras de reescrita escritas à mão. Um LLM aprende
    de centenas de bilhões de parâmetros ajustados sobre trilhões de tokens. A escala e
    a origem do conhecimento (artesanal vs. aprendido) são de natureza completamente diferente.

    (c) Não. Reduzir um LLM a ``só casamento de padrões'' é como reduzir um cérebro a
    ``só sinapses disparando'' — tecnicamente verdadeiro, mas enganoso. A questão relevante
    não é se o mecanismo básico é simples (é: multiplicação de matrizes + não-linearidades),
    mas se a escala produz comportamentos emergentes que merecem ser chamados de raciocínio.
    Esta é uma pergunta em aberto na pesquisa em 2025.

    \errocomum{Responder (c) com um simples ``sim'' ou ``não''. A pergunta exige distinguir
    mecanismo de comportamento — o que o sistema \textit{faz} vs. como ele \textit{faz}.}

    \competencia{Classificar sistemas de IA e analisar criticamente alegações de inteligência.}
  }
\end{questao}

% ════════════════ NÍVEL ★ (SÍNTESE) ════════════════
\section*{Nível ★ \quad Síntese}

\begin{questao}{9}{★}
  Projete um experimento mental: você herda um sistema especialista médico com
  10.000 regras, desenvolvido ao longo de 15 anos por uma equipe de 20 engenheiros
  de conhecimento. O sistema tem 92\% de acurácia em diagnóstico.

  \begin{enumerate}[label=(\alph*),nosep]
    \item Identifique \textbf{três} cenários em que este sistema falharia de forma
    sistemática que um modelo de ML treinado sobre os mesmos dados não falharia.
    \item Para cada cenário, explique \textbf{por que} a natureza artesanal das
    regras é a causa raiz da falha.
    \item Proponha uma arquitetura híbrida (simbólico + aprendizado) que combine
    o melhor dos dois mundos para este domínio. Descreva em no máximo 6 linhas.
  \end{enumerate}

  \sol{
    \textbf{Respostas:}
    (a) e (b):
    \begin{enumerate}[nosep]
      \item \textbf{Nova doença:} o SE falha porque nenhuma regra cobre o fenômeno novo.
      O ML pode detectar que os sintomas formam um cluster distinto nos dados e sinalizar
      uma anomalia — mesmo sem rótulos. Causa raiz: SE só reconhece o que foi explicitamente
      programado.
      \item \textbf{Interação entre condições:} um paciente com diabetes \textit{e} COVID
      pode ativar regras contraditórias (uma recomenda corticoides, outra contraindica).
      O SE não tem como ponderar evidências conflitantes. O ML aprende pesos relativos
      dos features. Causa raiz: regras são discretas e não-compostas; não modelam interações.
      \item \textbf{Deriva de distribuição:} se a população de pacientes mudar (ex.:
      envelhecimento), a acurácia do SE cai porque as regras foram escritas para a
      distribuição antiga. O ML pode ser retreinado com dados novos. Causa raiz: SE
      não se adapta sem reengenharia manual.
    \end{enumerate}

    (c) \textbf{Arquitetura híbrida:} usar um modelo de ML (ex.: gradient-boosted trees
    ou rede neural) para produzir o diagnóstico base e a estimativa de confiança. As
    10.000 regras viram um sistema de \textbf{verificação}: antes de emitir o diagnóstico,
    o sistema checa se ele viola alguma regra de segurança crítica (ex.: ``nunca prescrever
    droga X para grávidas''). Se violar, escala para um médico humano. O ML lida com a
    generalização; as regras lidam com as restrições rígidas de segurança.

    \errocomum{Propor um híbrido em que as regras \textit{guiam} o aprendizado. Isso
    reintroduz o gargalo: as regras limitam o que o modelo pode aprender. O melhor
    híbrido usa regras como \textit{restrições de segurança}, não como indutores de viés.}

    \competencia{Sintetizar os conceitos de IA simbólica e aprendizado em uma arquitetura
    fundamentada.}
  }
\end{questao}

\begin{questao}{10}{★ 🌉}
  O treinamento do GPT usa a mesma estrutura de loop que o agente aspirador da Aula B:
  \textbf{percepção → decisão → ação}. A diferença está na complexidade de \texttt{decidir()}.

  \begin{enumerate}[label=(\alph*),nosep]
    \item No aspirador, \texttt{decidir()} é uma função com 4 regras \texttt{if}.
    No GPT, \texttt{decidir()} é uma rede neural com centenas de bilhões de
    parâmetros. \textbf{Por que}, conceitualmente, a mesma interface de 4 linhas
    funciona para ambos? O que isso revela sobre o paradigma de agentes?
    \item A Bitter Lesson argumenta que conhecimento artesanal atinge um platô.
    O prompt engineering — escrever instruções cuidadosas para um LLM — é uma
    forma de conhecimento artesanal? Justifique em até 5 linhas.
  \end{enumerate}

  \sol{
    \textbf{Respostas:}
    (a) A interface \texttt{obs → decidir(obs) → ação} funciona para ambos porque ela
    abstrai \textbf{o que} o agente faz de \textbf{como} ele decide. O paradigma de
    agentes é uma \textit{casca}: define a assinatura da decisão, não sua implementação.
    O que muda entre o aspirador e o GPT é apenas a função \texttt{decidir()} — e essa
    função pode ser desde 4 ifs até uma rede de 175 bilhões de parâmetros. Isso revela
    que o paradigma de agentes é \textbf{universal}: ele unifica toda a IA sob uma
    única abstração.

    (b) Sim, prompt engineering é conhecimento artesanal — mas de um tipo novo. É
    artesanal porque um humano está manualmente ajustando a entrada para melhorar
    a saída, sem que o modelo aprenda com esse ajuste. Mas difere do artesanal
    clássico porque (1) o modelo subjacente foi aprendido de dados, não programado,
    e (2) o custo de ajustar um prompt é ordens de grandeza menor que o de escrever
    regras. Dito isso, a Bitter Lesson sugere que, no longo prazo, métodos que
    \textit{aprendem} a fazer prompting (auto-prompting, RLHF) substituirão o
    prompting manual — e já há evidências disso.

    \errocomum{Em (b), responder simplesmente ``não'' porque o modelo foi treinado com ML.
    O prompt em si é conhecimento artesanal injetado \textit{na inferência}, não no treino.}

    \competencia{Reconhecer o loop percepção→decisão→ação como a interface universal
    da IA e aplicar a Bitter Lesson a desenvolvimentos recentes.}
  }
\end{questao}

\vspace{1cm}
\hrule
\vspace{0.3cm}
\begin{center}
  {\footnotesize Lista S01 · GBC063 · Prof. Marcelo Keese Albertini · FACOM/UFU}
  \ifdefined\GABARITO
    \ifnum\GABARITO=1
      \\[0.2cm] {\footnotesize\color{red}\textbf{GABARITO — USO EXCLUSIVO DO PROFESSOR}}
    \fi
  \fi
\end{center}

\end{document}
"
%   Gabarito:  pdflatex -jobname=gabarito_lista_s01_introducao_ia "\def\GABARITO{1}% ═══════════════════════════════════════════════════════════════
% Lista S01 — O que é Inteligência Artificial?
% Prof. Marcelo Keese Albertini · FACOM · UFU · GBC063
%
% Compilação:
%   Questões:  pdflatex -jobname=lista_s01_introducao_ia "\def\GABARITO{0}% ═══════════════════════════════════════════════════════════════
% Lista S01 — O que é Inteligência Artificial?
% Prof. Marcelo Keese Albertini · FACOM · UFU · GBC063
%
% Compilação:
%   Questões:  pdflatex -jobname=lista_s01_introducao_ia "\def\GABARITO{0}\input{lista_s01_introducao_ia.tex}"
%   Gabarito:  pdflatex -jobname=gabarito_lista_s01_introducao_ia "\def\GABARITO{1}\input{lista_s01_introducao_ia.tex}"
% ═══════════════════════════════════════════════════════════════

\documentclass[12pt,a4paper]{article}

\usepackage[utf8]{inputenc}
\usepackage[T1]{fontenc}
\usepackage[brazil]{babel}
\usepackage{geometry}
\geometry{margin=2.2cm}
\usepackage{enumitem}
\usepackage{amsmath,amssymb}
\usepackage{xcolor}
\usepackage{hyperref}
% Using standard fonts available on any TeX distribution
\usepackage{mathpazo}
\usepackage[scaled=0.9]{helvet}
\usepackage{courier}
\renewcommand{\familydefault}{\rmdefault}

% ── Conditional: gabarito on/off ──
\ifdefined\GABARITO
  \ifnum\GABARITO=1
    \newcommand{\sol}[1]{\textcolor{blue}{#1}}
    \newcommand{\errocomum}[1]{\textcolor{orange}{⚠️ Erro comum: #1}}
    \newcommand{\competencia}[1]{\textcolor{gray}{Competência: #1}}
  \else
    \newcommand{\sol}[1]{}
    \newcommand{\errocomum}[1]{}
    \newcommand{\competencia}[1]{}
  \fi
\else
  \newcommand{\sol}[1]{}
  \newcommand{\errocomum}[1]{}
  \newcommand{\competencia}[1]{}
\fi

% ── Question formatting ──
\newcommand{\nivel}[1]{\textcolor{gray}{[#1]}}
\newenvironment{questao}[2]{
  \vspace{0.5cm}
  \noindent\textbf{Q#1} \nivel{#2} \quad
}{
  \vspace{0.2cm}
}

\begin{document}

% ════════════════ HEADER ════════════════
\begin{center}
  {\sffamily\Large GBC063 · Semana 1 · Inteligência Artificial}

  \vspace{0.2cm}
  {\small Prof. Marcelo Keese Albertini · FACOM · UFU}

  \vspace{0.2cm}
  {\footnotesize Cobre as competências do resumo S01 ·
  \href{../semana01/material_introducao_ia.html}{material de estudo} ·
  \href{../semana01/resumo_s01_introducao_ia.md}{resumo}}

  \ifdefined\GABARITO
    \ifnum\GABARITO=1
      \vspace{0.3cm}
      \fbox{\textcolor{red}{\textbf{GABARITO — Documento Exclusivo do Professor}}}
    \fi
  \fi
\end{center}

\vspace{0.3cm}
\hrule
\vspace{0.5cm}

% ════════════════ NÍVEL ● (REATIVAÇÃO) ════════════════
\section*{Nível ● \quad Reativação}

\begin{questao}{1}{●}
  Defina \textbf{agente racional} segundo Russell \& Norvig e explique,
  em no máximo 3 linhas, por que esta definição — e não o Teste de Turing —
  é adotada como definição operacional do curso.

  \sol{
    \textbf{Resposta:} Um agente racional é aquele que, para cada percepção possível,
    seleciona a ação que maximiza a métrica de desempenho esperada. O curso adota esta
    definição porque ela (1) funciona com qualquer métrica de sucesso, (2) não depende
    de replicar comportamento humano, e (3) unifica desde um termostato até o AlphaGo.

    \textbf{Resolução:} Russell \& Norvig (2013), Cap. 2. A definição de agente racional
    é independente de humanidade — um agente pode ser racional e não-humano (ex.: AlphaGo
    joga Go de forma sobre-humana).

    \errocomum{Definir agente racional como "agente que age como um humano". Isso é a visão
    "agir humanamente" (Teste de Turing), não "agir racionalmente".}

    \competencia{Explicar a diferença entre IA simbólica e aprendizado de máquina —
    e definir agente racional.}
  }
\end{questao}

\begin{questao}{2}{●}
  As quatro visões clássicas da IA organizam-se em uma matriz 2×2:
  pensamento \textit{vs.} ação × desempenho humano \textit{vs.} racionalidade.
  Classifique cada um dos sistemas abaixo em \textbf{um} dos quatro quadrantes
  e justifique em 1 linha:

  \begin{enumerate}[label=(\alph*),nosep]
    \item ELIZA (Weizenbaum, 1966)
    \item Deep Blue (IBM, 1997)
    \item O agente aspirador de pó da Aula B
    \item Um modelo de linguagem treinado para prever o próximo token
  \end{enumerate}

  \sol{
    \textbf{Respostas:}
    (a) Agir humanamente — ELIZA simulava um terapeuta humano sem raciocinar.
    (b) Agir racionalmente — Deep Blue selecionava lances que maximizavam chance de vitória,
    sem imitar o pensamento de um grande mestre humano.
    (c) Agir racionalmente — o aspirador maximiza a limpeza dada sua percepção.
    (d) Pensar racionalmente (durante o treino: otimiza uma função de perda) e
    agir humanamente (na inferência: gera texto indistinguível do humano).

    \errocomum{Classificar ELIZA como "pensar humanamente". ELIZA não modelava pensamento —
    apenas casava padrões de string.}

    \competencia{Classificar qualquer sistema de IA nas quatro visões e justificar.}
  }
\end{questao}

\begin{questao}{3}{●}
  Verdadeiro ou falso. \textbf{Justifique cada resposta em até 2 linhas.}

  \begin{enumerate}[label=(\alph*),nosep]
    \item "O primeiro inverno da IA (1969–1980) foi causado exclusivamente por falta
    de poder computacional."
    \item "Redes neurais artificiais são uma ideia dos anos 2010."
  \end{enumerate}

  \sol{
    \textbf{Respostas:}
    (a) \textbf{Falso.} O inverno foi causado principalmente pelo livro \textit{Perceptrons}
    (Minsky \& Papert, 1969), que provou limitações do perceptron simples e levou o campo
    a concluir — erroneamente — que redes neurais eram um beco sem saída. A falta de
    computação contribuiu, mas o fator decisivo foi a perda de credibilidade e financiamento.

    (b) \textbf{Falso.} O primeiro modelo de neurônio artificial é de McCulloch \& Pitts (1943).
    O perceptron é de Rosenblatt (1959). O backpropagation que as treina é de 1986
    (Rumelhart, Hinton \& Williams). As ideias fundamentais têm 40–80 anos; o que mudou
    em 2012 foi a escala (dados + GPUs).

    \errocomum{Para (a): atribuir o inverno só à falta de hardware. O fator principal foi
    sociológico — a comunidade perdeu confiança nas redes neurais.}
    \errocomum{Para (b): achar que deep learning é uma invenção recente. As ideias são antigas;
    o que é novo é a escala.}

    \competencia{Enunciar a Bitter Lesson e dar exemplos históricos.}
  }
\end{questao}

% ════════════════ NÍVEL ●● (EXECUÇÃO) ════════════════
\section*{Nível ●● \quad Execução}

\begin{questao}{4}{●●}
  Considere o seguinte ambiente de aspirador de pó com \textbf{três} cômodos
  em linha (A–B–C). O agente começa na posição B. O cômodo A está sujo; B e C
  estão limpos. O agente segue exatamente a lógica abaixo:

  \begin{verbatim}
  def agente_aspirador(percepcao):
      posicao, sujo = percepcao
      if sujo:        return "ASPIRAR"
      if posicao == "A": return "DIREITA"
      if posicao == "B": return "DIREITA"
      return "ESQUERDA"
  \end{verbatim}

  Preencha a tabela de execução para os \textbf{6 primeiros passos}, indicando:
  posição atual, sensor (sujo/limpo), ação tomada, e estado dos cômodos após a ação.
  O agente limpa todos os cômodos? Em quantos passos?

  \sol{
    \textbf{Resposta:}

    \begin{tabular}{c|c|c|c|c}
      Passo & Posição & Sensor & Ação & Cômodos [A,B,C] \\
      \hline
      0 & B & limpo & DIREITA & [sujo, limpo, limpo] \\
      1 & C & limpo & ESQUERDA & [sujo, limpo, limpo] \\
      2 & B & limpo & DIREITA & [sujo, limpo, limpo] \\
      3 & C & limpo & ESQUERDA & [sujo, limpo, limpo] \\
      \dots & \dots & \dots & \dots & \dots \\
    \end{tabular}

    O agente \textbf{nunca alcança o cômodo A} porque sua lógica só vai de A→B→C→B→C→B→C\ldots
    Ele fica preso no ciclo B↔C. Para alcançar A, seria preciso uma regra adicional
    (ex.: "se estou em C e ambos A e B estão sujos, vá para B").

    \textbf{Resolução:} O bug está na lógica de movimento: quando a posição é B ou C,
    o agente nunca se move para A. A condição \texttt{posicao == "A"} nunca é verdadeira
    se o agente não consegue chegar a A.

    \errocomum{Assumir que o agente "dá a volta". O ambiente é uma linha, não um ciclo —
    de C só se vai para B.}

    \competencia{Implementar o agente aspirador e prever seu comportamento
    com 3 cômodos.}
  }
\end{questao}

\begin{questao}{5}{●●}
  O motor de encadeamento progressivo abaixo recebe uma base de fatos inicial
  e um conjunto de regras. Execute o algoritmo \textbf{passo a passo} com os
  dados fornecidos e indique o conjunto final de fatos.

  \begin{verbatim}
  fatos = {"chovendo", "temperatura < 15"}
  regras = [
    ({"chovendo", "temperatura < 15"}, "risco de gelo"),
    ({"risco de gelo"},              "fechar pista"),
    ({"chovendo"},                   "pista molhada"),
  ]
  \end{verbatim}

  \sol{
    \textbf{Resposta:} fatos finais = \{\texttt{"chovendo"}, \texttt{"temperatura < 15"},
    \texttt{"risco de gelo"}, \texttt{"fechar pista"}, \texttt{"pista molhada"}\}

    \textbf{Resolução:}
    \begin{enumerate}[nosep]
      \item Iteração 1: chovendo + temperatura < 15 → adiciona \texttt{"risco de gelo"}.
      \item Iteração 2: risco de gelo → adiciona \texttt{"fechar pista"}.
      \item Iteração 3: chovendo → adiciona \texttt{"pista molhada"}.
      \item Iteração 4: nenhuma regra nova dispara → fim.
    \end{enumerate}

    \errocomum{Esquecer que as regras disparam em qualquer ordem — o algoritmo testa TODAS
    as regras a cada iteração, não para na primeira que dispara.}

    \competencia{Explicar o funcionamento de um motor de inferência e por que regras
    não escalam.}
  }
\end{questao}

\begin{questao}{6}{●●}
  Em 2023, um LLM foi capaz de passar no exame da ordem dos advogados (Bar Exam)
  dos EUA sem ter recebido uma única regra explícita sobre direito. Usando os
  conceitos da Semana 1, explique \textbf{por que isso é um fato tão significativo}
  na história da IA. Se não obteve uma resposta completa, use o roteiro:
  (1) como um sistema especialista resolveria o Bar Exam? (2) por que essa abordagem
  falharia? (3) o que o LLM fez de diferente?

  \sol{
    \textbf{Resposta:} Um sistema especialista exigiria que engenheiros de conhecimento
    codificassem milhares de regras jurídicas à mão — um gargalo intransponível. O LLM
    não recebeu regra nenhuma: aprendeu padrões estatísticos de centenas de gigabytes
    de texto jurídico e geral. Isso é a Bitter Lesson em ação: conhecimento aprendido
    de dados venceu conhecimento artesanal — em um domínio (direito) que se supunha
    exigir raciocínio simbólico explícito.

    \errocomum{Dizer que o LLM "decorou" as respostas. LLMs não têm memória explícita do
    corpus de treino — eles aprendem padrões estatísticos, não fazem lookup de frases.}

    \competencia{Enunciar a Bitter Lesson e dar exemplos que a corroboram.}
  }
\end{questao}

% ════════════════ NÍVEL ●●● (ANÁLISE) ════════════════
\section*{Nível ●●● \quad Análise}

\begin{questao}{7}{●●●}
  Considere a afirmação de Sutton (2019):

  \begin{quote}
  \textit{``Métodos gerais que alavancam computação são, em última instância,
  os mais efetivos — e por uma grande margem. A maioria do esforço em IA é
  gasto tentando construir conhecimento em agentes; isso ajuda no curto prazo,
  mas no longo prazo atinge um platô.''}
  \end{quote}

  \begin{enumerate}[label=(\alph*),nosep]
    \item Por que conhecimento artesanal \textit{necessariamente} atinge um platô,
    enquanto métodos que aprendem de dados não? (Dica: pense no que acontece quando
    você dobra a quantidade de dados ou de computação em cada caso.)
    \item Dê um exemplo \textbf{concreto} de um sistema que atingiu esse platô e
    explique \textit{o que especificamente} o limitou.
  \end{enumerate}

  \sol{
    \textbf{Respostas:}
    (a) Conhecimento artesanal atinge um platô porque sua melhoria depende de \textbf{humanos
    adicionando mais regras} — um processo que é linear, caro e sujeito a contradições.
    Métodos que aprendem melhoram quando você adiciona mais dados ou computação, e esse
    processo é (em princípio) ilimitado: mais dados → melhores gradientes → melhor modelo.
    Não há um humano no loop ditando o teto.

    (b) Sistemas especialistas médicos (ex.: MYCIN, anos 80). Eram melhores que médicos
    juniores em diagnósticos específicos, mas cada nova doença exigia novas regras escritas
    por especialistas. Quando o número de regras passou de centenas, as interações entre
    elas tornaram-se imprevisíveis. O sistema não melhorava com mais dados de pacientes —
    só com mais regras escritas à mão. Enquanto isso, métodos de ML atuais melhoram
    automaticamente com mais exames, prontuários e artigos.

    \errocomum{Em (a), dizer que conhecimento artesanal atinge platô "porque humanos
    são limitados". A resposta precisa identificar o mecanismo: a taxa de melhoria é
    ditada pela velocidade com que humanos escrevem regras, não pela velocidade com
    que computação cresce.}

    \competencia{Enunciar a Bitter Lesson e explicar o mecanismo por trás dela.}
  }
\end{questao}

\begin{questao}{8}{●●●}
  O efeito ELIZA — atribuir inteligência onde há apenas casamento de padrões —
  foi documentado em 1966. Ele persiste em 2025 com LLMs?

  \begin{enumerate}[label=(\alph*),nosep]
    \item Aponte \textbf{uma semelhança} entre ELIZA e um LLM moderno que torna
    o efeito ELIZA relevante hoje.
    \item Aponte \textbf{uma diferença fundamental} que torna a comparação enganosa.
    \item Em no máximo 4 linhas: um LLM é ``só casamento de padrões''? Justifique.
  \end{enumerate}

  \sol{
    \textbf{Respostas:}
    (a) Ambos produzem texto fluente sem garantia de compreensão do conteúdo semântico.
    Um LLM pode gerar uma explicação convincente de um conceito que ele não ``entende''
    no sentido humano — assim como ELIZA simulava empatia sem sentir nada. O efeito
    ELIZA ocorre quando o usuário projeta compreensão onde há apenas correlação estatística.

    (b) ELIZA usava algumas dezenas de regras de reescrita escritas à mão. Um LLM aprende
    de centenas de bilhões de parâmetros ajustados sobre trilhões de tokens. A escala e
    a origem do conhecimento (artesanal vs. aprendido) são de natureza completamente diferente.

    (c) Não. Reduzir um LLM a ``só casamento de padrões'' é como reduzir um cérebro a
    ``só sinapses disparando'' — tecnicamente verdadeiro, mas enganoso. A questão relevante
    não é se o mecanismo básico é simples (é: multiplicação de matrizes + não-linearidades),
    mas se a escala produz comportamentos emergentes que merecem ser chamados de raciocínio.
    Esta é uma pergunta em aberto na pesquisa em 2025.

    \errocomum{Responder (c) com um simples ``sim'' ou ``não''. A pergunta exige distinguir
    mecanismo de comportamento — o que o sistema \textit{faz} vs. como ele \textit{faz}.}

    \competencia{Classificar sistemas de IA e analisar criticamente alegações de inteligência.}
  }
\end{questao}

% ════════════════ NÍVEL ★ (SÍNTESE) ════════════════
\section*{Nível ★ \quad Síntese}

\begin{questao}{9}{★}
  Projete um experimento mental: você herda um sistema especialista médico com
  10.000 regras, desenvolvido ao longo de 15 anos por uma equipe de 20 engenheiros
  de conhecimento. O sistema tem 92\% de acurácia em diagnóstico.

  \begin{enumerate}[label=(\alph*),nosep]
    \item Identifique \textbf{três} cenários em que este sistema falharia de forma
    sistemática que um modelo de ML treinado sobre os mesmos dados não falharia.
    \item Para cada cenário, explique \textbf{por que} a natureza artesanal das
    regras é a causa raiz da falha.
    \item Proponha uma arquitetura híbrida (simbólico + aprendizado) que combine
    o melhor dos dois mundos para este domínio. Descreva em no máximo 6 linhas.
  \end{enumerate}

  \sol{
    \textbf{Respostas:}
    (a) e (b):
    \begin{enumerate}[nosep]
      \item \textbf{Nova doença:} o SE falha porque nenhuma regra cobre o fenômeno novo.
      O ML pode detectar que os sintomas formam um cluster distinto nos dados e sinalizar
      uma anomalia — mesmo sem rótulos. Causa raiz: SE só reconhece o que foi explicitamente
      programado.
      \item \textbf{Interação entre condições:} um paciente com diabetes \textit{e} COVID
      pode ativar regras contraditórias (uma recomenda corticoides, outra contraindica).
      O SE não tem como ponderar evidências conflitantes. O ML aprende pesos relativos
      dos features. Causa raiz: regras são discretas e não-compostas; não modelam interações.
      \item \textbf{Deriva de distribuição:} se a população de pacientes mudar (ex.:
      envelhecimento), a acurácia do SE cai porque as regras foram escritas para a
      distribuição antiga. O ML pode ser retreinado com dados novos. Causa raiz: SE
      não se adapta sem reengenharia manual.
    \end{enumerate}

    (c) \textbf{Arquitetura híbrida:} usar um modelo de ML (ex.: gradient-boosted trees
    ou rede neural) para produzir o diagnóstico base e a estimativa de confiança. As
    10.000 regras viram um sistema de \textbf{verificação}: antes de emitir o diagnóstico,
    o sistema checa se ele viola alguma regra de segurança crítica (ex.: ``nunca prescrever
    droga X para grávidas''). Se violar, escala para um médico humano. O ML lida com a
    generalização; as regras lidam com as restrições rígidas de segurança.

    \errocomum{Propor um híbrido em que as regras \textit{guiam} o aprendizado. Isso
    reintroduz o gargalo: as regras limitam o que o modelo pode aprender. O melhor
    híbrido usa regras como \textit{restrições de segurança}, não como indutores de viés.}

    \competencia{Sintetizar os conceitos de IA simbólica e aprendizado em uma arquitetura
    fundamentada.}
  }
\end{questao}

\begin{questao}{10}{★ 🌉}
  O treinamento do GPT usa a mesma estrutura de loop que o agente aspirador da Aula B:
  \textbf{percepção → decisão → ação}. A diferença está na complexidade de \texttt{decidir()}.

  \begin{enumerate}[label=(\alph*),nosep]
    \item No aspirador, \texttt{decidir()} é uma função com 4 regras \texttt{if}.
    No GPT, \texttt{decidir()} é uma rede neural com centenas de bilhões de
    parâmetros. \textbf{Por que}, conceitualmente, a mesma interface de 4 linhas
    funciona para ambos? O que isso revela sobre o paradigma de agentes?
    \item A Bitter Lesson argumenta que conhecimento artesanal atinge um platô.
    O prompt engineering — escrever instruções cuidadosas para um LLM — é uma
    forma de conhecimento artesanal? Justifique em até 5 linhas.
  \end{enumerate}

  \sol{
    \textbf{Respostas:}
    (a) A interface \texttt{obs → decidir(obs) → ação} funciona para ambos porque ela
    abstrai \textbf{o que} o agente faz de \textbf{como} ele decide. O paradigma de
    agentes é uma \textit{casca}: define a assinatura da decisão, não sua implementação.
    O que muda entre o aspirador e o GPT é apenas a função \texttt{decidir()} — e essa
    função pode ser desde 4 ifs até uma rede de 175 bilhões de parâmetros. Isso revela
    que o paradigma de agentes é \textbf{universal}: ele unifica toda a IA sob uma
    única abstração.

    (b) Sim, prompt engineering é conhecimento artesanal — mas de um tipo novo. É
    artesanal porque um humano está manualmente ajustando a entrada para melhorar
    a saída, sem que o modelo aprenda com esse ajuste. Mas difere do artesanal
    clássico porque (1) o modelo subjacente foi aprendido de dados, não programado,
    e (2) o custo de ajustar um prompt é ordens de grandeza menor que o de escrever
    regras. Dito isso, a Bitter Lesson sugere que, no longo prazo, métodos que
    \textit{aprendem} a fazer prompting (auto-prompting, RLHF) substituirão o
    prompting manual — e já há evidências disso.

    \errocomum{Em (b), responder simplesmente ``não'' porque o modelo foi treinado com ML.
    O prompt em si é conhecimento artesanal injetado \textit{na inferência}, não no treino.}

    \competencia{Reconhecer o loop percepção→decisão→ação como a interface universal
    da IA e aplicar a Bitter Lesson a desenvolvimentos recentes.}
  }
\end{questao}

\vspace{1cm}
\hrule
\vspace{0.3cm}
\begin{center}
  {\footnotesize Lista S01 · GBC063 · Prof. Marcelo Keese Albertini · FACOM/UFU}
  \ifdefined\GABARITO
    \ifnum\GABARITO=1
      \\[0.2cm] {\footnotesize\color{red}\textbf{GABARITO — USO EXCLUSIVO DO PROFESSOR}}
    \fi
  \fi
\end{center}

\end{document}
"
%   Gabarito:  pdflatex -jobname=gabarito_lista_s01_introducao_ia "\def\GABARITO{1}% ═══════════════════════════════════════════════════════════════
% Lista S01 — O que é Inteligência Artificial?
% Prof. Marcelo Keese Albertini · FACOM · UFU · GBC063
%
% Compilação:
%   Questões:  pdflatex -jobname=lista_s01_introducao_ia "\def\GABARITO{0}\input{lista_s01_introducao_ia.tex}"
%   Gabarito:  pdflatex -jobname=gabarito_lista_s01_introducao_ia "\def\GABARITO{1}\input{lista_s01_introducao_ia.tex}"
% ═══════════════════════════════════════════════════════════════

\documentclass[12pt,a4paper]{article}

\usepackage[utf8]{inputenc}
\usepackage[T1]{fontenc}
\usepackage[brazil]{babel}
\usepackage{geometry}
\geometry{margin=2.2cm}
\usepackage{enumitem}
\usepackage{amsmath,amssymb}
\usepackage{xcolor}
\usepackage{hyperref}
% Using standard fonts available on any TeX distribution
\usepackage{mathpazo}
\usepackage[scaled=0.9]{helvet}
\usepackage{courier}
\renewcommand{\familydefault}{\rmdefault}

% ── Conditional: gabarito on/off ──
\ifdefined\GABARITO
  \ifnum\GABARITO=1
    \newcommand{\sol}[1]{\textcolor{blue}{#1}}
    \newcommand{\errocomum}[1]{\textcolor{orange}{⚠️ Erro comum: #1}}
    \newcommand{\competencia}[1]{\textcolor{gray}{Competência: #1}}
  \else
    \newcommand{\sol}[1]{}
    \newcommand{\errocomum}[1]{}
    \newcommand{\competencia}[1]{}
  \fi
\else
  \newcommand{\sol}[1]{}
  \newcommand{\errocomum}[1]{}
  \newcommand{\competencia}[1]{}
\fi

% ── Question formatting ──
\newcommand{\nivel}[1]{\textcolor{gray}{[#1]}}
\newenvironment{questao}[2]{
  \vspace{0.5cm}
  \noindent\textbf{Q#1} \nivel{#2} \quad
}{
  \vspace{0.2cm}
}

\begin{document}

% ════════════════ HEADER ════════════════
\begin{center}
  {\sffamily\Large GBC063 · Semana 1 · Inteligência Artificial}

  \vspace{0.2cm}
  {\small Prof. Marcelo Keese Albertini · FACOM · UFU}

  \vspace{0.2cm}
  {\footnotesize Cobre as competências do resumo S01 ·
  \href{../semana01/material_introducao_ia.html}{material de estudo} ·
  \href{../semana01/resumo_s01_introducao_ia.md}{resumo}}

  \ifdefined\GABARITO
    \ifnum\GABARITO=1
      \vspace{0.3cm}
      \fbox{\textcolor{red}{\textbf{GABARITO — Documento Exclusivo do Professor}}}
    \fi
  \fi
\end{center}

\vspace{0.3cm}
\hrule
\vspace{0.5cm}

% ════════════════ NÍVEL ● (REATIVAÇÃO) ════════════════
\section*{Nível ● \quad Reativação}

\begin{questao}{1}{●}
  Defina \textbf{agente racional} segundo Russell \& Norvig e explique,
  em no máximo 3 linhas, por que esta definição — e não o Teste de Turing —
  é adotada como definição operacional do curso.

  \sol{
    \textbf{Resposta:} Um agente racional é aquele que, para cada percepção possível,
    seleciona a ação que maximiza a métrica de desempenho esperada. O curso adota esta
    definição porque ela (1) funciona com qualquer métrica de sucesso, (2) não depende
    de replicar comportamento humano, e (3) unifica desde um termostato até o AlphaGo.

    \textbf{Resolução:} Russell \& Norvig (2013), Cap. 2. A definição de agente racional
    é independente de humanidade — um agente pode ser racional e não-humano (ex.: AlphaGo
    joga Go de forma sobre-humana).

    \errocomum{Definir agente racional como "agente que age como um humano". Isso é a visão
    "agir humanamente" (Teste de Turing), não "agir racionalmente".}

    \competencia{Explicar a diferença entre IA simbólica e aprendizado de máquina —
    e definir agente racional.}
  }
\end{questao}

\begin{questao}{2}{●}
  As quatro visões clássicas da IA organizam-se em uma matriz 2×2:
  pensamento \textit{vs.} ação × desempenho humano \textit{vs.} racionalidade.
  Classifique cada um dos sistemas abaixo em \textbf{um} dos quatro quadrantes
  e justifique em 1 linha:

  \begin{enumerate}[label=(\alph*),nosep]
    \item ELIZA (Weizenbaum, 1966)
    \item Deep Blue (IBM, 1997)
    \item O agente aspirador de pó da Aula B
    \item Um modelo de linguagem treinado para prever o próximo token
  \end{enumerate}

  \sol{
    \textbf{Respostas:}
    (a) Agir humanamente — ELIZA simulava um terapeuta humano sem raciocinar.
    (b) Agir racionalmente — Deep Blue selecionava lances que maximizavam chance de vitória,
    sem imitar o pensamento de um grande mestre humano.
    (c) Agir racionalmente — o aspirador maximiza a limpeza dada sua percepção.
    (d) Pensar racionalmente (durante o treino: otimiza uma função de perda) e
    agir humanamente (na inferência: gera texto indistinguível do humano).

    \errocomum{Classificar ELIZA como "pensar humanamente". ELIZA não modelava pensamento —
    apenas casava padrões de string.}

    \competencia{Classificar qualquer sistema de IA nas quatro visões e justificar.}
  }
\end{questao}

\begin{questao}{3}{●}
  Verdadeiro ou falso. \textbf{Justifique cada resposta em até 2 linhas.}

  \begin{enumerate}[label=(\alph*),nosep]
    \item "O primeiro inverno da IA (1969–1980) foi causado exclusivamente por falta
    de poder computacional."
    \item "Redes neurais artificiais são uma ideia dos anos 2010."
  \end{enumerate}

  \sol{
    \textbf{Respostas:}
    (a) \textbf{Falso.} O inverno foi causado principalmente pelo livro \textit{Perceptrons}
    (Minsky \& Papert, 1969), que provou limitações do perceptron simples e levou o campo
    a concluir — erroneamente — que redes neurais eram um beco sem saída. A falta de
    computação contribuiu, mas o fator decisivo foi a perda de credibilidade e financiamento.

    (b) \textbf{Falso.} O primeiro modelo de neurônio artificial é de McCulloch \& Pitts (1943).
    O perceptron é de Rosenblatt (1959). O backpropagation que as treina é de 1986
    (Rumelhart, Hinton \& Williams). As ideias fundamentais têm 40–80 anos; o que mudou
    em 2012 foi a escala (dados + GPUs).

    \errocomum{Para (a): atribuir o inverno só à falta de hardware. O fator principal foi
    sociológico — a comunidade perdeu confiança nas redes neurais.}
    \errocomum{Para (b): achar que deep learning é uma invenção recente. As ideias são antigas;
    o que é novo é a escala.}

    \competencia{Enunciar a Bitter Lesson e dar exemplos históricos.}
  }
\end{questao}

% ════════════════ NÍVEL ●● (EXECUÇÃO) ════════════════
\section*{Nível ●● \quad Execução}

\begin{questao}{4}{●●}
  Considere o seguinte ambiente de aspirador de pó com \textbf{três} cômodos
  em linha (A–B–C). O agente começa na posição B. O cômodo A está sujo; B e C
  estão limpos. O agente segue exatamente a lógica abaixo:

  \begin{verbatim}
  def agente_aspirador(percepcao):
      posicao, sujo = percepcao
      if sujo:        return "ASPIRAR"
      if posicao == "A": return "DIREITA"
      if posicao == "B": return "DIREITA"
      return "ESQUERDA"
  \end{verbatim}

  Preencha a tabela de execução para os \textbf{6 primeiros passos}, indicando:
  posição atual, sensor (sujo/limpo), ação tomada, e estado dos cômodos após a ação.
  O agente limpa todos os cômodos? Em quantos passos?

  \sol{
    \textbf{Resposta:}

    \begin{tabular}{c|c|c|c|c}
      Passo & Posição & Sensor & Ação & Cômodos [A,B,C] \\
      \hline
      0 & B & limpo & DIREITA & [sujo, limpo, limpo] \\
      1 & C & limpo & ESQUERDA & [sujo, limpo, limpo] \\
      2 & B & limpo & DIREITA & [sujo, limpo, limpo] \\
      3 & C & limpo & ESQUERDA & [sujo, limpo, limpo] \\
      \dots & \dots & \dots & \dots & \dots \\
    \end{tabular}

    O agente \textbf{nunca alcança o cômodo A} porque sua lógica só vai de A→B→C→B→C→B→C\ldots
    Ele fica preso no ciclo B↔C. Para alcançar A, seria preciso uma regra adicional
    (ex.: "se estou em C e ambos A e B estão sujos, vá para B").

    \textbf{Resolução:} O bug está na lógica de movimento: quando a posição é B ou C,
    o agente nunca se move para A. A condição \texttt{posicao == "A"} nunca é verdadeira
    se o agente não consegue chegar a A.

    \errocomum{Assumir que o agente "dá a volta". O ambiente é uma linha, não um ciclo —
    de C só se vai para B.}

    \competencia{Implementar o agente aspirador e prever seu comportamento
    com 3 cômodos.}
  }
\end{questao}

\begin{questao}{5}{●●}
  O motor de encadeamento progressivo abaixo recebe uma base de fatos inicial
  e um conjunto de regras. Execute o algoritmo \textbf{passo a passo} com os
  dados fornecidos e indique o conjunto final de fatos.

  \begin{verbatim}
  fatos = {"chovendo", "temperatura < 15"}
  regras = [
    ({"chovendo", "temperatura < 15"}, "risco de gelo"),
    ({"risco de gelo"},              "fechar pista"),
    ({"chovendo"},                   "pista molhada"),
  ]
  \end{verbatim}

  \sol{
    \textbf{Resposta:} fatos finais = \{\texttt{"chovendo"}, \texttt{"temperatura < 15"},
    \texttt{"risco de gelo"}, \texttt{"fechar pista"}, \texttt{"pista molhada"}\}

    \textbf{Resolução:}
    \begin{enumerate}[nosep]
      \item Iteração 1: chovendo + temperatura < 15 → adiciona \texttt{"risco de gelo"}.
      \item Iteração 2: risco de gelo → adiciona \texttt{"fechar pista"}.
      \item Iteração 3: chovendo → adiciona \texttt{"pista molhada"}.
      \item Iteração 4: nenhuma regra nova dispara → fim.
    \end{enumerate}

    \errocomum{Esquecer que as regras disparam em qualquer ordem — o algoritmo testa TODAS
    as regras a cada iteração, não para na primeira que dispara.}

    \competencia{Explicar o funcionamento de um motor de inferência e por que regras
    não escalam.}
  }
\end{questao}

\begin{questao}{6}{●●}
  Em 2023, um LLM foi capaz de passar no exame da ordem dos advogados (Bar Exam)
  dos EUA sem ter recebido uma única regra explícita sobre direito. Usando os
  conceitos da Semana 1, explique \textbf{por que isso é um fato tão significativo}
  na história da IA. Se não obteve uma resposta completa, use o roteiro:
  (1) como um sistema especialista resolveria o Bar Exam? (2) por que essa abordagem
  falharia? (3) o que o LLM fez de diferente?

  \sol{
    \textbf{Resposta:} Um sistema especialista exigiria que engenheiros de conhecimento
    codificassem milhares de regras jurídicas à mão — um gargalo intransponível. O LLM
    não recebeu regra nenhuma: aprendeu padrões estatísticos de centenas de gigabytes
    de texto jurídico e geral. Isso é a Bitter Lesson em ação: conhecimento aprendido
    de dados venceu conhecimento artesanal — em um domínio (direito) que se supunha
    exigir raciocínio simbólico explícito.

    \errocomum{Dizer que o LLM "decorou" as respostas. LLMs não têm memória explícita do
    corpus de treino — eles aprendem padrões estatísticos, não fazem lookup de frases.}

    \competencia{Enunciar a Bitter Lesson e dar exemplos que a corroboram.}
  }
\end{questao}

% ════════════════ NÍVEL ●●● (ANÁLISE) ════════════════
\section*{Nível ●●● \quad Análise}

\begin{questao}{7}{●●●}
  Considere a afirmação de Sutton (2019):

  \begin{quote}
  \textit{``Métodos gerais que alavancam computação são, em última instância,
  os mais efetivos — e por uma grande margem. A maioria do esforço em IA é
  gasto tentando construir conhecimento em agentes; isso ajuda no curto prazo,
  mas no longo prazo atinge um platô.''}
  \end{quote}

  \begin{enumerate}[label=(\alph*),nosep]
    \item Por que conhecimento artesanal \textit{necessariamente} atinge um platô,
    enquanto métodos que aprendem de dados não? (Dica: pense no que acontece quando
    você dobra a quantidade de dados ou de computação em cada caso.)
    \item Dê um exemplo \textbf{concreto} de um sistema que atingiu esse platô e
    explique \textit{o que especificamente} o limitou.
  \end{enumerate}

  \sol{
    \textbf{Respostas:}
    (a) Conhecimento artesanal atinge um platô porque sua melhoria depende de \textbf{humanos
    adicionando mais regras} — um processo que é linear, caro e sujeito a contradições.
    Métodos que aprendem melhoram quando você adiciona mais dados ou computação, e esse
    processo é (em princípio) ilimitado: mais dados → melhores gradientes → melhor modelo.
    Não há um humano no loop ditando o teto.

    (b) Sistemas especialistas médicos (ex.: MYCIN, anos 80). Eram melhores que médicos
    juniores em diagnósticos específicos, mas cada nova doença exigia novas regras escritas
    por especialistas. Quando o número de regras passou de centenas, as interações entre
    elas tornaram-se imprevisíveis. O sistema não melhorava com mais dados de pacientes —
    só com mais regras escritas à mão. Enquanto isso, métodos de ML atuais melhoram
    automaticamente com mais exames, prontuários e artigos.

    \errocomum{Em (a), dizer que conhecimento artesanal atinge platô "porque humanos
    são limitados". A resposta precisa identificar o mecanismo: a taxa de melhoria é
    ditada pela velocidade com que humanos escrevem regras, não pela velocidade com
    que computação cresce.}

    \competencia{Enunciar a Bitter Lesson e explicar o mecanismo por trás dela.}
  }
\end{questao}

\begin{questao}{8}{●●●}
  O efeito ELIZA — atribuir inteligência onde há apenas casamento de padrões —
  foi documentado em 1966. Ele persiste em 2025 com LLMs?

  \begin{enumerate}[label=(\alph*),nosep]
    \item Aponte \textbf{uma semelhança} entre ELIZA e um LLM moderno que torna
    o efeito ELIZA relevante hoje.
    \item Aponte \textbf{uma diferença fundamental} que torna a comparação enganosa.
    \item Em no máximo 4 linhas: um LLM é ``só casamento de padrões''? Justifique.
  \end{enumerate}

  \sol{
    \textbf{Respostas:}
    (a) Ambos produzem texto fluente sem garantia de compreensão do conteúdo semântico.
    Um LLM pode gerar uma explicação convincente de um conceito que ele não ``entende''
    no sentido humano — assim como ELIZA simulava empatia sem sentir nada. O efeito
    ELIZA ocorre quando o usuário projeta compreensão onde há apenas correlação estatística.

    (b) ELIZA usava algumas dezenas de regras de reescrita escritas à mão. Um LLM aprende
    de centenas de bilhões de parâmetros ajustados sobre trilhões de tokens. A escala e
    a origem do conhecimento (artesanal vs. aprendido) são de natureza completamente diferente.

    (c) Não. Reduzir um LLM a ``só casamento de padrões'' é como reduzir um cérebro a
    ``só sinapses disparando'' — tecnicamente verdadeiro, mas enganoso. A questão relevante
    não é se o mecanismo básico é simples (é: multiplicação de matrizes + não-linearidades),
    mas se a escala produz comportamentos emergentes que merecem ser chamados de raciocínio.
    Esta é uma pergunta em aberto na pesquisa em 2025.

    \errocomum{Responder (c) com um simples ``sim'' ou ``não''. A pergunta exige distinguir
    mecanismo de comportamento — o que o sistema \textit{faz} vs. como ele \textit{faz}.}

    \competencia{Classificar sistemas de IA e analisar criticamente alegações de inteligência.}
  }
\end{questao}

% ════════════════ NÍVEL ★ (SÍNTESE) ════════════════
\section*{Nível ★ \quad Síntese}

\begin{questao}{9}{★}
  Projete um experimento mental: você herda um sistema especialista médico com
  10.000 regras, desenvolvido ao longo de 15 anos por uma equipe de 20 engenheiros
  de conhecimento. O sistema tem 92\% de acurácia em diagnóstico.

  \begin{enumerate}[label=(\alph*),nosep]
    \item Identifique \textbf{três} cenários em que este sistema falharia de forma
    sistemática que um modelo de ML treinado sobre os mesmos dados não falharia.
    \item Para cada cenário, explique \textbf{por que} a natureza artesanal das
    regras é a causa raiz da falha.
    \item Proponha uma arquitetura híbrida (simbólico + aprendizado) que combine
    o melhor dos dois mundos para este domínio. Descreva em no máximo 6 linhas.
  \end{enumerate}

  \sol{
    \textbf{Respostas:}
    (a) e (b):
    \begin{enumerate}[nosep]
      \item \textbf{Nova doença:} o SE falha porque nenhuma regra cobre o fenômeno novo.
      O ML pode detectar que os sintomas formam um cluster distinto nos dados e sinalizar
      uma anomalia — mesmo sem rótulos. Causa raiz: SE só reconhece o que foi explicitamente
      programado.
      \item \textbf{Interação entre condições:} um paciente com diabetes \textit{e} COVID
      pode ativar regras contraditórias (uma recomenda corticoides, outra contraindica).
      O SE não tem como ponderar evidências conflitantes. O ML aprende pesos relativos
      dos features. Causa raiz: regras são discretas e não-compostas; não modelam interações.
      \item \textbf{Deriva de distribuição:} se a população de pacientes mudar (ex.:
      envelhecimento), a acurácia do SE cai porque as regras foram escritas para a
      distribuição antiga. O ML pode ser retreinado com dados novos. Causa raiz: SE
      não se adapta sem reengenharia manual.
    \end{enumerate}

    (c) \textbf{Arquitetura híbrida:} usar um modelo de ML (ex.: gradient-boosted trees
    ou rede neural) para produzir o diagnóstico base e a estimativa de confiança. As
    10.000 regras viram um sistema de \textbf{verificação}: antes de emitir o diagnóstico,
    o sistema checa se ele viola alguma regra de segurança crítica (ex.: ``nunca prescrever
    droga X para grávidas''). Se violar, escala para um médico humano. O ML lida com a
    generalização; as regras lidam com as restrições rígidas de segurança.

    \errocomum{Propor um híbrido em que as regras \textit{guiam} o aprendizado. Isso
    reintroduz o gargalo: as regras limitam o que o modelo pode aprender. O melhor
    híbrido usa regras como \textit{restrições de segurança}, não como indutores de viés.}

    \competencia{Sintetizar os conceitos de IA simbólica e aprendizado em uma arquitetura
    fundamentada.}
  }
\end{questao}

\begin{questao}{10}{★ 🌉}
  O treinamento do GPT usa a mesma estrutura de loop que o agente aspirador da Aula B:
  \textbf{percepção → decisão → ação}. A diferença está na complexidade de \texttt{decidir()}.

  \begin{enumerate}[label=(\alph*),nosep]
    \item No aspirador, \texttt{decidir()} é uma função com 4 regras \texttt{if}.
    No GPT, \texttt{decidir()} é uma rede neural com centenas de bilhões de
    parâmetros. \textbf{Por que}, conceitualmente, a mesma interface de 4 linhas
    funciona para ambos? O que isso revela sobre o paradigma de agentes?
    \item A Bitter Lesson argumenta que conhecimento artesanal atinge um platô.
    O prompt engineering — escrever instruções cuidadosas para um LLM — é uma
    forma de conhecimento artesanal? Justifique em até 5 linhas.
  \end{enumerate}

  \sol{
    \textbf{Respostas:}
    (a) A interface \texttt{obs → decidir(obs) → ação} funciona para ambos porque ela
    abstrai \textbf{o que} o agente faz de \textbf{como} ele decide. O paradigma de
    agentes é uma \textit{casca}: define a assinatura da decisão, não sua implementação.
    O que muda entre o aspirador e o GPT é apenas a função \texttt{decidir()} — e essa
    função pode ser desde 4 ifs até uma rede de 175 bilhões de parâmetros. Isso revela
    que o paradigma de agentes é \textbf{universal}: ele unifica toda a IA sob uma
    única abstração.

    (b) Sim, prompt engineering é conhecimento artesanal — mas de um tipo novo. É
    artesanal porque um humano está manualmente ajustando a entrada para melhorar
    a saída, sem que o modelo aprenda com esse ajuste. Mas difere do artesanal
    clássico porque (1) o modelo subjacente foi aprendido de dados, não programado,
    e (2) o custo de ajustar um prompt é ordens de grandeza menor que o de escrever
    regras. Dito isso, a Bitter Lesson sugere que, no longo prazo, métodos que
    \textit{aprendem} a fazer prompting (auto-prompting, RLHF) substituirão o
    prompting manual — e já há evidências disso.

    \errocomum{Em (b), responder simplesmente ``não'' porque o modelo foi treinado com ML.
    O prompt em si é conhecimento artesanal injetado \textit{na inferência}, não no treino.}

    \competencia{Reconhecer o loop percepção→decisão→ação como a interface universal
    da IA e aplicar a Bitter Lesson a desenvolvimentos recentes.}
  }
\end{questao}

\vspace{1cm}
\hrule
\vspace{0.3cm}
\begin{center}
  {\footnotesize Lista S01 · GBC063 · Prof. Marcelo Keese Albertini · FACOM/UFU}
  \ifdefined\GABARITO
    \ifnum\GABARITO=1
      \\[0.2cm] {\footnotesize\color{red}\textbf{GABARITO — USO EXCLUSIVO DO PROFESSOR}}
    \fi
  \fi
\end{center}

\end{document}
"
% ═══════════════════════════════════════════════════════════════

\documentclass[12pt,a4paper]{article}

\usepackage[utf8]{inputenc}
\usepackage[T1]{fontenc}
\usepackage[brazil]{babel}
\usepackage{geometry}
\geometry{margin=2.2cm}
\usepackage{enumitem}
\usepackage{amsmath,amssymb}
\usepackage{xcolor}
\usepackage{hyperref}
% Using standard fonts available on any TeX distribution
\usepackage{mathpazo}
\usepackage[scaled=0.9]{helvet}
\usepackage{courier}
\renewcommand{\familydefault}{\rmdefault}

% ── Conditional: gabarito on/off ──
\ifdefined\GABARITO
  \ifnum\GABARITO=1
    \newcommand{\sol}[1]{\textcolor{blue}{#1}}
    \newcommand{\errocomum}[1]{\textcolor{orange}{⚠️ Erro comum: #1}}
    \newcommand{\competencia}[1]{\textcolor{gray}{Competência: #1}}
  \else
    \newcommand{\sol}[1]{}
    \newcommand{\errocomum}[1]{}
    \newcommand{\competencia}[1]{}
  \fi
\else
  \newcommand{\sol}[1]{}
  \newcommand{\errocomum}[1]{}
  \newcommand{\competencia}[1]{}
\fi

% ── Question formatting ──
\newcommand{\nivel}[1]{\textcolor{gray}{[#1]}}
\newenvironment{questao}[2]{
  \vspace{0.5cm}
  \noindent\textbf{Q#1} \nivel{#2} \quad
}{
  \vspace{0.2cm}
}

\begin{document}

% ════════════════ HEADER ════════════════
\begin{center}
  {\sffamily\Large GBC063 · Semana 1 · Inteligência Artificial}

  \vspace{0.2cm}
  {\small Prof. Marcelo Keese Albertini · FACOM · UFU}

  \vspace{0.2cm}
  {\footnotesize Cobre as competências do resumo S01 ·
  \href{../semana01/material_introducao_ia.html}{material de estudo} ·
  \href{../semana01/resumo_s01_introducao_ia.md}{resumo}}

  \ifdefined\GABARITO
    \ifnum\GABARITO=1
      \vspace{0.3cm}
      \fbox{\textcolor{red}{\textbf{GABARITO — Documento Exclusivo do Professor}}}
    \fi
  \fi
\end{center}

\vspace{0.3cm}
\hrule
\vspace{0.5cm}

% ════════════════ NÍVEL ● (REATIVAÇÃO) ════════════════
\section*{Nível ● \quad Reativação}

\begin{questao}{1}{●}
  Defina \textbf{agente racional} segundo Russell \& Norvig e explique,
  em no máximo 3 linhas, por que esta definição — e não o Teste de Turing —
  é adotada como definição operacional do curso.

  \sol{
    \textbf{Resposta:} Um agente racional é aquele que, para cada percepção possível,
    seleciona a ação que maximiza a métrica de desempenho esperada. O curso adota esta
    definição porque ela (1) funciona com qualquer métrica de sucesso, (2) não depende
    de replicar comportamento humano, e (3) unifica desde um termostato até o AlphaGo.

    \textbf{Resolução:} Russell \& Norvig (2013), Cap. 2. A definição de agente racional
    é independente de humanidade — um agente pode ser racional e não-humano (ex.: AlphaGo
    joga Go de forma sobre-humana).

    \errocomum{Definir agente racional como "agente que age como um humano". Isso é a visão
    "agir humanamente" (Teste de Turing), não "agir racionalmente".}

    \competencia{Explicar a diferença entre IA simbólica e aprendizado de máquina —
    e definir agente racional.}
  }
\end{questao}

\begin{questao}{2}{●}
  As quatro visões clássicas da IA organizam-se em uma matriz 2×2:
  pensamento \textit{vs.} ação × desempenho humano \textit{vs.} racionalidade.
  Classifique cada um dos sistemas abaixo em \textbf{um} dos quatro quadrantes
  e justifique em 1 linha:

  \begin{enumerate}[label=(\alph*),nosep]
    \item ELIZA (Weizenbaum, 1966)
    \item Deep Blue (IBM, 1997)
    \item O agente aspirador de pó da Aula B
    \item Um modelo de linguagem treinado para prever o próximo token
  \end{enumerate}

  \sol{
    \textbf{Respostas:}
    (a) Agir humanamente — ELIZA simulava um terapeuta humano sem raciocinar.
    (b) Agir racionalmente — Deep Blue selecionava lances que maximizavam chance de vitória,
    sem imitar o pensamento de um grande mestre humano.
    (c) Agir racionalmente — o aspirador maximiza a limpeza dada sua percepção.
    (d) Pensar racionalmente (durante o treino: otimiza uma função de perda) e
    agir humanamente (na inferência: gera texto indistinguível do humano).

    \errocomum{Classificar ELIZA como "pensar humanamente". ELIZA não modelava pensamento —
    apenas casava padrões de string.}

    \competencia{Classificar qualquer sistema de IA nas quatro visões e justificar.}
  }
\end{questao}

\begin{questao}{3}{●}
  Verdadeiro ou falso. \textbf{Justifique cada resposta em até 2 linhas.}

  \begin{enumerate}[label=(\alph*),nosep]
    \item "O primeiro inverno da IA (1969–1980) foi causado exclusivamente por falta
    de poder computacional."
    \item "Redes neurais artificiais são uma ideia dos anos 2010."
  \end{enumerate}

  \sol{
    \textbf{Respostas:}
    (a) \textbf{Falso.} O inverno foi causado principalmente pelo livro \textit{Perceptrons}
    (Minsky \& Papert, 1969), que provou limitações do perceptron simples e levou o campo
    a concluir — erroneamente — que redes neurais eram um beco sem saída. A falta de
    computação contribuiu, mas o fator decisivo foi a perda de credibilidade e financiamento.

    (b) \textbf{Falso.} O primeiro modelo de neurônio artificial é de McCulloch \& Pitts (1943).
    O perceptron é de Rosenblatt (1959). O backpropagation que as treina é de 1986
    (Rumelhart, Hinton \& Williams). As ideias fundamentais têm 40–80 anos; o que mudou
    em 2012 foi a escala (dados + GPUs).

    \errocomum{Para (a): atribuir o inverno só à falta de hardware. O fator principal foi
    sociológico — a comunidade perdeu confiança nas redes neurais.}
    \errocomum{Para (b): achar que deep learning é uma invenção recente. As ideias são antigas;
    o que é novo é a escala.}

    \competencia{Enunciar a Bitter Lesson e dar exemplos históricos.}
  }
\end{questao}

% ════════════════ NÍVEL ●● (EXECUÇÃO) ════════════════
\section*{Nível ●● \quad Execução}

\begin{questao}{4}{●●}
  Considere o seguinte ambiente de aspirador de pó com \textbf{três} cômodos
  em linha (A–B–C). O agente começa na posição B. O cômodo A está sujo; B e C
  estão limpos. O agente segue exatamente a lógica abaixo:

  \begin{verbatim}
  def agente_aspirador(percepcao):
      posicao, sujo = percepcao
      if sujo:        return "ASPIRAR"
      if posicao == "A": return "DIREITA"
      if posicao == "B": return "DIREITA"
      return "ESQUERDA"
  \end{verbatim}

  Preencha a tabela de execução para os \textbf{6 primeiros passos}, indicando:
  posição atual, sensor (sujo/limpo), ação tomada, e estado dos cômodos após a ação.
  O agente limpa todos os cômodos? Em quantos passos?

  \sol{
    \textbf{Resposta:}

    \begin{tabular}{c|c|c|c|c}
      Passo & Posição & Sensor & Ação & Cômodos [A,B,C] \\
      \hline
      0 & B & limpo & DIREITA & [sujo, limpo, limpo] \\
      1 & C & limpo & ESQUERDA & [sujo, limpo, limpo] \\
      2 & B & limpo & DIREITA & [sujo, limpo, limpo] \\
      3 & C & limpo & ESQUERDA & [sujo, limpo, limpo] \\
      \dots & \dots & \dots & \dots & \dots \\
    \end{tabular}

    O agente \textbf{nunca alcança o cômodo A} porque sua lógica só vai de A→B→C→B→C→B→C\ldots
    Ele fica preso no ciclo B↔C. Para alcançar A, seria preciso uma regra adicional
    (ex.: "se estou em C e ambos A e B estão sujos, vá para B").

    \textbf{Resolução:} O bug está na lógica de movimento: quando a posição é B ou C,
    o agente nunca se move para A. A condição \texttt{posicao == "A"} nunca é verdadeira
    se o agente não consegue chegar a A.

    \errocomum{Assumir que o agente "dá a volta". O ambiente é uma linha, não um ciclo —
    de C só se vai para B.}

    \competencia{Implementar o agente aspirador e prever seu comportamento
    com 3 cômodos.}
  }
\end{questao}

\begin{questao}{5}{●●}
  O motor de encadeamento progressivo abaixo recebe uma base de fatos inicial
  e um conjunto de regras. Execute o algoritmo \textbf{passo a passo} com os
  dados fornecidos e indique o conjunto final de fatos.

  \begin{verbatim}
  fatos = {"chovendo", "temperatura < 15"}
  regras = [
    ({"chovendo", "temperatura < 15"}, "risco de gelo"),
    ({"risco de gelo"},              "fechar pista"),
    ({"chovendo"},                   "pista molhada"),
  ]
  \end{verbatim}

  \sol{
    \textbf{Resposta:} fatos finais = \{\texttt{"chovendo"}, \texttt{"temperatura < 15"},
    \texttt{"risco de gelo"}, \texttt{"fechar pista"}, \texttt{"pista molhada"}\}

    \textbf{Resolução:}
    \begin{enumerate}[nosep]
      \item Iteração 1: chovendo + temperatura < 15 → adiciona \texttt{"risco de gelo"}.
      \item Iteração 2: risco de gelo → adiciona \texttt{"fechar pista"}.
      \item Iteração 3: chovendo → adiciona \texttt{"pista molhada"}.
      \item Iteração 4: nenhuma regra nova dispara → fim.
    \end{enumerate}

    \errocomum{Esquecer que as regras disparam em qualquer ordem — o algoritmo testa TODAS
    as regras a cada iteração, não para na primeira que dispara.}

    \competencia{Explicar o funcionamento de um motor de inferência e por que regras
    não escalam.}
  }
\end{questao}

\begin{questao}{6}{●●}
  Em 2023, um LLM foi capaz de passar no exame da ordem dos advogados (Bar Exam)
  dos EUA sem ter recebido uma única regra explícita sobre direito. Usando os
  conceitos da Semana 1, explique \textbf{por que isso é um fato tão significativo}
  na história da IA. Se não obteve uma resposta completa, use o roteiro:
  (1) como um sistema especialista resolveria o Bar Exam? (2) por que essa abordagem
  falharia? (3) o que o LLM fez de diferente?

  \sol{
    \textbf{Resposta:} Um sistema especialista exigiria que engenheiros de conhecimento
    codificassem milhares de regras jurídicas à mão — um gargalo intransponível. O LLM
    não recebeu regra nenhuma: aprendeu padrões estatísticos de centenas de gigabytes
    de texto jurídico e geral. Isso é a Bitter Lesson em ação: conhecimento aprendido
    de dados venceu conhecimento artesanal — em um domínio (direito) que se supunha
    exigir raciocínio simbólico explícito.

    \errocomum{Dizer que o LLM "decorou" as respostas. LLMs não têm memória explícita do
    corpus de treino — eles aprendem padrões estatísticos, não fazem lookup de frases.}

    \competencia{Enunciar a Bitter Lesson e dar exemplos que a corroboram.}
  }
\end{questao}

% ════════════════ NÍVEL ●●● (ANÁLISE) ════════════════
\section*{Nível ●●● \quad Análise}

\begin{questao}{7}{●●●}
  Considere a afirmação de Sutton (2019):

  \begin{quote}
  \textit{``Métodos gerais que alavancam computação são, em última instância,
  os mais efetivos — e por uma grande margem. A maioria do esforço em IA é
  gasto tentando construir conhecimento em agentes; isso ajuda no curto prazo,
  mas no longo prazo atinge um platô.''}
  \end{quote}

  \begin{enumerate}[label=(\alph*),nosep]
    \item Por que conhecimento artesanal \textit{necessariamente} atinge um platô,
    enquanto métodos que aprendem de dados não? (Dica: pense no que acontece quando
    você dobra a quantidade de dados ou de computação em cada caso.)
    \item Dê um exemplo \textbf{concreto} de um sistema que atingiu esse platô e
    explique \textit{o que especificamente} o limitou.
  \end{enumerate}

  \sol{
    \textbf{Respostas:}
    (a) Conhecimento artesanal atinge um platô porque sua melhoria depende de \textbf{humanos
    adicionando mais regras} — um processo que é linear, caro e sujeito a contradições.
    Métodos que aprendem melhoram quando você adiciona mais dados ou computação, e esse
    processo é (em princípio) ilimitado: mais dados → melhores gradientes → melhor modelo.
    Não há um humano no loop ditando o teto.

    (b) Sistemas especialistas médicos (ex.: MYCIN, anos 80). Eram melhores que médicos
    juniores em diagnósticos específicos, mas cada nova doença exigia novas regras escritas
    por especialistas. Quando o número de regras passou de centenas, as interações entre
    elas tornaram-se imprevisíveis. O sistema não melhorava com mais dados de pacientes —
    só com mais regras escritas à mão. Enquanto isso, métodos de ML atuais melhoram
    automaticamente com mais exames, prontuários e artigos.

    \errocomum{Em (a), dizer que conhecimento artesanal atinge platô "porque humanos
    são limitados". A resposta precisa identificar o mecanismo: a taxa de melhoria é
    ditada pela velocidade com que humanos escrevem regras, não pela velocidade com
    que computação cresce.}

    \competencia{Enunciar a Bitter Lesson e explicar o mecanismo por trás dela.}
  }
\end{questao}

\begin{questao}{8}{●●●}
  O efeito ELIZA — atribuir inteligência onde há apenas casamento de padrões —
  foi documentado em 1966. Ele persiste em 2025 com LLMs?

  \begin{enumerate}[label=(\alph*),nosep]
    \item Aponte \textbf{uma semelhança} entre ELIZA e um LLM moderno que torna
    o efeito ELIZA relevante hoje.
    \item Aponte \textbf{uma diferença fundamental} que torna a comparação enganosa.
    \item Em no máximo 4 linhas: um LLM é ``só casamento de padrões''? Justifique.
  \end{enumerate}

  \sol{
    \textbf{Respostas:}
    (a) Ambos produzem texto fluente sem garantia de compreensão do conteúdo semântico.
    Um LLM pode gerar uma explicação convincente de um conceito que ele não ``entende''
    no sentido humano — assim como ELIZA simulava empatia sem sentir nada. O efeito
    ELIZA ocorre quando o usuário projeta compreensão onde há apenas correlação estatística.

    (b) ELIZA usava algumas dezenas de regras de reescrita escritas à mão. Um LLM aprende
    de centenas de bilhões de parâmetros ajustados sobre trilhões de tokens. A escala e
    a origem do conhecimento (artesanal vs. aprendido) são de natureza completamente diferente.

    (c) Não. Reduzir um LLM a ``só casamento de padrões'' é como reduzir um cérebro a
    ``só sinapses disparando'' — tecnicamente verdadeiro, mas enganoso. A questão relevante
    não é se o mecanismo básico é simples (é: multiplicação de matrizes + não-linearidades),
    mas se a escala produz comportamentos emergentes que merecem ser chamados de raciocínio.
    Esta é uma pergunta em aberto na pesquisa em 2025.

    \errocomum{Responder (c) com um simples ``sim'' ou ``não''. A pergunta exige distinguir
    mecanismo de comportamento — o que o sistema \textit{faz} vs. como ele \textit{faz}.}

    \competencia{Classificar sistemas de IA e analisar criticamente alegações de inteligência.}
  }
\end{questao}

% ════════════════ NÍVEL ★ (SÍNTESE) ════════════════
\section*{Nível ★ \quad Síntese}

\begin{questao}{9}{★}
  Projete um experimento mental: você herda um sistema especialista médico com
  10.000 regras, desenvolvido ao longo de 15 anos por uma equipe de 20 engenheiros
  de conhecimento. O sistema tem 92\% de acurácia em diagnóstico.

  \begin{enumerate}[label=(\alph*),nosep]
    \item Identifique \textbf{três} cenários em que este sistema falharia de forma
    sistemática que um modelo de ML treinado sobre os mesmos dados não falharia.
    \item Para cada cenário, explique \textbf{por que} a natureza artesanal das
    regras é a causa raiz da falha.
    \item Proponha uma arquitetura híbrida (simbólico + aprendizado) que combine
    o melhor dos dois mundos para este domínio. Descreva em no máximo 6 linhas.
  \end{enumerate}

  \sol{
    \textbf{Respostas:}
    (a) e (b):
    \begin{enumerate}[nosep]
      \item \textbf{Nova doença:} o SE falha porque nenhuma regra cobre o fenômeno novo.
      O ML pode detectar que os sintomas formam um cluster distinto nos dados e sinalizar
      uma anomalia — mesmo sem rótulos. Causa raiz: SE só reconhece o que foi explicitamente
      programado.
      \item \textbf{Interação entre condições:} um paciente com diabetes \textit{e} COVID
      pode ativar regras contraditórias (uma recomenda corticoides, outra contraindica).
      O SE não tem como ponderar evidências conflitantes. O ML aprende pesos relativos
      dos features. Causa raiz: regras são discretas e não-compostas; não modelam interações.
      \item \textbf{Deriva de distribuição:} se a população de pacientes mudar (ex.:
      envelhecimento), a acurácia do SE cai porque as regras foram escritas para a
      distribuição antiga. O ML pode ser retreinado com dados novos. Causa raiz: SE
      não se adapta sem reengenharia manual.
    \end{enumerate}

    (c) \textbf{Arquitetura híbrida:} usar um modelo de ML (ex.: gradient-boosted trees
    ou rede neural) para produzir o diagnóstico base e a estimativa de confiança. As
    10.000 regras viram um sistema de \textbf{verificação}: antes de emitir o diagnóstico,
    o sistema checa se ele viola alguma regra de segurança crítica (ex.: ``nunca prescrever
    droga X para grávidas''). Se violar, escala para um médico humano. O ML lida com a
    generalização; as regras lidam com as restrições rígidas de segurança.

    \errocomum{Propor um híbrido em que as regras \textit{guiam} o aprendizado. Isso
    reintroduz o gargalo: as regras limitam o que o modelo pode aprender. O melhor
    híbrido usa regras como \textit{restrições de segurança}, não como indutores de viés.}

    \competencia{Sintetizar os conceitos de IA simbólica e aprendizado em uma arquitetura
    fundamentada.}
  }
\end{questao}

\begin{questao}{10}{★ 🌉}
  O treinamento do GPT usa a mesma estrutura de loop que o agente aspirador da Aula B:
  \textbf{percepção → decisão → ação}. A diferença está na complexidade de \texttt{decidir()}.

  \begin{enumerate}[label=(\alph*),nosep]
    \item No aspirador, \texttt{decidir()} é uma função com 4 regras \texttt{if}.
    No GPT, \texttt{decidir()} é uma rede neural com centenas de bilhões de
    parâmetros. \textbf{Por que}, conceitualmente, a mesma interface de 4 linhas
    funciona para ambos? O que isso revela sobre o paradigma de agentes?
    \item A Bitter Lesson argumenta que conhecimento artesanal atinge um platô.
    O prompt engineering — escrever instruções cuidadosas para um LLM — é uma
    forma de conhecimento artesanal? Justifique em até 5 linhas.
  \end{enumerate}

  \sol{
    \textbf{Respostas:}
    (a) A interface \texttt{obs → decidir(obs) → ação} funciona para ambos porque ela
    abstrai \textbf{o que} o agente faz de \textbf{como} ele decide. O paradigma de
    agentes é uma \textit{casca}: define a assinatura da decisão, não sua implementação.
    O que muda entre o aspirador e o GPT é apenas a função \texttt{decidir()} — e essa
    função pode ser desde 4 ifs até uma rede de 175 bilhões de parâmetros. Isso revela
    que o paradigma de agentes é \textbf{universal}: ele unifica toda a IA sob uma
    única abstração.

    (b) Sim, prompt engineering é conhecimento artesanal — mas de um tipo novo. É
    artesanal porque um humano está manualmente ajustando a entrada para melhorar
    a saída, sem que o modelo aprenda com esse ajuste. Mas difere do artesanal
    clássico porque (1) o modelo subjacente foi aprendido de dados, não programado,
    e (2) o custo de ajustar um prompt é ordens de grandeza menor que o de escrever
    regras. Dito isso, a Bitter Lesson sugere que, no longo prazo, métodos que
    \textit{aprendem} a fazer prompting (auto-prompting, RLHF) substituirão o
    prompting manual — e já há evidências disso.

    \errocomum{Em (b), responder simplesmente ``não'' porque o modelo foi treinado com ML.
    O prompt em si é conhecimento artesanal injetado \textit{na inferência}, não no treino.}

    \competencia{Reconhecer o loop percepção→decisão→ação como a interface universal
    da IA e aplicar a Bitter Lesson a desenvolvimentos recentes.}
  }
\end{questao}

\vspace{1cm}
\hrule
\vspace{0.3cm}
\begin{center}
  {\footnotesize Lista S01 · GBC063 · Prof. Marcelo Keese Albertini · FACOM/UFU}
  \ifdefined\GABARITO
    \ifnum\GABARITO=1
      \\[0.2cm] {\footnotesize\color{red}\textbf{GABARITO — USO EXCLUSIVO DO PROFESSOR}}
    \fi
  \fi
\end{center}

\end{document}
"
% ═══════════════════════════════════════════════════════════════

\documentclass[12pt,a4paper]{article}

\usepackage[utf8]{inputenc}
\usepackage[T1]{fontenc}
\usepackage[brazil]{babel}
\usepackage{geometry}
\geometry{margin=2.2cm}
\usepackage{enumitem}
\usepackage{amsmath,amssymb}
\usepackage{xcolor}
\usepackage{hyperref}
% Using standard fonts available on any TeX distribution
\usepackage{mathpazo}
\usepackage[scaled=0.9]{helvet}
\usepackage{courier}
\renewcommand{\familydefault}{\rmdefault}

% ── Conditional: gabarito on/off ──
\ifdefined\GABARITO
  \ifnum\GABARITO=1
    \newcommand{\sol}[1]{\textcolor{blue}{#1}}
    \newcommand{\errocomum}[1]{\textcolor{orange}{⚠️ Erro comum: #1}}
    \newcommand{\competencia}[1]{\textcolor{gray}{Competência: #1}}
  \else
    \newcommand{\sol}[1]{}
    \newcommand{\errocomum}[1]{}
    \newcommand{\competencia}[1]{}
  \fi
\else
  \newcommand{\sol}[1]{}
  \newcommand{\errocomum}[1]{}
  \newcommand{\competencia}[1]{}
\fi

% ── Question formatting ──
\newcommand{\nivel}[1]{\textcolor{gray}{[#1]}}
\newenvironment{questao}[2]{
  \vspace{0.5cm}
  \noindent\textbf{Q#1} \nivel{#2} \quad
}{
  \vspace{0.2cm}
}

\begin{document}

% ════════════════ HEADER ════════════════
\begin{center}
  {\sffamily\Large GBC063 · Semana 1 · Inteligência Artificial}

  \vspace{0.2cm}
  {\small Prof. Marcelo Keese Albertini · FACOM · UFU}

  \vspace{0.2cm}
  {\footnotesize Cobre as competências do resumo S01 ·
  \href{../semana01/material_introducao_ia.html}{material de estudo} ·
  \href{../semana01/resumo_s01_introducao_ia.md}{resumo}}

  \ifdefined\GABARITO
    \ifnum\GABARITO=1
      \vspace{0.3cm}
      \fbox{\textcolor{red}{\textbf{GABARITO — Documento Exclusivo do Professor}}}
    \fi
  \fi
\end{center}

\vspace{0.3cm}
\hrule
\vspace{0.5cm}

% ════════════════ NÍVEL ● (REATIVAÇÃO) ════════════════
\section*{Nível ● \quad Reativação}

\begin{questao}{1}{●}
  Defina \textbf{agente racional} segundo Russell \& Norvig e explique,
  em no máximo 3 linhas, por que esta definição — e não o Teste de Turing —
  é adotada como definição operacional do curso.

  \sol{
    \textbf{Resposta:} Um agente racional é aquele que, para cada percepção possível,
    seleciona a ação que maximiza a métrica de desempenho esperada. O curso adota esta
    definição porque ela (1) funciona com qualquer métrica de sucesso, (2) não depende
    de replicar comportamento humano, e (3) unifica desde um termostato até o AlphaGo.

    \textbf{Resolução:} Russell \& Norvig (2013), Cap. 2. A definição de agente racional
    é independente de humanidade — um agente pode ser racional e não-humano (ex.: AlphaGo
    joga Go de forma sobre-humana).

    \errocomum{Definir agente racional como "agente que age como um humano". Isso é a visão
    "agir humanamente" (Teste de Turing), não "agir racionalmente".}

    \competencia{Explicar a diferença entre IA simbólica e aprendizado de máquina —
    e definir agente racional.}
  }
\end{questao}

\begin{questao}{2}{●}
  As quatro visões clássicas da IA organizam-se em uma matriz 2×2:
  pensamento \textit{vs.} ação × desempenho humano \textit{vs.} racionalidade.
  Classifique cada um dos sistemas abaixo em \textbf{um} dos quatro quadrantes
  e justifique em 1 linha:

  \begin{enumerate}[label=(\alph*),nosep]
    \item ELIZA (Weizenbaum, 1966)
    \item Deep Blue (IBM, 1997)
    \item O agente aspirador de pó da Aula B
    \item Um modelo de linguagem treinado para prever o próximo token
  \end{enumerate}

  \sol{
    \textbf{Respostas:}
    (a) Agir humanamente — ELIZA simulava um terapeuta humano sem raciocinar.
    (b) Agir racionalmente — Deep Blue selecionava lances que maximizavam chance de vitória,
    sem imitar o pensamento de um grande mestre humano.
    (c) Agir racionalmente — o aspirador maximiza a limpeza dada sua percepção.
    (d) Pensar racionalmente (durante o treino: otimiza uma função de perda) e
    agir humanamente (na inferência: gera texto indistinguível do humano).

    \errocomum{Classificar ELIZA como "pensar humanamente". ELIZA não modelava pensamento —
    apenas casava padrões de string.}

    \competencia{Classificar qualquer sistema de IA nas quatro visões e justificar.}
  }
\end{questao}

\begin{questao}{3}{●}
  Verdadeiro ou falso. \textbf{Justifique cada resposta em até 2 linhas.}

  \begin{enumerate}[label=(\alph*),nosep]
    \item "O primeiro inverno da IA (1969–1980) foi causado exclusivamente por falta
    de poder computacional."
    \item "Redes neurais artificiais são uma ideia dos anos 2010."
  \end{enumerate}

  \sol{
    \textbf{Respostas:}
    (a) \textbf{Falso.} O inverno foi causado principalmente pelo livro \textit{Perceptrons}
    (Minsky \& Papert, 1969), que provou limitações do perceptron simples e levou o campo
    a concluir — erroneamente — que redes neurais eram um beco sem saída. A falta de
    computação contribuiu, mas o fator decisivo foi a perda de credibilidade e financiamento.

    (b) \textbf{Falso.} O primeiro modelo de neurônio artificial é de McCulloch \& Pitts (1943).
    O perceptron é de Rosenblatt (1959). O backpropagation que as treina é de 1986
    (Rumelhart, Hinton \& Williams). As ideias fundamentais têm 40–80 anos; o que mudou
    em 2012 foi a escala (dados + GPUs).

    \errocomum{Para (a): atribuir o inverno só à falta de hardware. O fator principal foi
    sociológico — a comunidade perdeu confiança nas redes neurais.}
    \errocomum{Para (b): achar que deep learning é uma invenção recente. As ideias são antigas;
    o que é novo é a escala.}

    \competencia{Enunciar a Bitter Lesson e dar exemplos históricos.}
  }
\end{questao}

% ════════════════ NÍVEL ●● (EXECUÇÃO) ════════════════
\section*{Nível ●● \quad Execução}

\begin{questao}{4}{●●}
  Considere o seguinte ambiente de aspirador de pó com \textbf{três} cômodos
  em linha (A–B–C). O agente começa na posição B. O cômodo A está sujo; B e C
  estão limpos. O agente segue exatamente a lógica abaixo:

  \begin{verbatim}
  def agente_aspirador(percepcao):
      posicao, sujo = percepcao
      if sujo:        return "ASPIRAR"
      if posicao == "A": return "DIREITA"
      if posicao == "B": return "DIREITA"
      return "ESQUERDA"
  \end{verbatim}

  Preencha a tabela de execução para os \textbf{6 primeiros passos}, indicando:
  posição atual, sensor (sujo/limpo), ação tomada, e estado dos cômodos após a ação.
  O agente limpa todos os cômodos? Em quantos passos?

  \sol{
    \textbf{Resposta:}

    \begin{tabular}{c|c|c|c|c}
      Passo & Posição & Sensor & Ação & Cômodos [A,B,C] \\
      \hline
      0 & B & limpo & DIREITA & [sujo, limpo, limpo] \\
      1 & C & limpo & ESQUERDA & [sujo, limpo, limpo] \\
      2 & B & limpo & DIREITA & [sujo, limpo, limpo] \\
      3 & C & limpo & ESQUERDA & [sujo, limpo, limpo] \\
      \dots & \dots & \dots & \dots & \dots \\
    \end{tabular}

    O agente \textbf{nunca alcança o cômodo A} porque sua lógica só vai de A→B→C→B→C→B→C\ldots
    Ele fica preso no ciclo B↔C. Para alcançar A, seria preciso uma regra adicional
    (ex.: "se estou em C e ambos A e B estão sujos, vá para B").

    \textbf{Resolução:} O bug está na lógica de movimento: quando a posição é B ou C,
    o agente nunca se move para A. A condição \texttt{posicao == "A"} nunca é verdadeira
    se o agente não consegue chegar a A.

    \errocomum{Assumir que o agente "dá a volta". O ambiente é uma linha, não um ciclo —
    de C só se vai para B.}

    \competencia{Implementar o agente aspirador e prever seu comportamento
    com 3 cômodos.}
  }
\end{questao}

\begin{questao}{5}{●●}
  O motor de encadeamento progressivo abaixo recebe uma base de fatos inicial
  e um conjunto de regras. Execute o algoritmo \textbf{passo a passo} com os
  dados fornecidos e indique o conjunto final de fatos.

  \begin{verbatim}
  fatos = {"chovendo", "temperatura < 15"}
  regras = [
    ({"chovendo", "temperatura < 15"}, "risco de gelo"),
    ({"risco de gelo"},              "fechar pista"),
    ({"chovendo"},                   "pista molhada"),
  ]
  \end{verbatim}

  \sol{
    \textbf{Resposta:} fatos finais = \{\texttt{"chovendo"}, \texttt{"temperatura < 15"},
    \texttt{"risco de gelo"}, \texttt{"fechar pista"}, \texttt{"pista molhada"}\}

    \textbf{Resolução:}
    \begin{enumerate}[nosep]
      \item Iteração 1: chovendo + temperatura < 15 → adiciona \texttt{"risco de gelo"}.
      \item Iteração 2: risco de gelo → adiciona \texttt{"fechar pista"}.
      \item Iteração 3: chovendo → adiciona \texttt{"pista molhada"}.
      \item Iteração 4: nenhuma regra nova dispara → fim.
    \end{enumerate}

    \errocomum{Esquecer que as regras disparam em qualquer ordem — o algoritmo testa TODAS
    as regras a cada iteração, não para na primeira que dispara.}

    \competencia{Explicar o funcionamento de um motor de inferência e por que regras
    não escalam.}
  }
\end{questao}

\begin{questao}{6}{●●}
  Em 2023, um LLM foi capaz de passar no exame da ordem dos advogados (Bar Exam)
  dos EUA sem ter recebido uma única regra explícita sobre direito. Usando os
  conceitos da Semana 1, explique \textbf{por que isso é um fato tão significativo}
  na história da IA. Se não obteve uma resposta completa, use o roteiro:
  (1) como um sistema especialista resolveria o Bar Exam? (2) por que essa abordagem
  falharia? (3) o que o LLM fez de diferente?

  \sol{
    \textbf{Resposta:} Um sistema especialista exigiria que engenheiros de conhecimento
    codificassem milhares de regras jurídicas à mão — um gargalo intransponível. O LLM
    não recebeu regra nenhuma: aprendeu padrões estatísticos de centenas de gigabytes
    de texto jurídico e geral. Isso é a Bitter Lesson em ação: conhecimento aprendido
    de dados venceu conhecimento artesanal — em um domínio (direito) que se supunha
    exigir raciocínio simbólico explícito.

    \errocomum{Dizer que o LLM "decorou" as respostas. LLMs não têm memória explícita do
    corpus de treino — eles aprendem padrões estatísticos, não fazem lookup de frases.}

    \competencia{Enunciar a Bitter Lesson e dar exemplos que a corroboram.}
  }
\end{questao}

% ════════════════ NÍVEL ●●● (ANÁLISE) ════════════════
\section*{Nível ●●● \quad Análise}

\begin{questao}{7}{●●●}
  Considere a afirmação de Sutton (2019):

  \begin{quote}
  \textit{``Métodos gerais que alavancam computação são, em última instância,
  os mais efetivos — e por uma grande margem. A maioria do esforço em IA é
  gasto tentando construir conhecimento em agentes; isso ajuda no curto prazo,
  mas no longo prazo atinge um platô.''}
  \end{quote}

  \begin{enumerate}[label=(\alph*),nosep]
    \item Por que conhecimento artesanal \textit{necessariamente} atinge um platô,
    enquanto métodos que aprendem de dados não? (Dica: pense no que acontece quando
    você dobra a quantidade de dados ou de computação em cada caso.)
    \item Dê um exemplo \textbf{concreto} de um sistema que atingiu esse platô e
    explique \textit{o que especificamente} o limitou.
  \end{enumerate}

  \sol{
    \textbf{Respostas:}
    (a) Conhecimento artesanal atinge um platô porque sua melhoria depende de \textbf{humanos
    adicionando mais regras} — um processo que é linear, caro e sujeito a contradições.
    Métodos que aprendem melhoram quando você adiciona mais dados ou computação, e esse
    processo é (em princípio) ilimitado: mais dados → melhores gradientes → melhor modelo.
    Não há um humano no loop ditando o teto.

    (b) Sistemas especialistas médicos (ex.: MYCIN, anos 80). Eram melhores que médicos
    juniores em diagnósticos específicos, mas cada nova doença exigia novas regras escritas
    por especialistas. Quando o número de regras passou de centenas, as interações entre
    elas tornaram-se imprevisíveis. O sistema não melhorava com mais dados de pacientes —
    só com mais regras escritas à mão. Enquanto isso, métodos de ML atuais melhoram
    automaticamente com mais exames, prontuários e artigos.

    \errocomum{Em (a), dizer que conhecimento artesanal atinge platô "porque humanos
    são limitados". A resposta precisa identificar o mecanismo: a taxa de melhoria é
    ditada pela velocidade com que humanos escrevem regras, não pela velocidade com
    que computação cresce.}

    \competencia{Enunciar a Bitter Lesson e explicar o mecanismo por trás dela.}
  }
\end{questao}

\begin{questao}{8}{●●●}
  O efeito ELIZA — atribuir inteligência onde há apenas casamento de padrões —
  foi documentado em 1966. Ele persiste em 2025 com LLMs?

  \begin{enumerate}[label=(\alph*),nosep]
    \item Aponte \textbf{uma semelhança} entre ELIZA e um LLM moderno que torna
    o efeito ELIZA relevante hoje.
    \item Aponte \textbf{uma diferença fundamental} que torna a comparação enganosa.
    \item Em no máximo 4 linhas: um LLM é ``só casamento de padrões''? Justifique.
  \end{enumerate}

  \sol{
    \textbf{Respostas:}
    (a) Ambos produzem texto fluente sem garantia de compreensão do conteúdo semântico.
    Um LLM pode gerar uma explicação convincente de um conceito que ele não ``entende''
    no sentido humano — assim como ELIZA simulava empatia sem sentir nada. O efeito
    ELIZA ocorre quando o usuário projeta compreensão onde há apenas correlação estatística.

    (b) ELIZA usava algumas dezenas de regras de reescrita escritas à mão. Um LLM aprende
    de centenas de bilhões de parâmetros ajustados sobre trilhões de tokens. A escala e
    a origem do conhecimento (artesanal vs. aprendido) são de natureza completamente diferente.

    (c) Não. Reduzir um LLM a ``só casamento de padrões'' é como reduzir um cérebro a
    ``só sinapses disparando'' — tecnicamente verdadeiro, mas enganoso. A questão relevante
    não é se o mecanismo básico é simples (é: multiplicação de matrizes + não-linearidades),
    mas se a escala produz comportamentos emergentes que merecem ser chamados de raciocínio.
    Esta é uma pergunta em aberto na pesquisa em 2025.

    \errocomum{Responder (c) com um simples ``sim'' ou ``não''. A pergunta exige distinguir
    mecanismo de comportamento — o que o sistema \textit{faz} vs. como ele \textit{faz}.}

    \competencia{Classificar sistemas de IA e analisar criticamente alegações de inteligência.}
  }
\end{questao}

% ════════════════ NÍVEL ★ (SÍNTESE) ════════════════
\section*{Nível ★ \quad Síntese}

\begin{questao}{9}{★}
  Projete um experimento mental: você herda um sistema especialista médico com
  10.000 regras, desenvolvido ao longo de 15 anos por uma equipe de 20 engenheiros
  de conhecimento. O sistema tem 92\% de acurácia em diagnóstico.

  \begin{enumerate}[label=(\alph*),nosep]
    \item Identifique \textbf{três} cenários em que este sistema falharia de forma
    sistemática que um modelo de ML treinado sobre os mesmos dados não falharia.
    \item Para cada cenário, explique \textbf{por que} a natureza artesanal das
    regras é a causa raiz da falha.
    \item Proponha uma arquitetura híbrida (simbólico + aprendizado) que combine
    o melhor dos dois mundos para este domínio. Descreva em no máximo 6 linhas.
  \end{enumerate}

  \sol{
    \textbf{Respostas:}
    (a) e (b):
    \begin{enumerate}[nosep]
      \item \textbf{Nova doença:} o SE falha porque nenhuma regra cobre o fenômeno novo.
      O ML pode detectar que os sintomas formam um cluster distinto nos dados e sinalizar
      uma anomalia — mesmo sem rótulos. Causa raiz: SE só reconhece o que foi explicitamente
      programado.
      \item \textbf{Interação entre condições:} um paciente com diabetes \textit{e} COVID
      pode ativar regras contraditórias (uma recomenda corticoides, outra contraindica).
      O SE não tem como ponderar evidências conflitantes. O ML aprende pesos relativos
      dos features. Causa raiz: regras são discretas e não-compostas; não modelam interações.
      \item \textbf{Deriva de distribuição:} se a população de pacientes mudar (ex.:
      envelhecimento), a acurácia do SE cai porque as regras foram escritas para a
      distribuição antiga. O ML pode ser retreinado com dados novos. Causa raiz: SE
      não se adapta sem reengenharia manual.
    \end{enumerate}

    (c) \textbf{Arquitetura híbrida:} usar um modelo de ML (ex.: gradient-boosted trees
    ou rede neural) para produzir o diagnóstico base e a estimativa de confiança. As
    10.000 regras viram um sistema de \textbf{verificação}: antes de emitir o diagnóstico,
    o sistema checa se ele viola alguma regra de segurança crítica (ex.: ``nunca prescrever
    droga X para grávidas''). Se violar, escala para um médico humano. O ML lida com a
    generalização; as regras lidam com as restrições rígidas de segurança.

    \errocomum{Propor um híbrido em que as regras \textit{guiam} o aprendizado. Isso
    reintroduz o gargalo: as regras limitam o que o modelo pode aprender. O melhor
    híbrido usa regras como \textit{restrições de segurança}, não como indutores de viés.}

    \competencia{Sintetizar os conceitos de IA simbólica e aprendizado em uma arquitetura
    fundamentada.}
  }
\end{questao}

\begin{questao}{10}{★ 🌉}
  O treinamento do GPT usa a mesma estrutura de loop que o agente aspirador da Aula B:
  \textbf{percepção → decisão → ação}. A diferença está na complexidade de \texttt{decidir()}.

  \begin{enumerate}[label=(\alph*),nosep]
    \item No aspirador, \texttt{decidir()} é uma função com 4 regras \texttt{if}.
    No GPT, \texttt{decidir()} é uma rede neural com centenas de bilhões de
    parâmetros. \textbf{Por que}, conceitualmente, a mesma interface de 4 linhas
    funciona para ambos? O que isso revela sobre o paradigma de agentes?
    \item A Bitter Lesson argumenta que conhecimento artesanal atinge um platô.
    O prompt engineering — escrever instruções cuidadosas para um LLM — é uma
    forma de conhecimento artesanal? Justifique em até 5 linhas.
  \end{enumerate}

  \sol{
    \textbf{Respostas:}
    (a) A interface \texttt{obs → decidir(obs) → ação} funciona para ambos porque ela
    abstrai \textbf{o que} o agente faz de \textbf{como} ele decide. O paradigma de
    agentes é uma \textit{casca}: define a assinatura da decisão, não sua implementação.
    O que muda entre o aspirador e o GPT é apenas a função \texttt{decidir()} — e essa
    função pode ser desde 4 ifs até uma rede de 175 bilhões de parâmetros. Isso revela
    que o paradigma de agentes é \textbf{universal}: ele unifica toda a IA sob uma
    única abstração.

    (b) Sim, prompt engineering é conhecimento artesanal — mas de um tipo novo. É
    artesanal porque um humano está manualmente ajustando a entrada para melhorar
    a saída, sem que o modelo aprenda com esse ajuste. Mas difere do artesanal
    clássico porque (1) o modelo subjacente foi aprendido de dados, não programado,
    e (2) o custo de ajustar um prompt é ordens de grandeza menor que o de escrever
    regras. Dito isso, a Bitter Lesson sugere que, no longo prazo, métodos que
    \textit{aprendem} a fazer prompting (auto-prompting, RLHF) substituirão o
    prompting manual — e já há evidências disso.

    \errocomum{Em (b), responder simplesmente ``não'' porque o modelo foi treinado com ML.
    O prompt em si é conhecimento artesanal injetado \textit{na inferência}, não no treino.}

    \competencia{Reconhecer o loop percepção→decisão→ação como a interface universal
    da IA e aplicar a Bitter Lesson a desenvolvimentos recentes.}
  }
\end{questao}

\vspace{1cm}
\hrule
\vspace{0.3cm}
\begin{center}
  {\footnotesize Lista S01 · GBC063 · Prof. Marcelo Keese Albertini · FACOM/UFU}
  \ifdefined\GABARITO
    \ifnum\GABARITO=1
      \\[0.2cm] {\footnotesize\color{red}\textbf{GABARITO — USO EXCLUSIVO DO PROFESSOR}}
    \fi
  \fi
\end{center}

\end{document}
"
% ═══════════════════════════════════════════════════════════════

\documentclass[12pt,a4paper]{article}

\usepackage[utf8]{inputenc}
\usepackage[T1]{fontenc}
\usepackage[brazil]{babel}
\usepackage{geometry}
\geometry{margin=2.2cm}
\usepackage{enumitem}
\usepackage{amsmath,amssymb}
\usepackage{xcolor}
\usepackage{hyperref}
% Using standard fonts available on any TeX distribution
\usepackage{mathpazo}
\usepackage[scaled=0.9]{helvet}
\usepackage{courier}
\renewcommand{\familydefault}{\rmdefault}

% ── Conditional: gabarito on/off ──
\ifdefined\GABARITO
  \ifnum\GABARITO=1
    \newcommand{\sol}[1]{\textcolor{blue}{#1}}
    \newcommand{\errocomum}[1]{\textcolor{orange}{⚠️ Erro comum: #1}}
    \newcommand{\competencia}[1]{\textcolor{gray}{Competência: #1}}
  \else
    \newcommand{\sol}[1]{}
    \newcommand{\errocomum}[1]{}
    \newcommand{\competencia}[1]{}
  \fi
\else
  \newcommand{\sol}[1]{}
  \newcommand{\errocomum}[1]{}
  \newcommand{\competencia}[1]{}
\fi

% ── Question formatting ──
\newcommand{\nivel}[1]{\textcolor{gray}{[#1]}}
\newenvironment{questao}[2]{
  \vspace{0.5cm}
  \noindent\textbf{Q#1} \nivel{#2} \quad
}{
  \vspace{0.2cm}
}

\begin{document}

% ════════════════ HEADER ════════════════
\begin{center}
  {\sffamily\Large GBC063 · Semana 1 · Inteligência Artificial}

  \vspace{0.2cm}
  {\small Prof. Marcelo Keese Albertini · FACOM · UFU}

  \vspace{0.2cm}
  {\footnotesize Cobre as competências do resumo S01 ·
  \href{../semana01/material_introducao_ia.html}{material de estudo} ·
  \href{../semana01/resumo_s01_introducao_ia.md}{resumo}}

  \ifdefined\GABARITO
    \ifnum\GABARITO=1
      \vspace{0.3cm}
      \fbox{\textcolor{red}{\textbf{GABARITO — Documento Exclusivo do Professor}}}
    \fi
  \fi
\end{center}

\vspace{0.3cm}
\hrule
\vspace{0.5cm}

% ════════════════ NÍVEL ● (REATIVAÇÃO) ════════════════
\section*{Nível ● \quad Reativação}

\begin{questao}{1}{●}
  Defina \textbf{agente racional} segundo Russell \& Norvig e explique,
  em no máximo 3 linhas, por que esta definição — e não o Teste de Turing —
  é adotada como definição operacional do curso.

  \sol{
    \textbf{Resposta:} Um agente racional é aquele que, para cada percepção possível,
    seleciona a ação que maximiza a métrica de desempenho esperada. O curso adota esta
    definição porque ela (1) funciona com qualquer métrica de sucesso, (2) não depende
    de replicar comportamento humano, e (3) unifica desde um termostato até o AlphaGo.

    \textbf{Resolução:} Russell \& Norvig (2013), Cap. 2. A definição de agente racional
    é independente de humanidade — um agente pode ser racional e não-humano (ex.: AlphaGo
    joga Go de forma sobre-humana).

    \errocomum{Definir agente racional como "agente que age como um humano". Isso é a visão
    "agir humanamente" (Teste de Turing), não "agir racionalmente".}

    \competencia{Explicar a diferença entre IA simbólica e aprendizado de máquina —
    e definir agente racional.}
  }
\end{questao}

\begin{questao}{2}{●}
  As quatro visões clássicas da IA organizam-se em uma matriz 2×2:
  pensamento \textit{vs.} ação × desempenho humano \textit{vs.} racionalidade.
  Classifique cada um dos sistemas abaixo em \textbf{um} dos quatro quadrantes
  e justifique em 1 linha:

  \begin{enumerate}[label=(\alph*),nosep]
    \item ELIZA (Weizenbaum, 1966)
    \item Deep Blue (IBM, 1997)
    \item O agente aspirador de pó da Aula B
    \item Um modelo de linguagem treinado para prever o próximo token
  \end{enumerate}

  \sol{
    \textbf{Respostas:}
    (a) Agir humanamente — ELIZA simulava um terapeuta humano sem raciocinar.
    (b) Agir racionalmente — Deep Blue selecionava lances que maximizavam chance de vitória,
    sem imitar o pensamento de um grande mestre humano.
    (c) Agir racionalmente — o aspirador maximiza a limpeza dada sua percepção.
    (d) Pensar racionalmente (durante o treino: otimiza uma função de perda) e
    agir humanamente (na inferência: gera texto indistinguível do humano).

    \errocomum{Classificar ELIZA como "pensar humanamente". ELIZA não modelava pensamento —
    apenas casava padrões de string.}

    \competencia{Classificar qualquer sistema de IA nas quatro visões e justificar.}
  }
\end{questao}

\begin{questao}{3}{●}
  Verdadeiro ou falso. \textbf{Justifique cada resposta em até 2 linhas.}

  \begin{enumerate}[label=(\alph*),nosep]
    \item "O primeiro inverno da IA (1969–1980) foi causado exclusivamente por falta
    de poder computacional."
    \item "Redes neurais artificiais são uma ideia dos anos 2010."
  \end{enumerate}

  \sol{
    \textbf{Respostas:}
    (a) \textbf{Falso.} O inverno foi causado principalmente pelo livro \textit{Perceptrons}
    (Minsky \& Papert, 1969), que provou limitações do perceptron simples e levou o campo
    a concluir — erroneamente — que redes neurais eram um beco sem saída. A falta de
    computação contribuiu, mas o fator decisivo foi a perda de credibilidade e financiamento.

    (b) \textbf{Falso.} O primeiro modelo de neurônio artificial é de McCulloch \& Pitts (1943).
    O perceptron é de Rosenblatt (1959). O backpropagation que as treina é de 1986
    (Rumelhart, Hinton \& Williams). As ideias fundamentais têm 40–80 anos; o que mudou
    em 2012 foi a escala (dados + GPUs).

    \errocomum{Para (a): atribuir o inverno só à falta de hardware. O fator principal foi
    sociológico — a comunidade perdeu confiança nas redes neurais.}
    \errocomum{Para (b): achar que deep learning é uma invenção recente. As ideias são antigas;
    o que é novo é a escala.}

    \competencia{Enunciar a Bitter Lesson e dar exemplos históricos.}
  }
\end{questao}

% ════════════════ NÍVEL ●● (EXECUÇÃO) ════════════════
\section*{Nível ●● \quad Execução}

\begin{questao}{4}{●●}
  Considere o seguinte ambiente de aspirador de pó com \textbf{três} cômodos
  em linha (A–B–C). O agente começa na posição B. O cômodo A está sujo; B e C
  estão limpos. O agente segue exatamente a lógica abaixo:

  \begin{verbatim}
  def agente_aspirador(percepcao):
      posicao, sujo = percepcao
      if sujo:        return "ASPIRAR"
      if posicao == "A": return "DIREITA"
      if posicao == "B": return "DIREITA"
      return "ESQUERDA"
  \end{verbatim}

  Preencha a tabela de execução para os \textbf{6 primeiros passos}, indicando:
  posição atual, sensor (sujo/limpo), ação tomada, e estado dos cômodos após a ação.
  O agente limpa todos os cômodos? Em quantos passos?

  \sol{
    \textbf{Resposta:}

    \begin{tabular}{c|c|c|c|c}
      Passo & Posição & Sensor & Ação & Cômodos [A,B,C] \\
      \hline
      0 & B & limpo & DIREITA & [sujo, limpo, limpo] \\
      1 & C & limpo & ESQUERDA & [sujo, limpo, limpo] \\
      2 & B & limpo & DIREITA & [sujo, limpo, limpo] \\
      3 & C & limpo & ESQUERDA & [sujo, limpo, limpo] \\
      \dots & \dots & \dots & \dots & \dots \\
    \end{tabular}

    O agente \textbf{nunca alcança o cômodo A} porque sua lógica só vai de A→B→C→B→C→B→C\ldots
    Ele fica preso no ciclo B↔C. Para alcançar A, seria preciso uma regra adicional
    (ex.: "se estou em C e ambos A e B estão sujos, vá para B").

    \textbf{Resolução:} O bug está na lógica de movimento: quando a posição é B ou C,
    o agente nunca se move para A. A condição \texttt{posicao == "A"} nunca é verdadeira
    se o agente não consegue chegar a A.

    \errocomum{Assumir que o agente "dá a volta". O ambiente é uma linha, não um ciclo —
    de C só se vai para B.}

    \competencia{Implementar o agente aspirador e prever seu comportamento
    com 3 cômodos.}
  }
\end{questao}

\begin{questao}{5}{●●}
  O motor de encadeamento progressivo abaixo recebe uma base de fatos inicial
  e um conjunto de regras. Execute o algoritmo \textbf{passo a passo} com os
  dados fornecidos e indique o conjunto final de fatos.

  \begin{verbatim}
  fatos = {"chovendo", "temperatura < 15"}
  regras = [
    ({"chovendo", "temperatura < 15"}, "risco de gelo"),
    ({"risco de gelo"},              "fechar pista"),
    ({"chovendo"},                   "pista molhada"),
  ]
  \end{verbatim}

  \sol{
    \textbf{Resposta:} fatos finais = \{\texttt{"chovendo"}, \texttt{"temperatura < 15"},
    \texttt{"risco de gelo"}, \texttt{"fechar pista"}, \texttt{"pista molhada"}\}

    \textbf{Resolução:}
    \begin{enumerate}[nosep]
      \item Iteração 1: chovendo + temperatura < 15 → adiciona \texttt{"risco de gelo"}.
      \item Iteração 2: risco de gelo → adiciona \texttt{"fechar pista"}.
      \item Iteração 3: chovendo → adiciona \texttt{"pista molhada"}.
      \item Iteração 4: nenhuma regra nova dispara → fim.
    \end{enumerate}

    \errocomum{Esquecer que as regras disparam em qualquer ordem — o algoritmo testa TODAS
    as regras a cada iteração, não para na primeira que dispara.}

    \competencia{Explicar o funcionamento de um motor de inferência e por que regras
    não escalam.}
  }
\end{questao}

\begin{questao}{6}{●●}
  Em 2023, um LLM foi capaz de passar no exame da ordem dos advogados (Bar Exam)
  dos EUA sem ter recebido uma única regra explícita sobre direito. Usando os
  conceitos da Semana 1, explique \textbf{por que isso é um fato tão significativo}
  na história da IA. Se não obteve uma resposta completa, use o roteiro:
  (1) como um sistema especialista resolveria o Bar Exam? (2) por que essa abordagem
  falharia? (3) o que o LLM fez de diferente?

  \sol{
    \textbf{Resposta:} Um sistema especialista exigiria que engenheiros de conhecimento
    codificassem milhares de regras jurídicas à mão — um gargalo intransponível. O LLM
    não recebeu regra nenhuma: aprendeu padrões estatísticos de centenas de gigabytes
    de texto jurídico e geral. Isso é a Bitter Lesson em ação: conhecimento aprendido
    de dados venceu conhecimento artesanal — em um domínio (direito) que se supunha
    exigir raciocínio simbólico explícito.

    \errocomum{Dizer que o LLM "decorou" as respostas. LLMs não têm memória explícita do
    corpus de treino — eles aprendem padrões estatísticos, não fazem lookup de frases.}

    \competencia{Enunciar a Bitter Lesson e dar exemplos que a corroboram.}
  }
\end{questao}

% ════════════════ NÍVEL ●●● (ANÁLISE) ════════════════
\section*{Nível ●●● \quad Análise}

\begin{questao}{7}{●●●}
  Considere a afirmação de Sutton (2019):

  \begin{quote}
  \textit{``Métodos gerais que alavancam computação são, em última instância,
  os mais efetivos — e por uma grande margem. A maioria do esforço em IA é
  gasto tentando construir conhecimento em agentes; isso ajuda no curto prazo,
  mas no longo prazo atinge um platô.''}
  \end{quote}

  \begin{enumerate}[label=(\alph*),nosep]
    \item Por que conhecimento artesanal \textit{necessariamente} atinge um platô,
    enquanto métodos que aprendem de dados não? (Dica: pense no que acontece quando
    você dobra a quantidade de dados ou de computação em cada caso.)
    \item Dê um exemplo \textbf{concreto} de um sistema que atingiu esse platô e
    explique \textit{o que especificamente} o limitou.
  \end{enumerate}

  \sol{
    \textbf{Respostas:}
    (a) Conhecimento artesanal atinge um platô porque sua melhoria depende de \textbf{humanos
    adicionando mais regras} — um processo que é linear, caro e sujeito a contradições.
    Métodos que aprendem melhoram quando você adiciona mais dados ou computação, e esse
    processo é (em princípio) ilimitado: mais dados → melhores gradientes → melhor modelo.
    Não há um humano no loop ditando o teto.

    (b) Sistemas especialistas médicos (ex.: MYCIN, anos 80). Eram melhores que médicos
    juniores em diagnósticos específicos, mas cada nova doença exigia novas regras escritas
    por especialistas. Quando o número de regras passou de centenas, as interações entre
    elas tornaram-se imprevisíveis. O sistema não melhorava com mais dados de pacientes —
    só com mais regras escritas à mão. Enquanto isso, métodos de ML atuais melhoram
    automaticamente com mais exames, prontuários e artigos.

    \errocomum{Em (a), dizer que conhecimento artesanal atinge platô "porque humanos
    são limitados". A resposta precisa identificar o mecanismo: a taxa de melhoria é
    ditada pela velocidade com que humanos escrevem regras, não pela velocidade com
    que computação cresce.}

    \competencia{Enunciar a Bitter Lesson e explicar o mecanismo por trás dela.}
  }
\end{questao}

\begin{questao}{8}{●●●}
  O efeito ELIZA — atribuir inteligência onde há apenas casamento de padrões —
  foi documentado em 1966. Ele persiste em 2025 com LLMs?

  \begin{enumerate}[label=(\alph*),nosep]
    \item Aponte \textbf{uma semelhança} entre ELIZA e um LLM moderno que torna
    o efeito ELIZA relevante hoje.
    \item Aponte \textbf{uma diferença fundamental} que torna a comparação enganosa.
    \item Em no máximo 4 linhas: um LLM é ``só casamento de padrões''? Justifique.
  \end{enumerate}

  \sol{
    \textbf{Respostas:}
    (a) Ambos produzem texto fluente sem garantia de compreensão do conteúdo semântico.
    Um LLM pode gerar uma explicação convincente de um conceito que ele não ``entende''
    no sentido humano — assim como ELIZA simulava empatia sem sentir nada. O efeito
    ELIZA ocorre quando o usuário projeta compreensão onde há apenas correlação estatística.

    (b) ELIZA usava algumas dezenas de regras de reescrita escritas à mão. Um LLM aprende
    de centenas de bilhões de parâmetros ajustados sobre trilhões de tokens. A escala e
    a origem do conhecimento (artesanal vs. aprendido) são de natureza completamente diferente.

    (c) Não. Reduzir um LLM a ``só casamento de padrões'' é como reduzir um cérebro a
    ``só sinapses disparando'' — tecnicamente verdadeiro, mas enganoso. A questão relevante
    não é se o mecanismo básico é simples (é: multiplicação de matrizes + não-linearidades),
    mas se a escala produz comportamentos emergentes que merecem ser chamados de raciocínio.
    Esta é uma pergunta em aberto na pesquisa em 2025.

    \errocomum{Responder (c) com um simples ``sim'' ou ``não''. A pergunta exige distinguir
    mecanismo de comportamento — o que o sistema \textit{faz} vs. como ele \textit{faz}.}

    \competencia{Classificar sistemas de IA e analisar criticamente alegações de inteligência.}
  }
\end{questao}

% ════════════════ NÍVEL ★ (SÍNTESE) ════════════════
\section*{Nível ★ \quad Síntese}

\begin{questao}{9}{★}
  Projete um experimento mental: você herda um sistema especialista médico com
  10.000 regras, desenvolvido ao longo de 15 anos por uma equipe de 20 engenheiros
  de conhecimento. O sistema tem 92\% de acurácia em diagnóstico.

  \begin{enumerate}[label=(\alph*),nosep]
    \item Identifique \textbf{três} cenários em que este sistema falharia de forma
    sistemática que um modelo de ML treinado sobre os mesmos dados não falharia.
    \item Para cada cenário, explique \textbf{por que} a natureza artesanal das
    regras é a causa raiz da falha.
    \item Proponha uma arquitetura híbrida (simbólico + aprendizado) que combine
    o melhor dos dois mundos para este domínio. Descreva em no máximo 6 linhas.
  \end{enumerate}

  \sol{
    \textbf{Respostas:}
    (a) e (b):
    \begin{enumerate}[nosep]
      \item \textbf{Nova doença:} o SE falha porque nenhuma regra cobre o fenômeno novo.
      O ML pode detectar que os sintomas formam um cluster distinto nos dados e sinalizar
      uma anomalia — mesmo sem rótulos. Causa raiz: SE só reconhece o que foi explicitamente
      programado.
      \item \textbf{Interação entre condições:} um paciente com diabetes \textit{e} COVID
      pode ativar regras contraditórias (uma recomenda corticoides, outra contraindica).
      O SE não tem como ponderar evidências conflitantes. O ML aprende pesos relativos
      dos features. Causa raiz: regras são discretas e não-compostas; não modelam interações.
      \item \textbf{Deriva de distribuição:} se a população de pacientes mudar (ex.:
      envelhecimento), a acurácia do SE cai porque as regras foram escritas para a
      distribuição antiga. O ML pode ser retreinado com dados novos. Causa raiz: SE
      não se adapta sem reengenharia manual.
    \end{enumerate}

    (c) \textbf{Arquitetura híbrida:} usar um modelo de ML (ex.: gradient-boosted trees
    ou rede neural) para produzir o diagnóstico base e a estimativa de confiança. As
    10.000 regras viram um sistema de \textbf{verificação}: antes de emitir o diagnóstico,
    o sistema checa se ele viola alguma regra de segurança crítica (ex.: ``nunca prescrever
    droga X para grávidas''). Se violar, escala para um médico humano. O ML lida com a
    generalização; as regras lidam com as restrições rígidas de segurança.

    \errocomum{Propor um híbrido em que as regras \textit{guiam} o aprendizado. Isso
    reintroduz o gargalo: as regras limitam o que o modelo pode aprender. O melhor
    híbrido usa regras como \textit{restrições de segurança}, não como indutores de viés.}

    \competencia{Sintetizar os conceitos de IA simbólica e aprendizado em uma arquitetura
    fundamentada.}
  }
\end{questao}

\begin{questao}{10}{★ 🌉}
  O treinamento do GPT usa a mesma estrutura de loop que o agente aspirador da Aula B:
  \textbf{percepção → decisão → ação}. A diferença está na complexidade de \texttt{decidir()}.

  \begin{enumerate}[label=(\alph*),nosep]
    \item No aspirador, \texttt{decidir()} é uma função com 4 regras \texttt{if}.
    No GPT, \texttt{decidir()} é uma rede neural com centenas de bilhões de
    parâmetros. \textbf{Por que}, conceitualmente, a mesma interface de 4 linhas
    funciona para ambos? O que isso revela sobre o paradigma de agentes?
    \item A Bitter Lesson argumenta que conhecimento artesanal atinge um platô.
    O prompt engineering — escrever instruções cuidadosas para um LLM — é uma
    forma de conhecimento artesanal? Justifique em até 5 linhas.
  \end{enumerate}

  \sol{
    \textbf{Respostas:}
    (a) A interface \texttt{obs → decidir(obs) → ação} funciona para ambos porque ela
    abstrai \textbf{o que} o agente faz de \textbf{como} ele decide. O paradigma de
    agentes é uma \textit{casca}: define a assinatura da decisão, não sua implementação.
    O que muda entre o aspirador e o GPT é apenas a função \texttt{decidir()} — e essa
    função pode ser desde 4 ifs até uma rede de 175 bilhões de parâmetros. Isso revela
    que o paradigma de agentes é \textbf{universal}: ele unifica toda a IA sob uma
    única abstração.

    (b) Sim, prompt engineering é conhecimento artesanal — mas de um tipo novo. É
    artesanal porque um humano está manualmente ajustando a entrada para melhorar
    a saída, sem que o modelo aprenda com esse ajuste. Mas difere do artesanal
    clássico porque (1) o modelo subjacente foi aprendido de dados, não programado,
    e (2) o custo de ajustar um prompt é ordens de grandeza menor que o de escrever
    regras. Dito isso, a Bitter Lesson sugere que, no longo prazo, métodos que
    \textit{aprendem} a fazer prompting (auto-prompting, RLHF) substituirão o
    prompting manual — e já há evidências disso.

    \errocomum{Em (b), responder simplesmente ``não'' porque o modelo foi treinado com ML.
    O prompt em si é conhecimento artesanal injetado \textit{na inferência}, não no treino.}

    \competencia{Reconhecer o loop percepção→decisão→ação como a interface universal
    da IA e aplicar a Bitter Lesson a desenvolvimentos recentes.}
  }
\end{questao}

\vspace{1cm}
\hrule
\vspace{0.3cm}
\begin{center}
  {\footnotesize Lista S01 · GBC063 · Prof. Marcelo Keese Albertini · FACOM/UFU}
  \ifdefined\GABARITO
    \ifnum\GABARITO=1
      \\[0.2cm] {\footnotesize\color{red}\textbf{GABARITO — USO EXCLUSIVO DO PROFESSOR}}
    \fi
  \fi
\end{center}

\end{document}
